\documentclass[10pt, twoside]{scrartcl}
\usepackage[ansi]{umlaute}
\usepackage[german]{babel}
\usepackage{fancyhdr}
\usepackage[paperwidth=170mm, paperheight=220mm, top=25mm, bottom=25mm, inner=25mm, outer=20mm]{geometry}
\usepackage[T1]{fontenc}
%
\usepackage{lmodern}
\usepackage{microtype}
\usepackage{graphicx}
\usepackage{makeidx}
\usepackage{amsmath}
\usepackage{amssymb}
\usepackage{theorem}
\title{Softwareentwicklung mit FlowProtocol 2}
\author{Wolfgang Maier}
%%
%% Layout
%%
\setlength{\parindent}{0mm}
\setlength{\parskip}{0.7ex plus0.2ex minus0.1ex}
\addtolength{\headheight}{2.5pt}
% Fancyhdr konfigurieren
\pagestyle{fancy}
\fancyhf{} % L�scht alle Standard-Kopf- und Fu�zeilen
% Kopfzeilen konfigurieren
\fancyhead[LE]{\nouppercase{\rightmark}} % Section auf geraden Seiten rechts (L=left, E=even)
\fancyhead[LO]{\thepage}
\fancyhead[RO]{\nouppercase{\leftmark}} % Subsection auf ungeraden Seiten links (R=right, O=odd)
\fancyhead[RE]{\thepage}
%%
%% Commands
%%
\newcommand{\FPZ}{\textit{Flow\-Proto\-col 2}}
\newcommand{\FPZb}{\textit{Flow\-Proto\-col-2}}
\newcommand{\fcmd}[1]{\texttt{\textasciitilde #1}\index{#1@\texttt{\textasciitilde #1}}}
\newcommand{\cidxvar}[1]{\texttt{\$#1}\index{#1@\texttt{\$#1}}}
\newcommand{\cidxfrag}[1]{\texttt{#1}\index{#1@\texttt{#1}}}
\newcounter{bspcount}
\newcommand{\Beispiel}[1]{\refstepcounter{bspcount}\subsubsection*{Beispiel~\arabic{bspcount} -- #1}}
\newcommand{\synvar}[1]{\ensuremath{\langle\text{#1}\rangle}}
\theorembodyfont{\rmfamily}
\newtheorem{grundelement}{Kernelement}[section]
\makeindex
\begin{document}
\thispagestyle{empty}
\maketitle
\begin{center}
    \includegraphics{FlowLogo}
\end{center}
\newpage
\tableofcontents
\newpage
\thispagestyle{empty}
\section*{Vorwort}
Die vorliegende Anleitung beschreibt M�glichkeiten, die Anwendung \FPZ\ in der professionellen Softwareentwicklung einzusetzen und damit die besonderen Herausforderungen zu meistern, die sich in diesem Bereich stellen. An vielen kleinen und auch auch gr��eren Beispielen wird gezeigt, wie man Skripte in \FPZ\ erstellt und welche Bandbreite an Hilfsmitteln sich damit bereitstellen lassen. Man lernt, Skripte als Ablageform f�r Wissen und Konventionen innerhalb eines Entwicklungsteams zu sehen und die eigenen F�higkeiten in dieser Form weiterzugeben.

\FPZ\ ist die komplett �berarbeitete Nachfolgeversion von FlowProtocol, das im Winter 2021/22 entwickelt wurde. Es handelt sich um eine kleine Anwendung. die �ber einen Browser bedient wird, und auf der Skripte ausgef�hrt werden k�nnen. Die Skripte bestehen aus einfachen Textdateien, die in einem beliebigen Editor erstellt werden k�nnen. Die Verwaltung der Skripte erfolgt in einer lokalen Verzeichnisstruktur. mit der eine organisatorische oder aufgabenbezogene Gliederung abgebildet werden kann.

Bei der Ausf�hrung eines Skriptes werden Informationen �ber Eingabefelder abgefragt und �ber die Anweisungen im Skript verarbeitet. Die Ausgabe erfolgt ausschlie�lich als Ergebnisdokument im Browser, aus dem dann z.B. Textpassagen �ber die Zwischenablage weiterverwendet werden k�nnen. Alle eingegebenen Daten werden ausschlie�lich als Parameter in der URL verwaltet, es gibt keine angebundene Datenbank, keine Benutzerverwaltung und es werden durch die Anwendung keine Dateien erstellt und ge�ndert.

Der Anwendungsfall, f�r den FlowProtocol urspr�nglich entwickelt wurde, ist die Erstellung von Checklisten, die durch Interaktion mit dem Benutzer auf einen individuellen Fall zugeschnitten werden, und so beliebig ins Detail gehen k�nnen, ohne unn�tige Eintr�ge aufzulisten. Daran werden sich auch die ersten Beispiele in dieser Anleitung orientieren. Auf dieser Basis entstanden zahlreiche Skripte f�r den Product Owner, die f�r die  verschiedenen Standardentwicklungen alle ben�tigten Einstellungen und Informationen abfragten, die durch das Designsystem und die damit assoziierten Framework-Komponenten verf�gbar waren. Der Vorteil f�r den Product Owner bestand darin, dass ihm aufgrund der in den Skripten hinterlegten Abh�ngigkeiten immer nur die Optionen angeboten wurden, die f�r den jeweiligen Fall sinnvoll waren, und er so auch an alle Entscheidungen herangef�hrt wurde, die an der jeweiligen Stelle getroffen werden mussten. Das Ergebnisdokument bestand in diesem Fall aus einem sehr umfangreichen Userstory-Entwurf, in dem schon alle abgefragten Informationen und Entscheidungen eingearbeitet waren, und der mehr oder weniger nur noch um einige Benennungen und Aufz�hlungen angereichert werden musste. Schon bald wurde FlowProtocol um die M�glichkeit erweitert, auch Texteingaben abzufragen und diese in das Ergebnisdokument einzuarbeiten.

Die Systematik, die auf Seiten der Story-Formulierungen m�glich ist, ist haupts�chlich begr�ndet durch die Komponenten des �ber viele Jahre aufgebauten eigenen Frameworks und deren M�glichkeiten und die Erfahrung des Teams, also das �ber zahlreiche Wiederholungen aufgebaute Wissen, wie man wiederkehrende Muster umsetzt, und was dabei zu beachten ist. Entsprechend lag es nahe, �hnliche Unterst�tzungswerkzeuge auch auf Entwicklungsseite zu schaffen, die den Programmierer dazu anleiten, die richtigen Klassen zu verwenden und am Ende idealerweise sogar fertigen Programmcode erstellen. Zu diesem Zweck wurde die Formatierungsform als Code geschaffen, bei der man einen Block per Schaltfl�che in die Zwischenablage kopieren kann, und der \fcmd{CamelCase}-Befehl, mit dem Namen f�r Felder, Funktionen oder Variablen erzeugen k�nnen. Die Verwendung interaktiver Anleitungen mit Code-Generierung steigert nicht nur die Effizienz bei der Umsetzung von Standardaufgaben, sie hilft auch bei der Sicherstellung von Einheitlichkeit und ist damit eine wichtige S�ule in der Qualit�tssicherung.

Auch wenn \FPZ\ selbst nicht direkt mit anderen Anwendungen interagiert, so kann in einem Skript praktisch jede URL mit beliebigen Parametern aufgebaut werden, so dass sowohl bei der Ausf�hrung, als auch aus dem Ergebnisdokument der parametrisierte Aufruf anderer Web-, oder Intranet-Anwendungen m�glich ist. Allen voran k�nnen Skripte auf diese Weise andere Skripte aufrufen, aber auch viele der im Entwicklungsumfeld eingesetzten Anwendungen f�r die Verwaltung von Vorg�ngen oder als Wiki bieten gute Steuerungsm�glichkeiten. Schon allein mit der �bergabe eines Suchbegriffs kommt man schon recht weit, und �ber einen MailTo-Link lassen sich sogar vollst�ndig ausformulierte E-Mails vorbereiten. Die Codierung der Parameter f�r die URL l�sst sich dabei sehr einfach mit dem UrlEncode-Befehl umsetzen.

Inzwischen umfasst der Befehlssatz von \FPZ\ alle notwendigen Befehle f�r den Aufbau von Programmen, wie If-Abfragen, For-Schleifen, sowie die Definition von Funktionen, die auch rekursiv aufgerufen werden k�nnen. Zusammen mit den verschiedenen Befehlen f�r das Rechnen mit Zahlen und Datum-Uhrzeit-Werten und zur Manipulation von Zeichenketten lassen sich mit sehr geringem Aufwand kleinere und gr��ere n�tzliche Hilfsanwendungen schreiben, die unmittelbar auf jedem Arbeitsplatz verf�gbar sind.

Der Nutzen von \FPZ\ innerhalb eines Teams oder einer Einrichtung erh�ht sich in besonderem Ma�e dadurch, dass die Skripte von einer gr��eren Zahl an Kollegen gepflegt und erweitert werden, und dass der bestehenden Erfahrungsschatz auf diese Weise permanent durch neues Wissen erweitert wird. In einfacher Form l�sst sich dies sogar hinbekommen, ohne dass daf�r Programmierkenntnisse vorausgesetzt werden, indem ein Skript die f�r seine eigene Erweiterung notwendigen Informationen abfragt, und daraus den resultierenden Programmcode samt Einbauanleitung selbst erzeugt. Auch diese Methode wird an einem Beispiel beschrieben.

Bis dahin werden aber erst einmal die grundlegenden Befehle und Funktionsweisen beschrieben, um einen einfachen Einstieg in \FPZ\ zu erm�glichen. Ich freue mich �ber jeden, der mit Hilfe dieser kleinen Anleitung auf Entdeckungstour geht, und w�nsche viel Spa� und einen hoffentlich nutzbringenden Einsatz.

Zuletzt m�chte ich noch Danke sagen, an alle Kollegen bei easySoft, die \FPZ\ bei der t�glichen Arbeit eingesetzt, und durch ihr Feedback zur permanenten Verbesserung beigetragen haben, insbesondere an \'Eric Louvard, der dort f�r eine permanent robuste und leicht aktualisierbare Installation der Anwendung gesorgt hat. Vielen Dank auch an alle Entwickler und Mitentwickler der zahlreichen Open-Source-Produkte, die ich bei der Entwicklung und der Erstellung dieser Dokumentation mit viel Freude genutzt habe.

\hspace*{\fill}Wolfgang Maier

\newpage

\section{Bezug und Konfiguration}\label{BezugKonfig}
\FPZ\ steht unter der MIT-Lizenz und ist unter folgender Adresse auf GitHub verf�gbar :
\begin{center}
    \texttt{https://github.com/maier-san/FlowProtocol2}
\end{center}

Der Programmcode kann direkt mit Git oder der Entwicklungsumgebung Visual Studio Code in ein lokales Verzeichnis, z.B. \verb|D:\Apps\FlowProtocol2|, �bertragen, und dort  mit dem dotnet-Befehl kompiliert werden, wobei als Zielframework\index{Zielframework} .NET 8 vorausgesetzt wird:
\begin{verbatim}
dotnet.exe build
    D:\Apps\FlowProtocol2/FlowProtocol2/FlowProtocol2.csproj
\end{verbatim}


Die Konfiguration erfolgt �ber die Datei \emph{appsettings.json}\index{appsettings.json}, in der haupts�chlich der Pfad auf das Skripte-Verzeichnis\index{Skripte-Verzeichnis} mit dem Parameter \cidxfrag{ScriptPath} eingestellt werden muss. Dieses Verzeichnis enth�lt die Skripte und Unterverzeichnisse, die von \FPZ\ auf der Startseite angezeigt werden, und ist damit der Dreh- und Angelpunkt der Skriptverwaltung. F�r die ersten Versuche kann man den Parameter auf den Scripts-Ordner innerhalb des Projektes setzen, etwa so:
\begin{verbatim}
"ScriptPath": "D:\\Anwendungen\\FlowProtocol2\\Scripts", 
\end{verbatim}
In einem Unternehmen wird man das Verzeichnis so w�hlen, dass die Mitarbeiter, die aktiv an den Skripten arbeiten, dort direkt Dateien editieren und erstellen k�nnen, z.B. indem man dieses auf einem Netzlaufwerk verf�gbar macht. Zus�tzlich wird man die Skripte regelm��ig sichern, idealerweise mit Hilfe einer Versionsverwaltung. Wenn man Manipulation bef�rchtet, kann man das Editieren auch vollst�ndig auf den Weg �ber die Versionsverwaltung beschr�nken, verbaut sich damit aber die M�glichkeit, die Wirkung von �nderungen an einem Skript unmittelbar nach dem Speichern durch die Aktualisierung des Browser-Tabs zu �berpr�fen. In diesem Fall w�re dann eine getrennte Skript-Entwicklungsumgebung sinnvoll, analog zu den sonstigen Entwicklungsumgebungen im Umfeld der Softwareherstellung. \FPZ\  selbst ben�tigt auf dieses Verzeichnis nur Leserechte.

In den Programm- und Konfigurationsdateien \emph{launchSettings.json}\index{launchSettings.json}, \emph{appsettings.json} und \emph{Program.cs}\index{Program.cs} sind die Ports\index{Ports} 5000, bzw.\ 5001 eingetragen, die man je nach Installation auch nochmal anpassen m�chte.

Wie schon im Vorwort beschrieben ist die M�glichkeit, Links zu generieren, ein m�chtiges Mittel um die Interaktion mit anderen Anwendungen zu erm�glichen. Wie man jedoch aus jeder IT-Sicherheitsbelehrung wei�, kann das Anklicken von Links auch Gefahren mit sich bringen, und speziell in Phishing-E-Mails nutzen Angreifer vertrauensw�rdig aussehende Links, um den Empf�nger auf eine nachgebaute oder mit Schadcode gespickte Seite zu leiten. Auch mit den Links in einem Skript sind solche Angriffe m�glich, wenn auch aufgrund der Beschr�nkung auf die eigenen Mitarbeiter eher unwahrscheinlich. Aus diesem Grund werden alle durch ein Skript ausgegebenen Links, deren Dom�ne nicht in der Auf\-lis\-tung des \cidxfrag{LinkWhitelist}-Parameters steht, zus�tzlich zu dem Anzeigetext mit der vollst�ndigen URL ausgegeben. Damit kann bei der Konfiguration entschieden werden, welche Seitenaufrufe so vertrauensw�rdig sind, dass sie auch nur mit dem Anzeigetext dargestellt werden k�nnen, was zum einen die Lesbarkeit erh�ht und zum anderen die Aufmerksamkeit des Anwenders auf die Links konzentriert, die nicht diese Einstufung haben.

F�r die regelm��ige lokale Bereitstellung von \FPZ\ ist es am einfachsten, vom Basisrepository auf GitHub mittels Fork zu verzweigen und die lokalen Anpassungen im eigen Zweig zu verwalten.

\section{Grundlagen}
\subsection{Das erste Beispiel}
In diesem Abschnitt bleibt der Bezug zur Softwareentwicklung erst einmal darauf beschr�nkt, dass wohl jeder, der in dieser Branche t�tig ist, schon an dem einen oder anderen Hallo-Welt-Beispiel vorbeigekommen ist. In dieser Tradition starten auch wir und beginnen mit dem einfachsten aller Beispiele:

\Beispiel{Hallo Welt}
\begin{verbatim}
>> Hallo Welt!
\end{verbatim}

Das bedeutet, wir erstellen eine neue Textdatei mit dem Dateinamen \emph{Hallo Welt.fp2} und speichern diese direkt, oder in einem Unterordner im Skripte-Verzeichnis der \FPZb-Installation (vgl.\ Abschnitt \ref{BezugKonfig}).  Und ja, sp�testens jetzt ist es Zeit, \FPZ\ in einer eigenen Umgebungen zum Laufen zu bringen, denn nur durch Ausf�hren der gezeigten Beispiele und viel eigenem Rumprobieren ist es m�glich, den bestm�glichen Nutzen aus dieser Anleitung zu ziehen. Nat�rlich sind alle Beispiele aus dieser Anleitung auch als Textdatei im Projekt enthalten, so dass man diese nur auszuw�hlen braucht.

Nach erfolgreicher Konfiguration wird man von \FPZ\ nach dem Start mit dem Text \emph{Willkommen bei FlowProtocol 2} begr��t und man sieht das Logo, das auch das Titelblatt dieser Anleitung ziert. �ber die Schaltfl�che \emph{Zur Skriptauswahl} oder dem Men�punkt \emph{Start} gelangt man von dort aus zur Auf\-lis\-tung der Skript-Hauptgruppen, also den Ordnern im Skripte-Verzeichnis. F�r das Durcharbeiten dieser Anleitung ist es am besten, man kopiert den Ordner \emph{Doku-Beispiele} in das Skripte-Verzeichnis, so dass man die einzelnen Beispiele direkt aufrufen kann.

Die Ausf�hrung des oben genannten Beispiels f�hrt zur Ausgabe des Textes \emph{1.\ Hallo Welt!} unter der �berschrift \emph{Hallo Welt}.

Als �berschrift wird standardm��ig der Dateinahme ohne Endung ausgegeben, und da wir keinen Befehl angegeben haben, der explizit eine �berschrift setzt, ist das auch hier der Fall. Der Befehl \cidxfrag{>\mbox{}>} gibt den dahinter stehenden Text aus, standardm��ig als nummerierte Aufz�hlung, was die \emph{1.} erkl�rt.

Die Schaltfl�chen \emph{Neue Ausf�hrung} und \emph{Zur�ck zur Startseite} werden immer angezeigt, wenn die Ausf�hrung abgeschlossen ist. Mit \emph{Neue Ausf�hrung} kann das aktuell ausgew�hlte Skript neu gestartet werden und mit \emph{Zur�ck zur Startseite} kommt man zur�ck zur Auswahl der Skript-Hauptgruppen, was auch mit dem Men�punkt \emph{Start} m�glich ist. Wir werden noch einige Anwendungsf�lle sehen, die sich am besten mit einem Skript umsetzen lassen, das mehrmals hintereinander ausgef�hrt wird.

\subsection{Eine erste Auswahlabfrage}
Bis jetzt haben wir nur Text in Form gebracht und in einem Browser-Fenster ausgegeben, was in Anbetracht der recht einfachen Syntax und der sch�nen Formatierungsm�glichkeiten auch schon Mehrwerte bieten kann, aber der Hauptnutzen von \FPZ\ liegt ganz klar in der Interaktivit�t. Diese l�sst sich umsetzen mit Hilfe von Auswahlabfragen und Texteingaben. Das nachfolgende Beispiel zeigt eine einfache Auswahlabfrage.

\Beispiel{Hallo Welt mit Auswahl}
\begin{verbatim}
?Q: Wie soll die Welt begr��t werden?
    #h: Mit "Hallo Welt!"
        >> Hallo Welt!
    #a: Mit "Aloah Welt!"
        >> Aloah Welt!
\end{verbatim}

Beim Ausf�hren dieses Skriptes wird man als Anwender zuerst mit der Frage \emph{Wie soll die Welt begr��t werden?} konfrontiert, mit den beiden Antwortm�glichkeiten \emph{Mit ''Hallo Welt!''} und \emph{Mit ''Aloah Welt!''}. Es muss einer der beiden vorgegebenen Auswahlm�glichkeiten gew�hlt werden, denn wenn man die Skriptausf�hrung mit \emph{Weiter} fortsetzt, werden die nicht beantworten Fragen einfach erneut gestellt und man gelangt nicht zur Ausgabeseite.  Nach getroffener Wahl wird auf der Ausgabeseite dann entweder \emph{1.\ Hallo Welt!} oder \emph{1.\ Aloah Welt!} ausgegeben. Dort hat man dann die M�glichkeit, das eben ausgef�hrte Skript mittels \emph{Neue Ausf�hrung} erneut auszuf�hren oder zur Startseite der Skriptauswahl zu wechseln.

Das Grundprinzip, dass auf einer oder mehreren Eingabeseiten\index{Eingabeseiten} die eventuell auch voneinander abh�ngigen Eingaben abgefragt werden, um dann am Ende die daraus resultierenden Ausgaben auf einer Ausgabeseite anzuzeigen, gilt f�r alle Skripte.

Ein Blick in der Adresszeile des Browsers zeigt, dass die Antwort auf die Frage bei Wahl von \emph{Hallo Welt} als Parametersequenz \verb|Q=h| in der URL gespeichert wurde. Das ist auch der einzige Ort, wo sich diese Information wiederfindet, denn \FPZ\ selbst speichert keine Eingabedaten auf Serverseite. Die beiden Buchstaben \verb|Q| und \verb|h|, bzw.\ \verb|a| sind hierbei durch den Skriptcode festgelegt und hei�en Schl�ssel\index{Schl�ssel}. 
Die Auswahlabfrage wird durch das \cidxfrag{?}-Zeichen eingeleitet, unmittelbar gefolgt vom Schl�ssel der Frage, gefolgt von einem Doppelpunkt an den sich die als Text ausformulierte Fragestellung oder Eingabeaufforderung anschlie�t. Einger�ckt auf erster Ebene stehen unter der Frage die verschiedenen Antwortm�glichkeiten, die jeweils mit dem \cidxfrag{\#}-Zeichen beginnen, analog gefolgt vom Schl�ssel der Antwort, einem Doppelpunkt und der Ausformulierung der jeweiligen Antwortm�glichkeit oder Option. 

Der wiederum unterhalb einer Antwortm�glichkeit einger�ckte Skriptcode kann wieder aus Befehlen bestehen, die auch f�r sich alleine stehen k�nnen, Dieser Code wird wird nur dann ausgef�hrt, wenn die dazugeh�rende Antwortm�glichkeit ausgew�hlt wurde.

Die Schl�ssel f�r Auswahlabfragen und die sp�ter noch kommenden Texteingaben k�nnen prim�r aus einer Folge aus Buchstaben und Ziffern, sowie runden Klammern bestehen, die Schl�ssel f�r Auswahlwerte nur aus Buchstaben und Ziffern. Sie dienen wie schon beschreiben dazu, die Eingaben des Anwenders in der URL zu speichern und auch innerhalb des Skriptes abrufbar zu machen. Der Schl�ssel f�r eine Auswahlabfrage muss eindeutig innerhalb des gesamten Skriptes sein, der einer Antwortm�glichkeit muss eindeutig innerhalb einer Auswahlabfrage sein. 

Da die Gesamtl�nge der URL f�r Seitenaufrufe auf knapp �ber 2000 Zeichen begrenzt ist, muss man sich �ber die L�nge von Schl�sseln erst Gedanken machen, wenn man wirklich viele Auswahlabfragen in ein Skript einbaut.

\subsection{Allgemeine Erg�nzungen zur Syntax}
In \FPZ\ wird jede nicht leere Zeile einem Befehl zugeordnet. Die Einr�ckung\index{Einr�ckung} am Anfang wird getrennt erfasst und verarbeitet, Leerzeilen\index{Leerzeilen} und Leerraum am Ende wird ignoriert. Die Einr�ckung kann sowohl mit dem Tabulatorzeichen\index{Tabulatorzeichen}, also auch mit Leerzeichen erfolgen, wobei das Tabulatorzeichen intern in vier Leerzeichen umgerechnet wird. Es sollte also vermieden werden, innerhalb einer Skriptdatei sowohl mit Tabulatorzeichen, als auch mit Lerzeichen einzur�cken.

Die Zeichenfolge \cidxfrag{//} leitet einen Kommentar\index{Kommentar} ein, jedoch nur zu Beginn einer Zeile.

Lange Zeilen k�nnen auf mehrere Zeilen umgebrochen werden, wobei der Umbruch\index{Umbruch} durch die Zeichenfolge \cidxfrag{\_\_} zu Beginn jeder Folgezeile kenntlich gemacht werden muss.

\Beispiel{Syntaxerg�nzungen}
\begin{verbatim}
// Das ist ein Kommentar

>> Das ist eine Ausgabe,
    __ die im Scriptcode
    __ auf drei Zeilen verteilt wird.
\end{verbatim}

\subsection{Die FlowProtocol-2-Befehlsreferenz}
Um einen �berblick �ber alle in \FPZ\ vorhandenen Befehle zu bekommen und f�r jeden die genaue Syntax, die dabei m�glichen Fehler und die Verwendung anhand eines kleinen Beispiels nachschlagen zu k�nnen, ist die \FPZb-Befehlsreferenz\index{Befehlsreferenz} der richtige Ort. Gemeint ist der Ordner \emph{FP2-Tutorial}, der im \emph{Scipts}-Ordner zusammen mit der Anwendung bereitgestellt wird. Dort ist f�r jeden Befehl ein Skript vorhanden, das diesen soweit es m�glich ist, von anderen Befehlen unabh�ngig beschreibt. In vielen F�llen erzeugt das Skript direkt eine Ausgabe, die auch das Beispiel mit einschlie�t. Bei Befehlen die sich auf den Eingabeseiten auswirken, muss zuerst eine Eingabe durchlaufen werden, sodass man die Wirkung des dokumentierten Befehls direkt in der Programmoberfl�che sehen kann.

Erg�nzend dazu gibt es in Abschnitt~\ref{Befehlsreferenz} ebenfalls eine Befehlsreferenz, die alle Befehle auf\-listet, und auch auf deren Verwendung in dieser Dokumentation hier verweist.

\subsection{Fehlermeldungen und Debug-Techniken}
Das Entwickeln von Skripten wird nicht ohne Fehler\index{Fehler} ablaufen, wie auch? Wichtig ist, dass man diese schnell erkennt und korrigieren kann, und zielsicher an sein Ergebnis kommt.

\FPZ\ kann zwar nicht mit einem syntaxunterst�tzenden Editor aufwarten, liefert daf�r aber klare Hinweise auf Syntax- und Laufzeitfehler. Diese werden mit Dateiname und Zeilennummer benannt. Der am h�ufigsten auftretende Fehler d�rfte die \emph{Parsing Exception}\index{Parsing Exception} sein, also das Problem, dass eine Zeile nicht interpretiert werden kann. Da in \FPZ\ jeder Befehl in einer eigenen Zeile angegeben wird, muss man dementsprechend nur noch schauen, welchen Befehl man in der jeweiligen Zeile verwenden wollte, und wo die Zeile syntaktisch abweicht.

Hier ein Beispiel f�r so einen Fehler:

\Beispiel{Fehler im Skript}
\begin{verbatim}
?W: Wo ist der Fehler?
    #A: In Antwort A
        >> Antwort A
    #B: In Antwort B
        Antwort B
\end{verbatim}

Laufzeitfehler\index{Laufzeitfehler} treten dagegen meist im Zusammenhang mit Variablen auf, z.B.\ wenn mit diesen Rechenoperationen durchgef�hrt werden sollen, diese aber keinen Zahlenwert enthalten. Auch hier kann man anhand der Angaben in der Fehlermeldung die Stelle lokalisieren und sich an die Kl�rung der Ursache machen.

In den meisten restlichen F�llen erscheint keine Fehlermeldung, aber das Skript produziert nicht die gew�nschte Ausgabe. Und dann geht es an das Debuggen\index{Debuggen}, also das Entfernen der Fehler. Hier gibt es bei \FPZ\ den gro�en Vorteil, dass das Feedback nach einer �nderung sehr schnell zu bekommen ist. Man �ndert einfach eine Stelle im Skriptcode, speichert und dr�ckt die Aktualisieren-Schaltfl�che im Browser und erh�lt umgehend das Resultat der aktualisierten Version. Man muss weder neu kompilieren, noch irgendwelche Eingaben wiederholen.

Der beste Ansatz, um inhaltlichen Fehlern auf die Spur zu kommen, besteht darin, an vielen Stellen Hilfsausgaben zu erzeugen, z.B.\ mit den Werten von Variablen, der Information, ob eine bestimmten Stelle durchlaufen wurde, und so weiter. Diese k�nnen ganz normal als Ausgaben in das Skript eingebaut werden, idealerweise gesammelt in einem eigenen Debug-Abschnitt (siehe Abschnitt~\ref{Abschnitte}). Sobald die Entwicklung des Skriptes abgeschlossen ist, werden diese Hilfsausgaben auskommentiert oder entfernt.

\section{Entscheidungsb�ume}
Aus mehreren ineinander verschachtelten Auswahlabfragen lassen sich gro�e Entscheidungsb�ume und Flussdiagramme aufbauen, deren St�rke darin liegt, an jeder Stelle im Skriptcode die vollst�ndige bis dahin getroffene Auswahl zu kennen, und so die n�chste Fragestellung ganz exakt auf diese Situation abzustimmen. Die Anwendungsf�lle daf�r sind enorm vielf�ltig, aber um im Bereich der Softwareentwicklung zu bleiben, sind hier einige typische Beispiele.

\subsection{Strukturabbildung einer Software}\label{Strukturabbildung}
Durch den Nachbau der Struktur, bzw.\ des Men�baums\index{Men�baum} einer Software von der Programmoberfl�che aus betrachtet, kann man es dem Anwender des Skriptes enorm erleichtern, eine Stelle innerhalb dieses Softwareproduktes eindeutig zu benennen. Der daraus resultierende Pfad oder eine diesem zugeordnete Kennung kann f�r die Suche nach Vorg�ngen und Informationen verwendet werden. Da diese Art von Information oft Teil abteilungs�bergreifender Kommunikation ist, k�nnen mit einer solchen Auswahl sowohl Missverst�ndnissen vermieden und auch l�ngerer Tastatureingaben eingespart werden.

Beispiel~\ref{BspProgrammauswahl} zeigt, wie man einen solchen Auswahlbaum als wiederverwendbare Komponente anlegt.

\subsection{Kategorisierungshilfen}\label{Kategorisierungshilfen}
Eine wohl in allen Softwareh�usern stattfindende Kategorisierung\index{Kategorisierung} ist die des Schweregrades von neu erfassten Fehlern. Diese variieren in der Regel von \emph{hoch kritisch} bis hin zu \emph{niedrig} oder einer numerischen Abstufung mit dieser Entsprechung. Die subjektive und oft situationsbedingte Einsch�tzung der beteiligten Personen sollte hierbei nicht der Gradmesser sein, sondern vielmehr die objektiven Kriterien und der Vergleich mit Referenzbeispielen. Auf diese Weise k�nnen die verf�gbaren Ressourcen zielgerichtet eingesetzt werden.

Die nachfolgende Schweregradeinstufung\index{Schweregradeinstufung} beginnt damit, explizit nach Kriterien f�r die h�chste Einstufung zu fragen und geht weiter zu den Kriterien der zweiten Stufe, wenn diese nicht zutreffen. Nach der zweiten Stufe wird direkt nach Kriterien f�r die unterste Stufe gefragt, so dass automatisch f�r alles die Zuordnung zu Stufe 3 erfolgt, was auch dort nicht zugeordnet werden kann. Das macht die Aufstellung der Kriterien leichter, da sich einfacher Kriterien f�r die ganz niedrige Priorit�t finden lassen, als f�r die vorgelagerte Stufe.

\Beispiel{Schweregradeinstufung}
\begin{verbatim}
?S1: Treffen die Kriterien f�r einen Stufe-1-Fehler zu?    
    #k1: Ausgegebene Werte entsprechen nicht der Spezifikation.
        >> Stufe-1-Fehler: Spezifikationsverletzung
    #k2: Der Datenschutz ist an einer wesentlichen Stelle verletzt.
        >> Stufe-1-Fehler: Datenschutzverletzung
    #k3: Die Datensicherheit ist verletzt oder gef�hrdet.
        >> Stufe-1-Fehler: Datensicherheitssverletzung
    #no: Nein, die o.g. Kriterien treffen nicht zu.
        ?S2: Treffen die Kriterien f�r einen Stufe-2-Fehler zu?
            #k1: Der Fehler st�rt zentrale Programmfunktionen.
                >> Stufe-2-Fehler: St�rung zentraler Funktion.
            #k2: Es gibt zwei oder mehr Supportanfragen zu 
                __ diesem Fehler.
                >> Stufe-2-Fehler: Mindestens zwei Supportanfragen
            #no: Nein, die o.g. Kriterien treffen nicht zu.
                ?S4: Treffen die Kriterien f�r einen 
                    __ Stufe-4-Fehler zu?
                    #k1: Es handelt sich um einen erkennbaren
                        __ Tippfehler.
                        >> Stufe-4-Fehler: Tippfehler
                    #k2: Der Fehler tritt nur bei unsinnigen
                        __ Eingaben auf.
                        >> Stufe-4-Fehler: unsinnige Eingaben
                    #no: Nein, die o.g. Kriterien treffen nicht zu.
                        >> Stufe-3-Fehler
\end{verbatim}

F�r die Abarbeitung der Fehler-Aufgaben der verschiedenen Schweregradstufen kann man dann klare Regeln definieren, wie z.B.\ dass Stufe-1-Fehler umgehend behoben werden und dass die L�sung aller Stufe-2-Fehler zumindest in das n�chste Service Release einflie�en muss.

Insgesamt ist die Sortierung nach dem Schweregrad keine gute Idee, wenn sich viele Vorg�nge ansammeln, da es in diesem Fall erst dann zur Bearbeitung eines Stufe-4-Fehlers kommt, wenn alle Stufe-3-Fehler abgearbeitet sind, was unter Umst�nden nie der Fall ist. Ein System, das die Einstufung ber�cksichtigt, aber gleichzeitig die Abarbeitung jeder Aufgabe sicherstellt, ist angelehnt an die Spuren der Autobahn. Die Aufgaben jeder Schweregradstufe bilden jeweils eine Queue, in der neu erstellte Aufgaben unten angef�gt, und die �ltesten oben abgearbeitet werden. Um den Schweregrad zu ber�cksichtigen, l�sst man die linke Spur schneller laufen, indem man die Mengen f�r die Abarbeitung in ein vorgegebenes Verh�ltnis setzt, z.B.
\[M(2):M(3):M(4) = 6 : 3 : 1\]
d.h.\ f�r eine abgearbeitete Stufe-4-Aufgabe werden drei Stufe-3-Aufgaben und sechs Stufe-2-Aufgaben abgearbeitet.

In Beispiel~\ref{BspBugPlaner} wird ein Skript gegeben, dass die Menge der Vorg�nge in jedem Schweregrad f�r eine vorgegebene Anzahl so berechnet, dass sie so gut wie m�glich in diesem Verh�ltnis steht.

\subsection{Expertensysteme}\label{Expertensystem}
Mit Expertensystem\index{Expertensystem} oder auch Troubleshooting-System\index{Troubleshooting-System} oder Chatbot\index{Chatbot} ist hier eine Anwendung gemeint, die wie ein lebendiger Experte durch immer weiter in die Tiefe gehende Fragestellungen eine in der Regel problematische Situation m�glichst genau erfasst, um dann dem Anwender dazu passende Handlungsanweisungen und L�sungsans�tze aufzuzeigen. Dort wo diese nicht gelingt, l�sst sich anhand der beantworteten Fragen zumindest eine gute und aufs Wesentliche beschr�nkte Beschreibung der Problemsituation erstellen.
Probleme gibt es �berall und die zur L�sung bef�higten Experten erkennt man zuallererst daran, dass sie die richtigen Fragen stellen. Die M�glichkeit, einen gro�en Teil der Probleme auch in Abwesenheit menschlicher Experten l�sen zu k�nnen, ist unheimlich wertvoll, insbesondere dann, wenn eine gut gemachte Anleitung Zeit spart und die Sicherheit gibt, nicht irgendetwas riskantes auszuprobieren.

Wie schon bei dem zuletzt genannten Beispiel, das auch mit dem Expertensystem kombiniert werden kann, zeigt sich der enorme Nutzen, der sich schon aus diesen wenigen Skriptbausteinen ergibt. Der Grund daf�r ist nicht, dass diese Befehle irgendetwas raffiniert technisches machen, sondern dass sie es erm�glichen, strukturelles oder fachliches Wissen in eine Form zu bringen, so dass dieses Wissen vom Anwender unmittelbar genutzt werden kann, ohne dass er sich dieses Wissen aneignen muss. \FPZ\ �bernimmt die Anwendung des Wissens und der Anwender selbst muss sich immer nur mit einer einzelnen und idealerweise f�r ihn verst�ndlichen Fragestellung auseinandersetzen, bekommt F�hrung und Anleitung und muss nicht selbst m�hsam zum Experten werden.

Der Umfang solcher Entscheidungsb�ume\index{Entscheidungsb�ume} kann riesig werden und trotzdem muss der Anwender am Ende nur eine Handvoll Fragen beantworten, um zum Ziel zu kommen.

Hier ist ein sehr vereinfachtes Beispiel f�r ein Expertensystem, das eine Problemsituation analysiert und eine Handlungsempfehlung gibt:
\Beispiel{Mini-Expertensystem}\label{Bsp04}
\begin{verbatim}
?: Startet die Anwendung?
    #: Ja
        >> Anwendung startet
        ?: Ist eine Benutzeranmeldung m�glich?
            #: Ja
                >> Benutzeranmeldung m�glich.
                >> Wo liegt das Problem?
            #: Nein
                >> Benutzeranmeldung nicht m�glich
                >> Empfehlung: Passwort zur�cksetzen
    #: Nein
        >> Anwendung startet nicht
        ?: Wurde das System schon neu gestartet?
            #: Ja
                >> System wurde schon neu gestartet.
                >> Empfehlung: 2nd-Level-Support kontaktieren
            #: Nein
                >> System wurde noch nicht neu gestartet.
                >> Empfehlung: System neu starten
\end{verbatim}

Zuerst f�llt auf, dass hier keine Schl�ssel\index{Schl�ssel} angegeben sind, weder f�r die Auswahlabfragen, noch f�r die Antwortm�glichkeiten. Tats�chlich sind diese nicht zwingend erforderlich und werden von \FPZ\ selbst�ndig anhand der Reihenfolge im Code vergeben, wenn sie im Skript weggelassen werden. Dies ist dann m�glich, wenn man sich nicht an einer anderen Stelle im Skript auf eine Auswahl beziehen oder dieses wiederholen m�chte, und es ist dann besonders vorteilhaft, wenn der Baum sehr gro� werden kann und permanent erweitert wird und man den �berblick �ber die schon vergebenen Schl�ssel zu verlieren droht. Der Nachteil der impliziten Schl�sselvergabe liegt darin, dass sich durch das Einf�gen einer weiteren Auswahlabfrage am Anfang alle darunter liegenden Schl�ssel verschieben, und so eine URL, die die Antworten in Bezug auf die Vorversion enth�lt, in der neuen Version falsch interpretiert wird. 

Die automatische Nummerierung\index{automatische Nummerierung} von Schl�sseln f�r Auswahlabfragen und Texteingabe ist auch explizit m�glich, in dem man einem selbstgew�hlten Basisschl�ssel das \verb|'|-Zeichen anschlie�t. Mehrere nacheinander im Scriptcode vorkommenden Schl�sselangaben \verb|Q'| werden dann intern zu \verb|Q_1|, \verb|Q_2|, \verb|Q_3| usw.\ expandiert. Durch den Unterstrich werden �bereinstimmungen mit explizit vergebenen Schl�sseln vermieden.

Die Ausf�hrung dieses Beispiels zeigt dar�ber hinaus, wie das System mit ineinander verschachtelten Fragestellungen umgeht. In jedem Schritt werden immer genau die Fragen gestellt, die noch offen sind, und die aufgrund der vorangegangenen Schritte durchlaufen werden m�ssen. Sobald alle Fragen beantwortet sind, werden die gesammelten Ausgaben ausgegeben und die Ausf�hrung des Skriptes ist abgeschlossen.

\section{Formatierung der Ausgabe}\label{FormatierungAusgabe}
Dieser Abschnitt widmet sich der Formatierung der Ausgabe und beschreibt M�glichkeiten, Abschnitte zu bilden, Aufz�hlungen zu verschachteln, sowie Links und Code in die Ausgabe zu integrieren. 

\subsection{Ausgabe in Abschnitten}\label{Abschnitte}
Bis jetzt habe wir nur den Befehl \verb|>>| verwendet, um Aufz�hlungspunkte auszugeben, und dabei merkt man recht schnell, dass das Ergebnis zumeist nicht nur aus einer Liste besteht, oder bestehen k�nnte. Nehmen wir das Expertensystem aus Beispiel~\ref{Bsp04}. Die Beschreibung der Situation und der schon durchgef�hrten Ma�nahmen ist inhaltlich betrachtet eine Liste f�r sich, die Handlungsempfehlung eine andere und l�uft alles am Ende in eine Kontaktierung des 2nd-Level-Supports hinaus, k�nnte die Liste der ben�tigten Informationen und Materialien als dritte Liste ausgegeben werden.

Dieser Anwendungsfall zeigt auch schon, dass es durchaus �blich ist, dass verschiedenen Listen in einem Skript nicht nacheinander, sondern parallel zusammengestellt werden. Dies ist auch im folgenden Beispiel der Fall, das eine sehr einfache, nicht interaktive Implementierungsanleitung darstellt, die einen Entwickler bei der Implementierung eines neuen Moduls anleitet. Da jeder Entwicklungsschritt auch einen Integrationstest nach sich ziehen sollte, und der Entwickler in diesem Fall auch f�r die Erstellung der Testpunkte verantwortlich ist, wird zusammen mit der Anleitung auch gleich noch eine Liste von Testpunkten ausgegeben.

\Beispiel{Ausgabe in Abschnitten}
\begin{verbatim}
@Anleitung >> Erstelle eine neue Modul-Klasse
@Testpunkte >> Das Modul wird korrekt gestartet
@Anleitung >> Implementiere die Start-Methode
@Testpunkte >> Das Modul wird korrekt beendet
@Anleitung >> Implementiere die Beenden-Methode
@Anleitung >> Nehme das Modul in den Modulkatalog auf
@Anleitung >> �bernehme die Testpunkte von unten
\end{verbatim}

Die Ausf�hrung erzeugt einen Abschnitt\index{Abschnitt} \emph{Anleitung} mit entsprechender �berschrift und den f�nf Punkten der Anleitung, gefolgt von einem Abschnitt \emph{Testpunkte} mit den beiden Testpunkten. Im Gegensatz zur Zusammenstellung erfolgt die Ausgabe Abschnittsweise, wobei die Reihenfolge durch die jeweils erste Ausgabe des Abschnitts festgelegt wird.

Die Erzeugung einer Ausgabe in einem Abschnitt setzt gleichzeitig diesen Abschnitt f�r alle darauffolgenden Ausgaben, sofern diese nicht explizit einem Abschnitt zugeordnet sind. Um also nur die Anleitung auszugeben, reicht eine einzige Festlegung des Abschnitts am Anfang:

\Beispiel{Nur ein Abschnitt}
\begin{verbatim}
@Anleitung >> Erstelle eine neue Modul-Klasse
>> Implementiere die Start-Methode
>> Implementiere die Beenden-Methode
>> Nehme das Modul in den Modulkatalog auf
\end{verbatim}

\subsection{Unterpunkte und Absatzformate}
Eine gute Anleitung ist gleicherma�en geeignet f�r erfahrene und weniger erfahrene Benutzer und erg�nzt die prim�r durchzuf�hrenden Schritte mit einer detaillierten Beschreibung der dazugeh�renden Teilschritte. \FPZ\ unterst�tzt generell drei Aufz�hlungsebenen, die mit den Befehlen \cidxfrag{>\mbox{}>}, \cidxfrag{>} und \cidxfrag{>.} angesteuert werden, jeweils mit der optional m�glichen Angabe eines Abschnitts mittels \verb|@|. Ausgaben der zweiten Ebene gliedern sich der letzten Ausgabe der ersten Ebene des jeweiligen Abschnitts unter und Ausgaben der dritten Ebene dementsprechend der letzten Ausgabe der zweiten Ebene.

Zus�tzlich kann man f�r beide Ebenen angeben, ob eine Ausgabe als nummerierter Aufz�hlungspunkt (\verb|>>#|), nicht nummerierter Aufz�hlungspunkt (\verb|>>*|), einfache Textzeile (\verb|>>_|) oder Codezeile  (\verb=>>|=) ausgegeben wird.

\Beispiel{Unterpunkte}
\begin{verbatim}
@Anleitung >>_ Gesch�tzter Zeitaufwand ca. 15 min
>>* Erstelle eine neue Modul-Klasse.
    >* Erstelle eine C#-Datei mit dem Namen des Moduls.
>>* F�ge folgenden Code ein:
    >| public class NeuesModul : BaseModul
    >| {
    >| }
>>* Implementiere die Start-Methode.
    ># Gibt die Zeichenfolge "override" innerhalb der Klasse ein
        >.* Der Editor bietet Methoden zur Auswahl an.
    ># und w�hle in der Auswahl des Editors die Methode Start.
    ># Erg�nze den Start-Code des Moduls wie im Wiki beschrieben.
\end{verbatim}

Die komplette Ausgabe erfolgt im Abschnitt \emph{Anleitung}, die Angabe des gesch�tzten Zeitaufwands erfolgt als normale Textzeile, die Aufz�hlungspunkte\index{Aufz�hlungspunkte} der ersten Ebene sind entgegen dem Standard nicht nummeriert. Unter dem ersten Aufz�hlungspunkt wird ein ebenfalls nicht nummerierter Unterpunkt ausgegeben, unter dem zweiten ein Codeblock\index{Codeblock} mit drei Zeilen. Unter dem dritten Aufz�hlungspunkt sieht man drei Unterpunkte\index{Unterpunkte}, die mit Kleinbuchstaben durchnummeriert sind, der erste davon hat wiederum einen nicht nummerierten Unterpunkt.

Die drei aufeinanderfolgenden Codezeilen werden zu einem Codeblock zusammengefasst, der mit der automatisch darunter angeordneten Schaltfl�che \emph{In Zwischenablage kopieren} in die Zwischenablage\index{Zwischenablage} genommen werden kann.

Auf die vielf�ltigen M�glichkeiten der Codegenerierung kommen wir sp�ter noch zur�ck. Im n�chsten Abschnitt bleiben wir nochmal bei der Formatierung und zeigen, wie sich Text innerhalb einer Zeile formatieren l�sst.

\subsection{Links und Inline-Code}\label{LinksInline}
Gerade in Anleitungen wird auf Stellen im Code in Form von Klassen-, Methoden- oder Variablennamen verwiesen, und da erleichtert es den Lesefluss deutlich, wenn man diese Textelemente in der Ausgabe entsprechend hervorhebt. Dies ist mit dem Befehl \fcmd{AddCode} m�glich, der der Ausgabe der letzten Zeile Text hinzuf�gt, der als Code\index{Code} formatiert ist. Mit dem Befehl \fcmd{AddText} kann danach wieder weiterer Text hinzugef�gt werden. Man beachte das zus�tzliche Leerzeichen, das immer am Anfang des angeh�ngten Textes angef�gt werden muss, wenn zwischen diesem und dem vorangegangen Text ein Leerzeichen stehen soll, da Leerraum am Zeilenende generell ignoriert wird.

Als Anwendung im Browser kann \FPZ\ nat�rlich auch Links erzeugen, und damit den Anwender sowohl zu statischen Seiten, also auch zu anderen Anwendungen weiterleiten. Auch dieser M�glichkeit widmen wir sp�ter noch einen eigenen Abschnitt. Ein Link\index{Links} besteht dabei aus einer URL und dem anzuzeigenden Text, die in dieser Reihenfolge, getrennt von einem vertikalen Strich angegeben werden. Auch hier kann mit \fcmd{AddText} wieder weiterer Text im Anschluss angef�gt werden.

\Beispiel{Links und Inline-Code}\label{BspLinksUndInlineCode}
\begin{verbatim}
@Anleitung >> Implementiere die Start-Methode
    > Gibt die Zeichenfolge 
    ~AddCode  override 
    ~AddText  innerhalb der Klasse ein
    > und w�hle in der Auswahl des Editors die Methode Start.
    > �berschreibe die Start-Methode wie im
    ~AddLink  https://learn.microsoft.com/de-de/dotnet/csharp
        __/language-reference/keywords/override | Internet
    ~AddText  beschrieben.
\end{verbatim}

Das Beispiel formatiert den Text \emph{override} als Code und den Text \emph{Internet} als Link, also blau, unterstrichen und klickbar. Sofern die Dom�ne der angegebenen URL in der Konfiguration nicht als sicher angegeben ist (siehe Kernelement~\ref{defLinks}), wird die URL aus Sicherheitsgr�nden hinter dem Anzeigen-als-Text erg�nzt. Der Link wird standardm��ig immer in einem neuen Tab ge�ffnet.

\section{Variablen und Texteingaben}
Die erste Zielsetzung der Vorversion von \FPZ\ bestand tats�chlich nur darin, eine gro�e Menge an m�glichen Ausgaben mittels iterierten Fragestellungen auf eine spezifische Situation einzugrenzen und die damit verbundenen Anwendungsf�lle abzudecken. Sehr schnell wurde jedoch klar, dass mit ein paar grundlegenden Erweiterungen noch sehr viel mehr m�glich ist.

\subsection{Variablen setzen und verwenden}\label{VariablenSetzenUndVerwenden}
Unter einer Variablen\index{Variablen} versteht man kurz gesagt einen Platzhalter, der verschiedene Werte annehmen kann. Innerhalb von Programmier- oder Skriptsprachen lassen sich die Werte von Variablen setzen und auch abrufen. In \FPZ\ wird einer Variablen mit dem \fcmd{Set}-Befehl ein Wert zugewiesen. Die Variable kann dann an nahezu allen Stellen im Skript verwendet werden, indem man Sie mit vorangestelltem \verb|$|-Zeichen kennzeichnet.

\Beispiel{Hallo Welt mit Variable}\label{BspHalloWeltMitVariable}
\begin{verbatim}
?Q: Wie soll die Welt begr��t werden?
    #h: Mit "Hallo Welt!"
        ~Set Gruss=Hallo
    #a: Mit "Aloah Welt!"
        ~Set Gruss=Aloah
>> $Gruss Welt!
\end{verbatim}

Das Beispiel setzt  die Variable \verb|Gruss| je nach Auswahl auf den Wert \emph{Hallo} oder \emph{Aloah} und ruft den Wert in der Ausgabe am Ende auf, so dass entweder \emph{Hallo Welt!} oder \emph{Aloah Welt!} ausgegeben wird.

Eine Variable kann aus Buchstaben (ohne Umlaute), Zahlen, runden Klammern und anderen Variablen zusammengesetzt werden. Der Wert einer Variable kann jede beliebige Zeichenkette sein. 

Die Ersetzung von Variablen durch die jeweils zugeordneten Werte erfolgt durch einfache Ersetzung, d.h.\ das Ende der Variablen muss beim Aufruf nicht gekennzeichnet werden. Variablen, die nicht gesetzt wurden, werden auch nicht ersetzt.

\Beispiel{Variablenersetzung}\label{BspVariablenersetzung}
\begin{verbatim}
~Set X=abc
>> X = $X
>> Z = $Z (Z wurde nie zugeweisen)
>> XY = $XY (vor der Zuweisung von XY)
~Set XY=def
>> XY = $XY (nach der Zuweisung von XY)
~Set i=5
~Set F($i)=ghi
>> F(i) = F($i) = $F($i)
\end{verbatim}

Dieses Beispiel zeigt den Umgang mit Variablen, deren Anfang identisch mit einer anderen Variable ist, in diesem Fall \verb|XY|. Die Ersetzung erfolgt in diesem Fall absteigend sortiert, sodass \verb|$XY| vor \verb|$X| ersetzt wird. Die Verwendung der Variablen \verb|$i| im Ausdruck \verb|$F($i)| gibt schon einen Ausblick darauf, dass man auch ganze Felder von Variablen anlegen und durchlaufen kann. Bei der Auswertung wird hierbei zun�chst \verb|$i| durch \verb|5| und dann \verb|$F(5)| durch \verb|ghi| ersetzt. Die runden Klammern sind hier �brigens nur semantischer Zucker, um den Index-Teil der Variablen hervorzuheben und vom Feldnamen abzugrenzen.

Wir werden sp�ter noch viele andere Befehle kennenlernen, die Zeichenketten verarbeiten und die das Ergebnis in einer Variablen zur�ckgeben und lernen auch selbst Funktionen zu schreiben, die ihre Eingabewerte �ber Variablen �bergeben bekommen. 

\subsection{Bewertungssysteme}\label{BewSysteme}
Das Beispiel in diesem Abschnitt verwendet eine Variable, um einen Zahlenwert zu verwalten und in Abh�ngigkeit der Auswahlwerte Werte zu addieren. Dies ist z.B.\ dann erforderlich, wenn man mehrere gleichartige Dinge nach einem formalen Bewertungssystem\index{Bewertungssystem} bewerten m�chte, um sie anschlie�end anhand der daran angeschlossenen Metrik in eine Rangfolge zu bringen. 

Nehmen wir konkret die Bewertung von Aufgaben, die zur Stablilisierung bestimmter Programmfunktionen in einem Softwareprodukt erstellt wurden. Die Erstellung und Ausformulierung solcher Aufgaben werden meist den Entwicklern selbst �berlassen, da nur diese �ber das Wissen und den technischen Einblick verf�gen, um dort Verbesserungspotential auszumachen. Am Ende sollte es aber auch hier der Product Owner sein, der �ber Umfang und Priorit�t der Ma�nahmen entscheidet, und um die daf�r relevanten Aspekte der Entwicklung herauszuarbeiten, k�nnte er jeder dieser Aufgaben mit dem folgendem Bewertungssystem bewerten lassen:

\Beispiel{Bewertungssystem}\label{BspBewertungssystem}
\begin{verbatim}
~Set punkte=0
@Bewertung >>* Bewertungkrterien: 
?A1: Wie gro� ist der Verbesserungsbedarf an dieser Stelle?
    #w1: Gering. Es gibt wenig Meldungen zu Einschr�nkungen an
        __ dieser Stelle.
        >* Bedarf: gering (-)
        ~AddTo punkte+=10
    #w2: Mittel. Es gibt immer wieder Meldungen zu sp�rbaren 
        __  Einschr�nkungen an dieser Stelle.
        >* Bedarf: mittel
        ~AddTo punkte+=25
    #w3: Gro�. Es gibt permanent Meldungen zu st�renden
        __ Einschr�nkungen an dieser Stelle.
        >* Bedarf: gro� (+)
        ~AddTo punkte+=50
?A2: Wie klar ist die durchzuf�hrende Ma�nahme beschrieben?
    #w1: Unkrokret. Es sind noch Analysen und Vorarbeiten notwendig.
        >* Klarheit: unkonkret (-)
        ~AddTo punkte+=10
    #w2: Hinreichend konkret. Der Ansatz ist klar beschrieben,
        __ aber noch nicht erprobt.
        >* Klarheit: konkret
        ~AddTo punkte+=25
    #w3: �bertragbar. Der Ansatz kann von einer anderen Stelle
        __ hierher �bertragen werden.
        >* Klarheit: �bertragbar (+)
        ~AddTo punkte+=50
?A3: Wie effektiv wird die Ma�nahmen voraussichtlich sein?
    #w1: Unklar. Der tats�chliche Effekt ist erst nach der
        __ Umsetzung erkennbar.
        >* Effektivit�t: unklar (-)
        ~AddTo punkte+=10
    #w2: Gering bis mittel. Es bleiben Einschr�nkungen, aber
        __ weniger oft und weniger gro�.
        >* Effektivit�t: gering bis mittel
        ~AddTo punkte+=25
    #w3: Mittel bis gut. Die vorhandenen Einschr�nkungen werden
        __ weitestgehend behoben.
        >* Effektivit�t: mittel bis gut (+)
        ~AddTo punkte+=50
?A4: Wie aufw�ndig ist die Umsetzung der Ma�nahme?
    #w1: Sehr aufw�ndig. 20 Storypunkte oder mehr.
        >* Aufwand: hoch (-)
        ~AddTo punkte+=10
    #w2: Aufw�ndig. 13-20 Storypunkte.
        >* Aufwand: m��ig hoch
        ~AddTo punkte+=25
    #w3: Im Rahmen. Maximal 8 Storypunkte.
        >* Aufwand: im Rahmen (+)
        ~AddTo punkte+=50
?A5: Wie gut ist die Wiederverwendbarkeit?
    #w1: Gering. Die L�sung ist speziell auf diese eine Stelle
        __ zugeschnitten.
        >* Wiederverwendbarkeit: gering (-)
        ~AddTo punkte+=10
    #w2: �bertragbar. Die L�sung kann auf andere Stellen
        __ �bertragen werden.
        >* Wiederverwendbarkeit: �bertragbar
        ~AddTo punkte+=25
    #w3: Umfassend. Die L�sung wird zentral eingebaut und
        __ wirkt sich an mehren Stellen aus.
        >* Wiederverwendbarkeit: umfassend (+)
        ~AddTo punkte+=50
@Bewertung >>* Gesamtbewertung: $punkte Punkte.
\end{verbatim}

Am Anfang wird die Variable \verb|punkte| auf den Wert 0 gesetzt, was man auch weglassen k�nnte, und danach folgen f�nf Fragen, die jeweils ein Kriterium der in der Aufgabe beschriebenen Entwicklungsma�nahme abfragen. F�r jede Frage werden drei Auswahlm�glichkeiten angeboten, die jeweils aufsteigend eine geringe bis gute Erf�llung des genannten Kriteriums beschreiben. Der ausgew�hlte Wert wird dann sowohl in einer Zusammenfassung in Textform ausgegeben, als auch in der Gesamtpunktezahl ber�cksichtigt, indem zu dem in der Punkte-Variablen vorhandenen Wert mehr oder weniger hinzugez�hlt wird. Die Erh�hung einer Variablen um einen Wert erfolgt dabei mit dem Befehl \fcmd{AddTo}. Die Gesamtpunktezahl wird dann ebenfalls am Ende ausgegeben.

Die Verwendung eines solchen Bewertungsskriptes hat prim�r den Vorteil, dass die Bewertung sehr objektiv ausf�llt, da die Einsch�tzung in Bezug auf einzelnen Kriterien schon von sich aus reflektiert passieren muss und nicht mittels Bauchgef�hl, was zus�tzlich dadurch verst�rkt werden kann, dass die Formulierung der Auswahlwerte eine gewisse Belegbarkeit assoziiert, die insbesondere bei einer gemeinsamen Bewertung im Team auch ausdiskutiert werden kann. Durch das Festhalten der ausgew�hlten Kriterien zusammen mit der Bewertung ist deren Begr�ndung zudem gut dokumentiert. Die wiederholte Auseinandersetzung mit den zum gr��ten Teil wirtschaftlichen Kriterien sch�rft in diesem Beispiel auch gleich die Sichtweise der Entwickler auf diese Aspekte, so dass diese bei der Suche nach Verbesserungen immer mehr in den Fokus gelangen.

Ob man nun wie im Beispiel oben eine einfache additive Bewertung abbildet, oder komplexere Formeln verwendet, die auch ber�cksichtigen, dass sich Faktoren gegenseitig verst�rken k�nnen, bleibt der eigenen Kreativit�t �berlassen. Mit \FPZ\ lassen sich alle Formeln realisieren.

Viele zumeist abstrakte Fragestellungen lassen sich auf derartige Bewertungssysteme herunterbrechen: Wie stark wird durch eine Entwicklungsma�nahme die Produktstrategie verfolgt? Welcher �ber die Jahre gerechnete Mehraufwand wird durch Support und Aufrechterhaltung einer Funktion notwendig sein?
In der Regel sind es viele kleine Faktoren, wie die Abh�ngigkeit von Drittanbieter-Komponenten oder externen Diensten die L�nge der Kommunikationswege, der Grad der Konfigurierbarkeit, das Hinzukommen neuer Strukturen mit Semantik und Bedienkonzepten und die allgemeine Komplexit�t und Fehlertoleranz, die dauerhaft Aufwand verursachen. Die Identifikation dieser Faktoren kann zum einen schon bei der Planung der Entwicklung helfen, die Kosten-Nutzen-Verteilung besser einzusch�tzen, zum anderen k�nnen durch ein geeignetes Skript aber auch direkt geeignete Ma�nahmen abgeleitet werden, um diese Dinge bestm�glich und systematisch in den Griff zu bekommen.

\subsection{Texteingaben}\label{Texteingaben}
Schon bei der Erstellung der ersten Skripte gab es die Assoziation zu den zahlreichen Konfigurationsportalen, mit denen man im Internet sein Fahrrad oder Auto nach Wunsch zusammenstellen konnte, und daran angelehnt auch die ersten Versuche, die Konfiguration rund um die vielfach wiederkehrenden Standardentwicklungen in gleicher Weise abzufragen, um es so dem Product Owner zu erleichtern, die passende Konfiguration zu w�hlen. Da die Auswahl schon technisch auf auf die Menge der m�glichen Konfigurationen beschr�nkt blieb, gleichzeitig aber alle notwendigen Entscheidungen abgefragt wurden, konnte man so Aufgabenbeschreibungen generieren, die fast alle wesentlichen Konfigurationsaspekte auf\-listeten. Manuell nachgearbeitet werden mussten individuelle Beschreibungstexte und eben Begriffe und Benennungen, die daf�r meist mit einem gut erkennbaren Platzhalter \emph{xxx} in der Ausgabe ersetzt wurden, was insbesondere bei den wiederholt vorkommenden Begriffen den Komfort schm�lerte.

Inzwischen geh�ren Texteingaben\index{Texteingaben} zum Grundumfang von \FPZ\ und zusammen damit wurden auch zahlreiche Befehle hinzugef�gt, um Texteingaben zu verarbeiten.

\Beispiel{Textersetzung}\label{BspTextersetzung}
\begin{verbatim}
~Input si: Suche in
~Input sn: Suche nach
~Input ed: Ersetze durch
~Execute
~Replace erg=$si|$sn->$ed
@Aufgabe >> Ersetze "$sn" in "$si" durch "$ed".
@Ergebnis >> $erg
\end{verbatim}

In diesem Beispiel werden die drei Texte f�r eine einfache Ersetzungsaufgabe mit dem \fcmd{Input}-Befehl abgefragt und �ber die zu den Schl�sseln geh�renden Variablen \verb|si|, \verb|sn| und \verb|ed| an den \fcmd{Replace}-Befehl �bergeben, der das Ergebnis �ber die Variablen \verb|erg| bereitstellt.

Der Befehl \fcmd{Execute} sorgt daf�r, dass das Skript bis zu dieser Stelle ausgef�hrt wird, solange, bis alle dar�ber stehenden Eingabebefehle abgearbeitet wurden, und die Eingabewerte in den Variablen zur Verf�gung stehen. Er ist an dieser Stelle notwendig, da Variablen ohne Wert beim \fcmd{Replace}-Befehl zu Fehlern f�hren w�rden. Der Befehl kann auch bewusst zu Formatierung der Eingabe genutzt werden wie in Abschnitt~\ref{FormatierungDerEingabe} beschrieben wird.

Die M�glichkeit, Texte und Zahlen �ber eine einfach per Link adressierbare Webanwendung abzufragen und zur Durchf�hrung von Berechnungen zu verwenden, ist eine gute Alternative f�r so manche Kalkulationshilfe, die momentan noch in Form von Tabellenkalkulationsdateien verteilt werden.

\section{Weitere Formatierungsm�glichkeiten}
In Abschnitt~\ref{FormatierungAusgabe} haben wir ja schon gesehen, wie man das Format der Aufz�hlung variieren, und die Ausgabe in Abschnitte unterteilen kann. Hier gibt es nur noch wenig zu erg�nzen. Mindestens in gleicher Weise wichtig sind jedoch die Formatierungsm�glichkeiten, bei den Seiten, in denen die Eingabe abgefragt wird. Sie m�ssen den Bogen zum Anwendungskontext schlie�en und genau erkl�ren und dem Anwender dabei helfen, die richtigen Daten einzugeben und die korrekte Auswahl zu treffen.

\subsection{Ausgabetitel und Abschnittsverschiebungen}\label{AusgabetitelUndAbschnittsverschiebungen}
Der Titel\index{Titel}, der auf der Ausgabeseite angezeigt wird, kann mit dem Befehl \fcmd{SetTitle} unabh�ngig vom Dateinamen gesetzt werden, so dass man als Dateinamen und damit auch f�r die Auf\-lis\-tung im Men� k�rzere Titel w�hlen kann. Beim setzen des Titels im Scriptcode lassen sich auch Variablen verwenden, die die gemachten Eingaben ber�cksichtigen. Gerade wenn man die Ergebnisseite eines Skriptes an jemand anderen weitergeben m�chte, z.B. �ber den Link mit den Parametern, kann ein m�glichst passender Titel helfen, um das Dokument einzuordnen.

Beim Verschieben eines Abschnitts\index{Abschnitt!verschieben} wird der Inhalt eines Abschnitts an einen anderen Abschnitt angeh�ngt, wobei dieser ggf.\ zuvor erstellt wird. Dies erm�glicht es, mehrere Informationen zun�chst parallel in verschiedenen Abschnitten zu sammeln und diese dann in einem Abschnitt aneinanderzuf�gen. Das ist besonders n�tzlich, wenn die Reihenfolge der Informationen in der Ausgabe eine andere ist als in der Abfrage.

\Beispiel{Abschnittsverschiebung}\label{BspAbschnittsverschiebung}
\begin{verbatim}
~Input Bez: Wie soll das Feld hei�en?
@AnlEig >> Setze die Bezeichnung "$Bez"
?: Welche Art von Werten soll eingegeben werden?
    #: Texte
        ~Set Klasse=TextFeld        
    #: Zahlen
        ~Set Klasse=Zahlenfeld
    #: Datumsangaben
        ~Set Klasse=Datumsfeld
    #: Werte aus einer vorgegebenen Menge
        ~Set Klasse=Auswahlfeld
~SetTitle Implementerung der $Klasse-Instanz $Bez
@AnlInst >> Erstelle eine $Klasse-Instanz
?: Kann das Feld leer gelassen werden?
    #j: Ja
        ~Set KannLeerSein=true
    #n: Nein
        ~Set KannLeerSein=false
@AnlEig >> Setze KannLeerSein = $KannLeerSein
~MoveSection AnlInst -> Anleitung
~MoveSection AnlEig -> Anleitung
\end{verbatim}

In diesem Beispiel wird eine kleine Anleitung f�r den Einbau eines Feldes in ein Programm erstellt. Hierbei werden Bezeichnung, Typ und die Eigenschaft, ob das Feld leer gelassen werden kann, abgefragt. Sobald Bezeichnung und Typ festgelegt sind, wird der Titel unter Verwendung dieser Informationen neu gesetzt. Die Anleitung wird zun�chst in zwei Abschnitten getrennt zusammengestellt, \verb|@AnlEig| f�r die beiden Eigenschaften und \verb|@AnlInst| f�r die Instanziierung. Am Ende werden die beiden Abschnitte  mit dem \fcmd{MoveSection}-Befehl nacheinander in den neuen Abschnitt \emph{Anleitung} �berf�hrt, so dass die Punkte dort in einer logischen Reihenfolge stehen.

F�r dieses Minimalbeispiel h�tte man die gew�nschte Reihenfolge auch ohne die Verschiebung von Abschnitten hinbekommen, aber gerade bei gr��eren Anleitungsskripten kann eine Aufteilung helfen, die jeweils zu einem Anleitungsschritt dazugeh�renden Aspekte wie Programmcode und Testpunkte im Skript eng beieinanderzuhalten.

\subsection{Formatierung der Eingabe}\label{FormatierungDerEingabe}
So wie man den Titel\index{Titel} f�r die Ausgabeseite setzen kann, l�sst sich der Titel auch f�r die Eingabeseiten setzen und im Laufe der Skriptausf�hrung �ndern, ebenso lassen sich auch bei der Eingabe Abschnitte einf�gen, um die Eingabeseiten weiter zu gliedern. Zuletzt kann auch noch unterhalb der �berschrift eine Beschreibung des Skriptes ausgegeben werden, um auch den Anwender abzuholen, der nur durch die Skriptauswahl st�bert, und der bislang nur den Dateinamen kennt.

\Beispiel{Bug-Planer}\label{BspBugPlaner}
\begin{verbatim}
~SetTitle Bug-Planer
~SetInputTitle Bug-Planer 6-3-1
~SetInputDescription Der Bug-Planer berechnet
    __ f�r eine Anzahl von Bugs, wie viele Bugs von
    __ jedem Schweregrad eingeplant werden m�ssen,
    __ damit diese im Verh�ltnis 6:3:1 liegen.
~SetInputSection Vorhandene Bug-Vorg�nge
~Input A2: Anzahl Bugs mit Schweregrad 2
~Input A3: Anzahl Bugs mit Schweregrad 3
~Input A4: Anzahl Bugs mit Schweregrad 4
~SetInputSection Planung
~Input AP: Wie viele Bugs m�chtest du einplanen?
~Execute
~Set S2=0
~Set S3=0
~Set S4=0
~Set Sum=0
~Split Indexlist=2,3,2,4,2,3,2,2,3,2|,
~Set anybug = true
~DoWhile $Sum<$AP && $anybug
    ~ForEach idx in Indexlist
        ~If $Sum<$AP && $S$idx<$A$idx
            ~AddTo S$idx+=1        
            ~AddTo Sum+=1
    ~Loop
    ~EvalExpression anybug = $S4<$A4 || $S3<$A3 || $S2<$A2
~Loop
@Ergebnis >>* Eingeplant werden sollen $AP Vorg�nge.
>>* Folgende Aufteilung wird vorgeschlagen:
    > $S2 von $A2 Bugs mit Schweregrad 2
    > $S3 von $A3 Bugs mit Schweregrad 3
    > $S4 von $A4 Bugs mit Schweregrad 4
\end{verbatim}

Der Bug-Planer\index{Bug-Planer} aus dem Beispiel oben ist eine sch�nes Beispiel f�r ein Berechnungsskript, bei dem Werte abgefragt, und daraus andere Werte berechnet werden. In diesem Fall wird die Anzahl der Bug-Vorg�nge f�r jeden Schweregrad von 2 bis 4 abgefragt, der z.B. wie in Abschnitt~\ref{Kategorisierungshilfen} beschrieben festgelegt werden kann. Zus�tzlich wird noch die Anzahl der einzuplanenden Bug-Vorg�nge abgefragt und damit dann ein Vorschlag f�r eine Aufteilung berechnet, die m�glichst im Verh�ltnis $6:3:1$ steht. Eine solche Angleichung geht nat�rlich nur, wenn von allen Schweregraden ausreichend viele Vorg�nge vorhanden sind. Wenn nicht, wird das verplant, was da ist und das Sollverh�ltnis auf die anderen Schweregrade angewendet.

F�r die Eingabeseite wurde hier mit \fcmd{SetInputTitle} eine eigene �berschrift gew�hlt. Zus�tzlich wurde mit \fcmd{SetInputDescription} eine Skriptbeschreibung\index{Skriptbeschreibung} formuliert, die unterhalb der �berschrift ausgegeben wird. Die vier Eingabefelder wurden mit dem Befehl \fcmd{SetInputSection} in die beiden Abschnitte\index{Abschnitt!bei Eingabe} \emph{Vorhandene Bug-Vorg�nge} und \emph{Planung} unterteilt. Man beachte, dass gleich benannte Eingabeabschnitte an verschiedenen Stellen nicht zusammengef�hrt werden.

Nach der vierten Eingabe steht der Befehl \fcmd{Execute}, der ebenfalls zur Formatierung verwendet werden kann. Der Befehl bewirkt, dass der Aufbau der Eingabeseite an diesem Punkt abgebrochen, und die Eingabeseite im aktuellen Zustand angezeigt wird. Erst wenn alle Abfragen und Eingaben bis zu diesem Punkt get�tigt wurden, l�uft das Skript weiter bis zum Ende oder bis zum n�chsten \fcmd{Execute}-Befehl. Bei Skripten mit sehr vielen Abfragen kann man diesen Befehl nutzen, um die Fragen so aufzuteilen, dass sie jeweils auf einen Bildschirm passen oder man kann erzwingen, dass die von einer Auswahl abh�ngigen Eingaben zuerst vollst�ndig abgefragt werden, bevor der nachfolgenden Block ausgef�hrt wird. Im Beispiel oben sorgt der Befehl daf�r, dass die Berechnungen in der zweiten H�lfte des Skriptes erst ausgef�hrt werden, wenn die an die Eingaben gebundenen Variablen mit Werten belegt sind. Ansonsten w�rden die ersten mit diesen Variablen durchgef�hrten Vergleiche zu einem Fehler f�hren.

Die Berechnung selbst funktioniert wie folgt: Zun�chst werden Variablen \verb|S2|, \verb|S3| und \verb|S4| f�r die Zahl der zu verteilenden Aufgaben des jeweiligen Schweregrades auf $0$ gesetzt, ebenso eine Variable \verb|Sum| f�r die Zwischensumme der insgesamt schon verteilten Aufgaben. 

Die �u�ere \fcmd{DoWhile}-Schleife wiederholt so lange, wie diese Zwischensumme noch kleiner als die Anzahl der zu verteilenden Aufgaben ist, und noch ein Bug-Vorgang zum verteilen �brig ist. Letzteres wird mit dem Befehl \fcmd{EvalExpression} getrennt ausgewertet und in der Variablen \verb|anybug| abgespeichert.

In der Inneren Schleife wird die Indexsequenz $2,3,2,4,2,3,2,2,3,2$ durchlaufen, was die Index-Variable \verb|idx| nacheinander auf diese Werte setzt. Dies bewirkt, dass die Indices $2$, $3$ und $4$ gleichm��ig im Verh�ltnis $6:3:1$ durchlaufen werden. F�r den jeweils ausgew�hlten Index wird dann die Anzahl der zu verteilenden Aufgaben $S_i$ um $1$ erh�ht, sofern noch Aufgaben des jeweiligen Schweregrades �brig sind ($S_i<A_i$) und die Zwischensumme $S$ der schon verteilten Aufgaben noch kleiner ist, als die Anzahl $A$ der Aufgaben, die verteilt werden soll ($S<A$), also \verb|$S$idx<$A$idx| und \verb|$Sum<$AP|. F�r jede verteilte Aufgabe wird auch die Zwischensumme um $1$ erh�ht.

Eine nette Erweiterung dieses Skriptes k�nnte so aussehen, dass man die Gesamtzahl der zu verplanenden Vorg�nge nicht abfragt, sondern prozentual aus der Zahl der vorhandenen Bugs berechnet und dabei nach oben und unten durch einen absoluten Wert begrenzt. Zus�tzlich k�nnte man das Planungsverh�ltnis abfragen und das Gesamtpaket aus Prozentwert, obere und untere Grenze, sowie Verh�ltnis als fertige Planungsstrategie ausw�hlbar machen. 

F�r ein beliebiges Verh�ltnis $S_0:S_1:\dots :S_{v-1}$ f�r die Einplanung von $v$ verschiedenen Schweregraden und eine Folge $a=(a_i)_{i\in n}$ von Indices $a_i\in v$ kann man berechnen, f�r welchen n�chsten Index $a_n\in v$ die erweiterte Gesamtfolge $a'=(a_i)_{i\in n+1}$ die kleinste Abweichung $d(a_n)$ zum vorgegebenen Verh�ltnis hat, mit
\[d(a_n)=\min_{l\in v}\left(1+ \big|\{i\in n\colon a_i=a_n\}\big|-\frac{(n+1)S_l}{\sum_{j\in v}S_j}\right)\]
Diesen w�hlt man dann, um die Folge zu erweitern. Die Berechnung so oft wiederholt werden, wie man weitere Folgenglieder ben�tigt.

\subsection{Hilfezeilen}\label{Hilfezeilen}
Einer der gro�en Vorteile einer interaktiven Anleitung liegt darin, dass man eben nicht wissen muss, welche Schritte man ausf�hren muss und welche man weglassen kann, weil das die Anleitung ja zusammen mit dem Anwender herausarbeitet. Die Informationen, unter welcher Bedingung welcher Schritt erforderlich ist, und was dabei genau zu tun ist, kann alles vom Ersteller in die Struktur und die Ausgaben des Skriptes gepackt werden, so dass der Anwender nur noch richtig auf die gestellten Fragen antworten, und die ben�tigen Eingaben t�tigen muss.

Diese Fragen und Eingaben so einfach und verst�ndlich zu halten, so dass auch Anwender mit wenig Erfahrung in der Lage sind, das Skript auszuf�hren, ist eine Kunst f�r sich. In manchen F�llen wird das nicht m�glich sein, und man wird einfach ein gewisses Systemwissen\index{Systemwissen} beim Anwender voraussetzen m�ssen. In den meisten F�llen wird man jedoch mit ein oder zwei Zeilen Anleitung auskommen, um zu beschreiben, wo man einen Wert findet oder wie die gestellte Frage oder einzelne Auswahlwerte genau zu interpretieren sind.

\Beispiel{Hilfezeilen}\label{BspHilfezeilen}
\begin{verbatim}
~Input Bez: Trage die Lizenz f�r FlowProtocol 2 ein:
~AddHelpLine Du findest die zugeordnete Lizenz auf der 
~AddHelpLink https://github.com/maier-san/FlowProtocol2
    __| FlowProtocol2-Seite
~AddHelpText  bei GitHub auf der rechten Seite als 
    __ zweiten Punkt unter "About".
?K: Handelt es sich um eine freie Lizenz?
    ~AddHelpLine Suche im Internet nach Details der Lizenz.
    ~AddHelpLine "Freie Lizenz" bedeutet in diesem
        __ Zusammenhang, dass f�r die Nutzung der Software
        __ auch im kommerziellen Umfeld keine Kosten anfallen.
    #: Ja
    #: Nein
\end{verbatim}

Einzelne Hilfezeilen\index{Hilfezeilen} lassen sich f�r Texteingaben und Auswahlabfragen mit dem Befehl \fcmd{AddHelpLine} hinzuf�gen. Mit \fcmd{AddHelpLink} lassen sich, analog wie bei den Ausgaben, auch Links\index{Links} einbauen, so dass man auf Wiki-Seiten oder andere \FPZb-Skripte verweisen kann. Danach kann man mit \fcmd{AddHelpText} wieder normalen Text an die letzte Zeile anh�ngen.


\section{Programmierung}
Auch wenn es urspr�nglich nicht die Absicht war, aus \FPZ\ eine Programmierumgebung zu machen, gab es doch immer wieder Anforderungen, die mit den grundlegenden Befehlen einer �blichen Programmiersprache gut umsetzbar gewesen w�ren. Die Vorteile, den Aufbau von Anleitungen mit Programmlogik zu verkn�pfen, Bedingungen abzupr�fen und Werte zu berechnen, waren insgesamt so gro�, dass schlie�lich alle wichtigen Programmierbefehle in \FPZ\ Einzug fanden.

\subsection{Spr�nge}\label{Spruenge}
Wir starten das Thema Programmierung mit einem Befehl, der in den h�heren Programmiersprachen eher selten zum Einsatz kommt, dem \fcmd{GoTo}-Befehl. Dieser erm�glicht es einen Sprungmarke\index{Sprungmarke} anzuspringen, und die Skriptausf�hrung dort fortzusetzen. Die Sprungmarke wird wiederum mit dem Befehl \fcmd{JumpMark} gesetzt. Eine Sprungmarke kann auch ganz normal als n�chste Zeile in der Ausf�hrung durchlaufen werden. Ihr Durchlaufen selbst hat keinen Effekt.

Der \fcmd{GoTo}-Befehl ist f�r \FPZ\ dahingehend interessant, da er eng mit dem namensgebenden Ursprung der Anwendung zusammenh�ngt und sehr direkt die Abbildung von Flussdiagrammen\index{Flussdiagramm} erm�glicht. Im Gegensatz zur Verschachtelung von Abfragen kann man damit weitere Muster gut umsetzen, ohne dabei Skriptcode wiederholen zu m�ssen, wie etwa im folgenden Beispiel:

\Beispiel{Spr�nge}\label{BspSpruenge}
\begin{verbatim}
?Q1: Um welche Art von Entwicklung handelt es sich?
    #F: Neues Framework
        ~GoTo Framework
    #I: Individuelle Einzelentwicklung
        ~GoTo Einzelentwicklungen
    #V: Standard-Entwicklung auf Framework-Basis
        ~GoTo Alle Entwicklungen
~JumpMark Framework
>> Fragen f�r Framework-Entwicklungen
~JumpMark Einzelentwicklungen
>> Fragen f�r Einzelentwicklungen 
~JumpMark Alle Entwicklungen
>> Fragen f�r alle Entwicklungen 
\end{verbatim}

Zu Beginn wird nach der Kategorisierung einer Entwicklung gefragt, die entweder eine Framework-Entwicklung ist, eine individuelle Einzelentwicklung oder eine Standardentwicklung auf Basis des vorhandenen Frameworks ist. In Abh�ngigkeit der Antwort werden entweder alle drei Bl�cke, nur die letzten beiden Bl�cke oder nur der letzte Block durchlaufen.

Man kann Sprungmarken auch einfach dazu verwenden, um tief verschachtelte Abfragen in �bersichtlicher Art und Weise als Folgen anzuordnen.

\subsection{Berechnungen}\label{Berechnungen}
Aus Abschnitt~\ref{BewSysteme} kennen wir schon den \fcmd{AddTo}-Befehl, mit dem man einen Wert zu einer Variablen dazuz�hlen kann. Sp�ter wurde der \fcmd{Calculate}-Befehl erg�nzt, der eine Berechnung\index{Berechnung} mit zwei Argumenten und einem der Operatoren \verb|+|, \verb|-|, \verb|*|, \verb|/|, und \verb|%| 
(modulo) durchf�hren konnte. Umfangreichere Ausdr�cke mussten dementsprechend in mehreren Zeilen zerlegt und mit entsprechend vielen Hilfsvariablen berechnet werden.

Mit dem Befehl \fcmd{CalculateExpression} lassen sich inzwischen auch l�ngere Ausdr�cke auf Basis der Operatoren \verb|+|, \verb|-|, \verb|*|, \verb|/|, \verb|%| 
und \verb|^| und den Funktionen \cidxfrag{sqrt}, \cidxfrag{sin}, \cidxfrag{cos}, \cidxfrag{tan}, \cidxfrag{exp} und \cidxfrag{ln} und einer beliebigen Klammerung in einer Zeile berechnen. Die Berechnungen erfolgen nummerisch und die Ergebnisse werden in Dezimal- oder ggf.\ auch in Exponentialschreibweise\index{Exponentialschreibweise} dargestellt.

\Beispiel{Berechnungen}\label{BspBerechnungen}
\begin{verbatim}
~CalculateExpression U = 72 % 7
~CalculateExpression V = (-2,5 + 1,57)*(1,3 - 0,4)/18
~CalculateExpression W = 2^3*16^(1/2)+7
~CalculateExpression X = exp(sin(1/4))
~Round Z=$X|3
@Ausgabe >> U = $U
    >> V = $V
    >> W = $W
    >> X = $X
    >> Z = $Z
\end{verbatim}

In die letzten Zeile vor der Ausgabe wird der Wert der Variablen \verb|X| mit dem \fcmd{Round}-Befehl auf drei Stellen hinter dem Komma gerundet und das Ergebnis der Variablen \verb|Z| zugewiesen.

\subsection{Schleifen}\label{Schleifen}
Die gro�e St�rke von Computern liegt darin, die gleichen Dinge in leichter Variation oft zu wiederholen, ohne sich zu langweilen und ohne in der Ausf�hrung nachl�ssig zu werden. Mit Hilfe von Schleifen\index{Schleifen} lassen sich solche Mehrfachausf�hrungen einfach implementieren.

Die bekannten h�heren Programmiersprachen unterscheiden drei prinzipielle Arten von Schleifen: Die \emph{For-Schleife}\index{For-Schleife} wird mit einer Z�hlervariablen eingeleitet, die mit einem Startwert beginnt, diesen bei jedem Durchlauf um eine Schrittl�nge hoch z�hlt, und beendet wird, wenn ein Zielwert erreicht oder eine Abbruchbedingung erf�llt wird. Die \emph{Do-While-Schleife}\index{Do-While-Schleife} l�uft dagegen so lange durch, bis eine Abbruchbedingung erf�llt ist. Schlie�lich gibt es noch die \emph{For-Each-Schleife}\index{For-Each-Schleife}, bei der eine vorgegebene Menge von Elementen durchlaufen wird. Eine Indexvariable nimmt dabei nacheinander den Wert dieser Elemente an.

In \FPZ\ gibt es die Beiden Befehle \fcmd{DoWhile} und \fcmd{ForEach} f�r die Einleitung von Do-While-Schleifen und For-Each-Schleifen. Reine For-Schleifen lassen sich sehr direkt auch als Do-While-Schleife abbilden, deshalb gibt es f�r diese keinen eigenen Befehl. Das Schleifenende, also die Stelle, an der ggf.\ ein weiterer Durchlauf gestartet wird, wird in beiden F�llen mit dem \fcmd{Loop}-Befehl gekennzeichnet, der die gleiche Einr�ckung wie der Beginn der Schleife haben muss. Auf diese Weise lassen sich mehrere Schleifen ineinander verschachteln. Zus�tzlich gibt es noch den \fcmd{ExitLoop}-Befehl, mit dem die aktuell ausgef�hrte Schleife unmittelbar verlassen wird.

In Beispiel~\ref{BspBugPlaner} hatten beide Schleifenvarianten schon einen kleinen Gastauftritt, aber um die Funktionsweise nochmal genauer anzuschauen ist die folgende Primzahlberechnung\index{Primzahlberechnung} gut geeignet:

\Beispiel{Primzahlen}\label{BspPrimzahlen}
\begin{verbatim}
~Set n=2
~Set pidx=0
~DoWhile $n<20
    ~Set prim=ja
    ~ForEach p in PZ
        ~CalculateExpression r = $n % $p
        ~If $r==0
            ~Set prim=nein
            ~ExitLoop        
    ~Loop
    ~If $prim==ja
        ~AddTo pidx+=1
        >> P($pidx) = $n
        ~Set PZ($pidx)=$n
    ~AddTo n+=1
~Loop
\end{verbatim}

Die �u�eren Do-While-Schleife wird durchlaufen, solange die Variable \verb|n| kleiner als 20 ist. Der Wert von \verb|n| wird dabei jedes Mal vor Schleifenende um eins erh�ht, was vom Prinzip her einer For-Schleife entspricht. Gestartet wird mit \verb|n=2|. Auf die Formulierung von Bedingungen gehen wir in Abschnitt~\ref{IfUndBedingungen} n�her ein.

Die innere For-Each-Schleife durchl�uft eine Menge \verb|PZ| unter Verwendung der Variablen \verb|p|. \verb|PZ| wird hierbei als eindimensionales Feld interpretiert, das die fortlaufend nummerierten Variablen \verb|PZ(1)|, \verb|PZ(2)|, \verb|PZ(3)|, usw.\ enth�lt.
Beim ersten Durchlauf der �u�eren Schleife ist die Variable \verb|PZ(1)| noch nicht definiert und die Menge \verb|PZ| daher leer, d.h.\ die innere Schleife wird nicht durchlaufen.

F�r jede Zahl $n$, f�r die nach Durchlauf der inneren Schleife die Variable \verb|prim| noch auf \verb|ja| steht, wird die Menge der gefundenen Primzahlen \verb|PZ| um $n$ erweitert, weil $n$ in diesem Fall durch keine der bisher gefundenen Primzahlen teilbar ist. Ist $n$ dagegen durch eine dieser Primzahlen $p$ teilbar, also wenn der Rest bei Division durch $p$ null ergibt, dann wird die Variable \verb|prim| auf \verb|nein| gesetzt und die innere Schleife mit dem Befehl \fcmd{ExitLoop} verlassen, weil auch die folgenden Durchl�ufe das Ergebnis nicht mehr ver�ndern w�rden.

Das Programm geht also recht effizient vor und versucht soweit es geht, unn�tige Schleifendurchl�ufe zu vermeiden. Dies ist dahingehend kein Fehler, da die Ausf�hrung von \FPZb-Skripten, verglichen mit einer kompilierbaren Programmiersprache wie C\#, sehr langsam ist.

Die Programmierung von Schleifen bringt immer auch mit sich, dass man versehentlich eine Endlosschleife durchl�uft, z.B.\ wenn man das Hochz�hlen der Variable vergisst und so die Abbruchbedingung nie erreicht wird. Um dies und auch Endlosrekursionen fr�hzeitig zu erkennen, werden alle Durchl�ufe von Schleifen im Hintergrund mitgez�hlt, und die Ausf�hrung mit dem Fehler \emph{Maximale Anzahl Schleifendurchl�ufe erreicht} beendet, wenn der intern gesetzte Stop-Z�hler\index{Stop-Z�hler} f�r Schleifen oder die Gesamtzahl an durchlaufenen Befehlen erreicht wird. Mit dem Befehl \fcmd{SetStopCounter} lassen sich diese beiden Werte hochsetzen, wenn man absch�tzen kann, dass eine regul�re Skriptanwendung so viele Durchl�ufe und Befehlsausf�hrungen mit sich bringt. Raffinierte Berechnungen mit vielen Iterationen und Rekursionen sind zwar m�glich, liegen aber eher am Rand der Wohlf�hlzone von \FPZ.

Ein h�ufigerer Anwendungsfall, insbesondere wenn man Skripte schreibt, die Programmcode generieren, ist die Abfrage einer Sequenz von Elementen, z.B.  die Variablennamen f�r die Parameter einer Funktion. Im Beispiel unten wird �ber eine Schleife wiederholt abgefragt, ob es noch einen weiteren Parameter gibt, und wenn ja, nach dem dazugeh�renden Variablenname gefragt. 

\Beispiel{Abfrageschleife}\label{BspAbfrageschleife}
\begin{verbatim}
?V1: Hat die Funktion Parameter?
    #j: Ja
        ~Set i=1
        ~DoWhile $V$i==j
            ~Input B$i: Variable f�r den $i. Parameter:
            ~Set Vars($i)=$B$i
            ~AddTo i+=1
            ?V$i: Hat die Funktion einen $i. Parameter?
                #j: Ja            
                #n: Nein
            ~Execute
        ~Loop
        ~ForEach v in Vars
            @Variablen >> $v
        ~Loop
    #n: Nein
\end{verbatim}

Die Antworten und Eingaben werden in fortlaufend nummerierten Schl�sseln verwaltet und die Ausf�hrung nach jedem Schleifendurchlauf angehalten. 

\subsection{If-Abfragen und Bedingungen}\label{IfUndBedingungen}
Zu den wichtigsten Steuermechanismen innerhalb eines Programms geh�rt ganz klar die M�glichkeit, Bedingungen\index{Bedingungen} auszuwerten und Fallunterscheidungen\index{Fallunterscheidungen} abzubilden.  Der dazugeh�rende Befand lautet \fcmd{If} und kann kombiniert werden mit den Befehlen \fcmd{ElseIf} und \fcmd{Else}. Er ist so elementar, dass wir ihn schon in zahlreichen Beispielen verwendet haben, ohne n�her darauf einzugehen.

\fcmd{If} f�hrt den darunter einger�ckten Skriptcode genau dann aus, wenn die hinter dem \fcmd{If}-Befehl stehende Bedingung erf�llt ist. Mit \fcmd{ElseIf} kann eine weitere Bedingung mit einem weiteren Codeblock angeschlossen werden, die aber nur dann ausgewertet, bzw.\ ausgef�hrt werden, wenn keine der davorstehenden Bedingungen erf�llt ist. Auf diese Weise k�nnen beliebig viele ElseIf-Bedingungen aneinandergereiht werden.�Abschlie�end kann mit \fcmd{Else} noch ein Codeblock angeh�ngt werden, der nur dann ausgef�hrt wird, wenn keiner der davorstehenden Bedingungen erf�llt wurde.

\Beispiel{Bewertungsschema}\label{BspBewertungsschema}
\begin{verbatim}
~Input P: Wie hoch ist die Punktezahl (0-20)?
~Execute
~If $P>=18
    ~Set Note=sehr gut
~ElseIf $P>=14
    ~Set Note=gut
~ElseIf $P>=10
    ~Set Note=befriedigend
~ElseIf $P>=5
    ~Set Note=ausreichend
~Else
    ~Set Note=ungen�gend
>> Ergebnis: $P Punkte (Note $Note)
\end{verbatim}

Ein klassischer Anwendungsfall f�r mehrstufige If-Abfragen ist die Anwendung eines Bewertungsschemas\index{Bewertungsschema}, bei der f�r einen vorgegebenen Wert das dazu passende Intervall gesucht werden muss, mit dem diesem Wert dann ein Zielwert zugeordnet wird. Im Beispiel oben wird so einem Punktewert von 0 bis 20 eine Note in Textform zugeordnet.

Die Bedingungen im Beispiel sind einfache Gr��er-gleich-Vergleiche\index{Vergleiche}, aber es k�nnen auch komplexere Ausdr�cke formuliert werden. Die Bedingung muss dabei in der disjunktiven Normalform\index{disjunktive Normalform} angegeben werden, also als Oder-Verkn�pfung\index{Oder-Verkn�pfung} (\verb=||=) von Und-Verkn�pfungen\index{Und-Verkn�pfung} (\verb|&&|), wobei keine Klammerung notwendig ist.
Als Literale sind die Konstanten \verb|1| und \cidxfrag{true} (wahr\index{wahr}), sowie \verb|0| und \cidxfrag{false} (falsch\index{falsch}) verwendbar, sowie f�r Zeichenketten $s$ und $t$, Zahlen $x$ und $y$ und Variablen $v$ die folgenden Ausdr�cke zul�ssig: \verb|$s==$t| ($s$ ist gleich $t$), \verb|$s!=$t| ($s$ ist ungleich $t$), \verb|$x<>$y| ($x$ ist ungleich $y$), \verb|$x<$y| ($x$ ist kleiner als $y$), \verb|$x<=$y| ($x$ ist kleiner oder gleich $y$), \verb|$x>$y| ($x$ ist gr��er als $y$), \verb|$x>=$y| ($x$ ist gr��er oder gleich $y$), \verb|$s~$t| ($s$ enth�lt $t$), \verb|$s!~$t| ($s$ enth�lt $t$ nicht), \verb|?$v| ($v$ ist gesetzt), \verb|!?$v| ($v$ ist nicht gesetzt). Siehe auch Kernelement~\ref{defBedingungen}.

\subsection{Funktionen}\label{Funktionen}
Viele spannende und m�chtige Algorithmen gr�nden auf Funktionen\index{Funktionen}, die sich rekursiv selbst aufrufen. Es gibt Programmiersprachen wie Scheme, die einen dazu zwingen, so gut wie jede Problemstellung mittels Rekursion\index{Rekursion} zu l�sen. Auch \FPZ\ erlaubt die Definition von Funktionen und rekursive Aufrufe, obgleich wie schon in fr�heren Abschnitten erw�hnt, dort nicht unbedingt der Schwerpunkt dieser Skriptsprache liegt.

Das nachfolgende Beispiel zeigt den Klassiker unter den rekursiv l�sbaren Aufgaben, n�mlich die T�rme von Hanoi\index{T�rme von Hanoi}. Dabei geht es darum, einen Turm von mehreren absteigend gro�en Schreiben von einer Stange unter Zuhilfenahme einer zweiten Stange auf eine dritte Stange umzustapeln, wobei jeweils nur eine Scheibe bewegt werden darf und niemals eine gr��ere Scheibe auf einer kleineren Scheibe abgelegt werden darf.

\Beispiel{T�rme von Hanoi}\label{BspTuermeVonHanoi}
\begin{verbatim}
~Set MoveFromS=A
~Set MoveToS=C
~Set MoveCountS=4
~GoSub Move; BaseKey=S
~End

// Bewegt n Scheiben von Stange a nach b
// MoveFrom$BaseKey:  Ausgangsstange a (A, B, C)
// MoveTo$BaseKey:    Zielstange b (A, B, C)
// MoveCount$BaseKey: Anzahl Scheiben n
~DefineSub Move
    ~If $MoveCount$BaseKey>0
        ~Set Q1$BaseKey=$MoveFrom$BaseKey
        ~Set Q2$BaseKey=$MoveTo$BaseKey
        ~Set Q3$BaseKey=ABC
        ~Replace Q3$BaseKey=$Q3$BaseKey|$Q1$BaseKey->
        ~Replace Q3$BaseKey=$Q3$BaseKey|$Q2$BaseKey->
        ~Set MC$BaseKey=$MoveCount$BaseKey
        ~AddTo MC$BaseKey+=-1

        ~Set MoveFrom$BaseKeyA=$Q1$BaseKey
        ~Set MoveTo$BaseKeyA=$Q3$BaseKey
        ~Set MoveCount$BaseKeyA=$MC$BaseKey
        ~GoSub Move; BaseKey=$BaseKeyA
        
        @L�sung >> Von $Q1$BaseKey nach $Q2$BaseKey
            __ (BaseKey=$BaseKey)
        
        ~Set MoveFrom$BaseKeyB=$Q3$BaseKey
        ~Set MoveTo$BaseKeyB=$Q2$BaseKey
        ~Set MoveCount$BaseKeyB=$MC$BaseKey
        ~GoSub Move; BaseKey=$BaseKeyB
~Return
\end{verbatim}

Der L�sungsansatz besteht darin, das Umstapeln eines Turms mit $n$ Scheiben von Stange $a$ nach Stange $b$ in drei Schritte zu zerlegen: Zun�chst stapelt man den oberen Teil des Turms aus $n-1$ Schreiben um auf die freie Stange $c$, bewegt dann die $n$te Schreibe auf die Zielstange $b$ und stapelt dann wieder denselben Turm von Stange $c$ auf Stange $b$ oben auf diese Scheibe drauf. Auf diese Weise hat man das Problem auf das Umstapeln eines kleineren Turms zur�ckgef�hrt und kann das wiederholt anwenden, bis man schlie�lich bei einem Turm mit $0$ Scheiben angekommen ist, bei dem nichts mehr zu tun ist.

Die Definition der Funktionen erfolgt standardm��ig am Ende des Skriptes. Ihr Code wird nur durchlaufen, wenn sie aufgerufen werden, trotzdem wird hier zur besseren �bersichtlichkeit die Ausf�hrung vorher mit dem Befehl \fcmd{End} beendet. Jede Funktion wird mit dem Befehl \fcmd{DefineSub} eingeleitet und mit dem Befehl \fcmd{Return} beendet. Der dazwischen liegende, einger�ckte Codeblock wird bei jedem Aufruf durchlaufen. Ein Aufruf wird mit dem Befehl \fcmd{GoSub} in die Wege geleitet, bei dem auch der \cidxfrag{BaseKey}-Wert �bergeben werden kann. Nach dem Aufruf wird das Skript in der n�chsten Zeile fortgesetzt.

In jedem Schritt muss zun�chst die Stange herausgefunden werden, die aktuell weder Ausgangspunkt, noch Ziel der Umschichtung darstellt und die als Zwischenablage verwendet werden kann. Das wird dadurch bewerkstelligt, dass man im Ausdruck \verb|ABC| die Kennungen der beiden gegebenen Stangen mit dem \fcmd{Replace}-Befehl entfernt.

Da es in \FPZ\ keine lokalen Variablenbereiche gibt, muss man die Variablen, die f�r einen bestimmten Funktionsaufruf erhalten bleiben sollen, explizit von den anderen Variablen abgrenzen. Daf�r kann �ber die Sondervariable \cidxfrag{BaseKey} eine Zahl oder Zeichenfolge �bergeben werden, die f�r jeden Funktionsaufruf getrennt verwaltet wird und die in den innerhalb der Funktion genutzten Variablen als Teil der Bezeichnung eingebaut werden kann, um auch diese zu trennen. In gleicherweise k�nnen auch die Funktionsargumente vor dem Funktionsaufruf als Variablen definiert werden, die sich auf den �bergebenen \cidxfrag{BaseKey}-Wert beziehen. Im Beispiel wird vor dem ersten Aufruf im Skript \verb|MoveFromS=A| als einer von drei Funktionsparametern gesetzt und beim Funktionsaufruf wird \verb|BaseKey=S| �bergeben. Damit wird innerhalb der Funktion \verb|$MoveFrom$BaseKey| zu \verb|$MoveFromS| und schlie�lich zu \verb|A| ersetzt. Bei den rekursiven Aufrufen wird der \cidxfrag{BaseKey}-Wert dann jeweils um die Buchstaben A und B erweitert, so dass am Ende Funktionsparameter zu 15 verschiedenen \cidxfrag{BaseKey}-Werten verwaltet werden.

Der Umgang mit rekursiven Funktionen ist in \FPZ\ zugegebenerma�en etwas gew�hnungsbed�rftig, aber er l�sst sich mit einigen Debug-Ausgaben und etwas �bung doch gut hinbekommen. Analog zu den Variablen lassen sich auch die Schl�ssel f�r Eingaben mithilfe der \cidxfrag{BaseKey}-Variablen definieren, so dass in jeder Rekursion auch Benutzerinteraktionen stattfinden k�nnen. Auf diese Weise kann zum Beispiel ein Sortieralgorithmus wie Mergesort implementiert werden, der den einzelnen Vergleich zweier Elemente an den Anwender weitergibt. Damit kann etwa die Sortierung von Projekten oder Aufgaben nach Priorit�t auf Einzelvergleiche heruntergebrochen werden, die sich oft mit Bauchgef�hl ganz gut festlegen lassen.

Nat�rlich kann man Funktionen auch ohne Rekursion nutzen. In diesem Fall ben�tigt man keine \cidxfrag{BaseKey}-Variable und kann sowohl die Funktionsargumente, als auch die R�ckgabewerte ganz normal in Variablen �bergeben.

\section{Verarbeitung von Texten}
Der Zweck eines Skriptes ist es, eine Ausgabe zu erzeugen und der Zweck von Texteingaben liegt meist darin, diese Texte in irgendeiner Form in der Ausgabe einzuarbeiten. Sofern die Eingabe unver�ndert in der Ausgabe erscheinen soll, reicht es, einfach die entsprechende Variable in der Ausgabe zu verwenden. In manchen F�llen m�chte man jedoch einen Text in einer bestimmten Form abwandeln oder entsprechend einer vorgegebenen Struktur zerlegen. Mit den in diesem Abschnitt beschriebenen Befehlen lassen sich die meisten Textver�nderungen hinbekommen.

\subsection{Texte aufteilen und neu kombinieren}\label{TexteAufteilen}
Wie wir schon in Beispiel~\ref{BspBugPlaner} gesehen haben, Kann man mit Schleifen sehr gut Listen durchlaufen und eine Menge von Elementen bearbeiten. In diesem Beispiel wurde auch schon der \fcmd{Split}-Befehl verwendet, um aus einer kommagetrennten Aufz�hlung eine solche indizierte Liste zu erstellen, die mit dem \fcmd{ForEach}-Befehl durchlaufen werden kann. In der gleichen Weise kann auch eine ganze Aufz�hlung durch eine einzelne \fcmd{Input}-Abfrage als Benutzereingabe entgegengenommen, und in eine Liste umgewandelt werden. 

Im nachfolgenden Beispiel wird sogar eine verschachtelte Liste, also eine Liste von Listen\index{Listen!von Listen} als Eingabe abgefragt und verarbeitet. Der \fcmd{Split}-Befehl funktioniert so, dass er den �bergebenen Text anhand des angegebenen Trennzeichens aufteilt und in einer indizierten Liste speichert. 

\Beispiel{Kombinator}\label{BspKombinator}
\begin{verbatim}
~Input C: Kombinierbare Eigenschaften
    __ (z.B.: A,B,C : X,Y,Z)
~Split A=$C|:
~Set aidx=1
~Set QAnz(1)=1
~Set Q(1)(1)=#Start#
~ForEach a in A
    ~Set aprev=$aidx
    ~AddTo aidx+=1
    ~Split W=$a|,
    ~Set widx=0
    ~ForEach q in Q($aprev)
        ~ForEach w in W
            ~AddTo widx+=1
            ~Replace Q($aidx)($widx)=$q - $w|#Start# -->
        ~Loop
    ~Loop
~Loop
~ForEach q in Q($aidx)
    @Kombinationen >> $q
~Loop
\end{verbatim}

Der \emph{Kombinator}\index{Kombinator} nimmt eine Liste von kombinierbaren Eigenschaften entgegen, die jeweils verschiedene Werte annehmen k�nnen, z.B.\ \emph{Farbe gleich blau, Farbe gleich gr�n} und \emph{Material gleich Holz, Material gleich Metall}. Die Werte jeder Eigenschaft werden dabei durch Kommata getrennt, die einzelnen Wertlisten wiederum durch einen Doppelpunkt. Bei Ausf�hrung des Skriptes werden alle Kombinationen aus den Werten jeder Eigenschaft gebildet und aufgelistet. In unserem Beispiel ergibt das vier Kombinationen, angefangen mit \emph{Farbe gleich blau und Material gleich Holz}.

Der Umgang mit solchen kombinatorischen Auf\-lis\-tungen\index{kombinatorischen Auf\-lis\-tungen} ist in der Softwareentwicklung wichtiger, als man denkt, und ein gro�er Teil aller Fehler bei der Implementierung von Programmfunktionalit�t geht darauf zur�ck, dass F�lle �bersehen oder vergessen werden. Dies passiert oft schon beim Entwurf der Userstory, in der haupts�chlich die Geradeaus-F�lle betrachtet werden, und F�lle, die keinem normalen Anwendungsfall entsprechen, oft au�er Acht gelassen werden. Wenn solche F�lle jedoch theoretisch m�glich sind, muss auch die Software damit in irgendeiner Form umgehen k�nnen, und damit sie das kann, m�ssen auch diese F�lle von Anfang an betrachtet und behandelt werden. 

Gerade bei der Umsetzung von Formeln ist es ungeheuer wichtig, dass der komplette zul�ssige Definitionsbereich bei Implementierung und Test abgedeckt sind, und dass systematisch sichergestellt wird, dass alle nicht zul�ssigen Eingaben und F�lle sauber und f�r den Benutzer transparent durch die Programmoberfl�che abgefangen werden. Eine solche Systematik bekommt man am besten dadurch hin, dass man zun�chst die vorhandenen F�lle in ihre Eigenschaften zerlegt und diese dann mit Hilfe so eines Werkzeugs zu allen Fallunterscheidungen kombiniert, was allerdings recht umfangreich werden kann. 

In vielen F�llen h�ngen weitere Teilkombinationen nur an einzelnen Werten und die Bestimmung aller Kombinationen entspricht dem ausmultiplizieren eines Terms mit  den Operationen $+$ und $\cdot$ und einer beliebigen Klammerung, nur dass die Elemente Textfragmente sind und keine Zahlen. Das folgende Beispiel zeigt, wie die Summanden des Ergebnisausdruckes am Ende die einzelnen Kombinationen bilden:
\[
\begin{split}
K &= (\text{Material Holz} + \text{Material Metall})\\
    &\quad\cdot (\text{Farbe blau} + \text{Farbe gr�n}\cdot (\text{einfarbig} + \text{mit gelben Streifen}))\\
&= \text{Material Holz}\cdot\text{Farbe blau}\\
&\quad+ \text{Material Holz}\cdot\text{Farbe gr�n}\cdot\text{einfarbig}\\
&\quad+ \text{Material Holz}\cdot\text{Farbe gr�n}\cdot\text{mit gelben Streifen}\\
&\quad+ \text{Material Metall}\cdot\text{Farbe blau}\\
&\quad+ \text{Material Metall}\cdot\text{Farbe gr�n}\cdot\text{einfarbig}\\
&\quad+ \text{Material Metall}\cdot\text{Farbe gr�n}\cdot\text{mit gelben Streifen}
\end{split}
\]
Mit der in Abschnitt~\ref{RegUndDaten} beschriebenen Funktion zur Anwendung regul�rer Ausdr�cke lassen sich solche Klammerausdr�cke wiederholt in einem Eingabetext suchen und ausmultiplizieren, bis am Ende der gesamte Ausdruck expandiert ist. Ebenso l�sst sich nach definierten Textmarken suchen, die, wenn sie in einer Kombination aufeinandertreffen, diese aus dem Ergebnis entfernen. Zus�tzlich l�sst sich auch noch eine Textmarke f�r eine Ergebnisgruppierung definieren, so dass die Kombinationen als zweistufige Liste ausgegeben werden.


\subsection{Ersetzungen und Zufallsgenerierung}\label{ErsetzungenUndZufallsgenerierung}
Den Ersetzungsbefehl \fcmd{Replace} haben wir schon in vielen Skripten gesehen, etwa in Beispiel~\ref{BspTextersetzung}. Die Ersetzung\index{Ersetzung} von Texten in anderen Texten ist f�r viele Zwecke einsetzbar. Im nachfolgende Beispiel wird die Textersetzung dazu verwendet um systematisch einen Text aufzubauen. Hierbei kommt auch der Befehl \fcmd{Random} zum Einsatz, mit dem eine Zufallszahl\index{Zufallszahl} in einem Zahlenbereich generiert werden kann.

\Beispiel{Passwortgenerator}\label{Passwortgenerator}
\begin{verbatim}
~Set L=12
~Random AnzZiffer=1..2
~Random AnzGross=1..2
~Random AnzSonder=1..2
~Split SZ=33,35,36,37,38,42,43,47,64|,
~Set M="
~Set i=0
~DoWhile $i<$L
    ~Random r=0..$i
    ~AddTo i+=1
    ~If $r==0
        ~Set M="+$i$M
    ~Else
        ~Replace M=$M|"+$r"->"+$r"+$i"
~Loop
~Set i=0
~DoWhile $i<$L
    ~AddTo i+=1
    ~If $AnzZiffer>0
        ~Random z=48..57
        ~AddTo AnzZiffer+=-1
    ~ElseIf $AnzGross>0
        ~Random z=65..90
        ~AddTo AnzGross+=-1
    ~ElseIf $AnzSonder>0
        ~Random s=1..$SZ(0)
        ~Set z=$SZ($s)
        ~AddTo AnzSonder+=-1
    ~Else
        ~Random z=97..122
    ~Replace M=$M|"+$i"->"$Chr($z)"
~Loop
~Replace PW=$M|"->
@Ausgabe >>* Passwortvorschlag
    >|$PW
\end{verbatim}

Das Beispiel zeigt einen Passwortgenerator\index{Passwortgenerator}, der f�r eine vorgegebene L�nge ein Passwort aus zuf�lligen Zeichen erzeugt. Die von den meisten Passwortrichtlinien geforderten Zusatzbedingungen, dass das Passwort auch Ziffern, Gro�buchstaben und Sonderzeichen enthalten muss, sind ebenfalls ber�cksichtigt. Da diese Zeichen zum Eintippen oftmals l�stig sind, kann ihre Anzahl hier beschr�nkt werden. Ebenso kann die Menge der Sonderzeichen festgelegt werden, die zum Einsatz kommen soll.

Im ersten Teil des Skriptes wird eine Zeichenkette aus den einzelnen Zahlen $1$ bis $L$ f�r die gew�nschte L�nge $L$ des Passworts und einigen Trennzeichen erstellt. Diese wird bereits bei der Erstellung permutiert, indem das Einf�gen an einer beliebigen Stelle mit dem Ersetzungsbefehl durchgef�hrt wird. Im zweiten Teil wird f�r jede Zahl von $1$ bis $L$ ein Zeichen aus den verschiedenen Zeichengruppe per Zufallsgenerator bestimmt und die Zahl durch das entsprechende Zeichen ersetzt. Hierbei wird der \cidxvar{Chr(...)}-Befehl verwendet der das Zeichen mit dem dazu geh�renden ANSI-Code\index{ANSI-Code} erzeugt. Nacheinander wird so die geforderte Anzahl an Ziffern, Gro�buchstaben und Sonderzeichen eingef�gt. Die restlichen Zeichen werden mit Kleinbuchstaben aufgef�llt. Am Ende werden die Trennzeichen entfernt, sodass nur noch die Zeichen �brig bleiben, die das Passwort bilden.

F�r die Erzeugung der Sonderzeichen wird zun�chst eine Zufallszahl von $1$ bis zur L�nge der Liste \verb|SZ| generiert, wobei diese mit \verb|$SZ(0)| abgerufen wird. Mit dieser Zahl wird der entsprechende Listeneintrag ausgelesen, der dann in gleicher Weise wie die anderen Codes in ein Zeichen umgewandelt wird. Damit kann die Liste um weitere Codes erg�nzt werden, ohne dass an anderer Stelle die Anzahl der Liste angepasst werden muss.

\subsection{Regul�re Ausdr�cke und Datenabfragen}\label{RegUndDaten}
Dieser Abschnitt demonstriert zwei der wohl m�chtigsten Befehle im Besteckkasten von \FPZ: regul�re Ausdr�cke\index{regul�re Ausdr�cke} und Datenabfragen. Regul�re Ausdr�cke dienen dazu, Textmuster zu beschreiben, so dass man diese innerhalb von Texten auffinden kann, z.B.\ \emph{ein Gro�buchstabe gefolgt von zwei Ziffern} was mit \verb|[A-Z]\d{2}| beschrieben wird. Zur Formulierung von regul�ren Ausdr�cken gibt es zahlreiche Hilfsseiten und Online-Tools und nat�rlich auch die Chatbots der gro�en Sprachmodelle, denen man den gew�nschten Ausdruck umgangssprachlich beschreiben kann. Ein von mir gerne genutztes Tool zur �berpr�fung regul�rer Ausdr�cke ist \emph{https://regex101.com}. 

Datenabfragen\index{Datenabfragen} sind Zugriffe auf Dateien, die irgendwo im Skripte-Verzeichnis liegen, wie die Skripte, die sich �ber das Men� von \FPZ\ ausw�hlen lassen. Es ist zwar nicht m�glich, schreibend auf diese Dateien zuzugreifen, aber man kann sie auslesen. Entsprechend gibt es mit \fcmd{ForEachLine} einen noch nicht genannten Schleifen-Befehl, mit dem man �ber die Zeilen einer beliebigen Datei iterieren kann. Kombiniert man nun das Einlesen einer Datei mit den M�glichkeiten der regul�ren Ausdr�cke, so kann man mehr oder weniger jedes beliebige textbasierte Dateiformat verarbeiten. 

Das nachfolgende Beispiel ist ein einfaches Werkzeug zur Anwendung einer Terminologie\index{Terminologie} auf einen Eingabetext. Die Terminologie f�r einen bestimmten Kontext beschreibt die dort relevanten Begriffe und gibt an, welche der jeweils m�glichen Benennungen innerhalb des Kontextes zul�ssig sind und welche nicht. In der Regel einigt man sich f�r jeden Begriff auf eine zul�ssige Benennung und schlie�t mehrere andere aus. So entsteht mit der Zeit eine Terminologiedatenbank aus Begriffen, Benennungen, Definitionen, Beschreibungen und Ausschlussbegr�ndungen, die in einem sehr einfachen Fall so aussehen k�nnte:

\Beispiel{Terminologieliste.txt}
\begin{verbatim}
1    Passwort    ja  Zeichenfolge, deren Kenntnis Zugang...
1    Kennwort    nein    Weniger �blich.
2    Schaltfl�che    ja  Steuerelement, das bei Klick eine...
2    Button  nein    Anglizismus.
2    Knopf   nein    Wird eher mit Ger�ten assoziiert.
\end{verbatim}

Jede Zeile hat vier Eintr�ge, die durch Tabulatoren getrennt sind. In der ersten Spalte wird die ID des Begriffes angegeben, und da ein Begriff mehrere Benennungen haben kann, k�nnen auch mehrere Zeilen zur selben ID vorhanden sein. Die zweite Spalte gibt eine Benennung an. Die dritte Spalte enth�lt entweder den Inhalt \emph{ja} oder \emph{nein} und gibt an, ob die jeweilige Benennung zul�ssig ist oder nicht. Im ersten Fall gibt die vierte Spalte die Definition oder Beschreibung des Begriffs an, im zweiten Fall enth�lt sie die Begr�ndung, warum die entsprechende Benennung nicht verwendet werden soll. F�r die Verarbeitung ist es wichtig, dass die zul�ssige Benennung immer in der ersten Zeile f�r jede ID steht und dass alle Zeilen zur selben ID aufeinander folgen.

Der Terminologiepr�fer\index{Terminologiepr�fer} ist das dazu geh�rende Terminologie-Tool und sieht so aus:

\Beispiel{Terminologiepr�fer}\label{BspTerminologiepruefer}
\begin{verbatim}
~Input T:Text
@Eingabetext >>_ $T
~ForEachLine z in Terminologieliste.txt; NoFormat
    ~RegExMatch ti=$z|([0-9]*)\t([^\t]*)\t(ja|nein)\t(.*)
    ~If $ti(0)
        ~If $ti(3)==ja
            ~Set EmpfBen=$ti(2)
            ~Set DefBegr=$ti(4)
        ~If $T~$ti(2) && $ti(3)==ja
            @Erlaubte Benennungen >> $ti(2)
                > ID: $ti(1)
                > Definition: $ti(4)
        ~If $T~$ti(2) && $ti(3)==nein
            @Nicht erlaubte Benennungen >> $ti(2)
                > ID: $ti(1)
                > Empfohlen: $EmpfBen
                > Definition: $DefBegr
                > Begr�ndung: $ti(4)
~Loop
\end{verbatim}

Der Terminologiepr�fer durchl�uft alle Benennungen aus der Terminologiedatenbank und pr�ft, ob diese im eingegebenen Text vorhanden sind. Je nachdem ob es sich um eine zul�ssige Benennung handelt oder nicht, wird diese dann in einem entsprechenden Abschnitt aufgelistet. Bei einer nicht zul�ssigen Benennung wird die zum entsprechenden Begriff empfohlene Benennung samt Definition ausgegeben, ebenso die Begr�ndung warum die Benennung nicht verwendet werden soll. F�r die Eingabe \emph{nach Eingabe des Passworts muss der Button gedr�ckt werden} wird \emph{Passwort} als zul�ssige Benennung best�tigt und \emph{Button} mit Hinweis auf \emph{Schaltfl�che} und der Begr�ndung, dass es sich um einen Anglizismus handelt, als nicht zul�ssige Benennung abgelehnt. 

Der Befehl \fcmd{RegExMatch} liefert in der Ergebnisvariablen \verb|ti| eine Liste zur�ck. Das nullte Element \verb|ti(0)| gibt an, ob f�r den Ausdruck ein Treffer gefunden wurde, die nachfolgenden Elemente enthalten die Teiltexte des gefundenen Ausdrucks, die �ber die runden Klammern im Ausdruck als Gruppen angegeben wurden.

Das \cidxfrag{NoFormat}-Argument gibt an, dass die Datei nicht das Dateilisten-Format hat und daher alle Zeilen unver�ndert zur�ckgegeben werden sollen. Beispiele f�r Dateilisten findet man in Beispiel~\ref{BspBasisbefehle} und \ref{BspProgrammierung}.

Das Skript bietet noch viel Spielraum f�r Erweiterungen, beispielsweise k�nnten alle im Text vorhandenen gro�geschriebenen W�rter, die noch nicht gefunden wurden als Vorschl�ge f�r die Aufnahme in die Terminologie aufgelistet werden, direkt mit einem Link zu einem Terminologieerfassungskript, das die ben�tigten Informationen abfragt und daraus die entsprechenden Eintr�ge f�r die Terminologiedatenbank erstellt.

\subsection{Parametrisierte Links generieren}\label{ParametrisierteLinks}
Der gro�e Vorteil von Webanwendungen liegt darin, dass man innerhalb der Programmoberfl�che �berall Links auf andere Programmbereiche oder sogar andere Anwendungen einbauen kann. Diese lassen sich dann bequem in neuen Tabs �ffnen, um einen zwischengelagerten Arbeitsschritt auszuf�hren und danach kann man zu seiner urspr�nglichen Aufgabe zur�ckkehren. Umgekehrt kann eine Webanwendung wiederum �ber einen parametrisierten Link\index{Link!parametrisiert} aufgerufen werden, dem man einzelne Begriffe oder auch kleinere Texte als Argumente �bergeben kann.

Wie man Links in die Ausgabe einbauen kann, haben wir in Abschnitt~\ref{LinksInline} schon gesehen. Um jedoch beliebigen Text als Parameter in einem Link einzuf�gen, muss man diesen in einer speziellen Form kodieren. Hierf�r gibt es in \FPZ\ den Befehl \fcmd{UrlEncode}. Das nachfolgende Beispiel zeigt die Verwendung von URL-Parametern\index{URL-Parametern} am Beispiel eines Suchbegriffs, den man als Parameter an eine Suchmaschine �bergeben kann, hier am Beispiel von Google und Bing.

\Beispiel{Suchlink-E-Mail}\label{BspSuchlink}
\begin{verbatim}
~Input S:Suchbegriff
~UrlEncode uS=$S
@Suchanfragen >>
    ~AddLink https://www.google.com/search?q=$uS |
    __Suche "$S" bei Google
>>
    ~AddLink https://www.bing.com/search?q=$uS |
    __ Suche "$S" bei Bing
~UrlEncode usub=Suchlinks zu "$S"
~UrlEncode ubody=Hallo,$CRLF
    __hier sind zwei Suchlinks zu "$S":$CRLF
    __https://www.google.com/search?q=$uS$CRLF
    __https://www.bing.com/search?q=$uS$CRLF
    __MfG
>> Suchanfragen
    ~AddLink mailto:a@bc.de?subject=$usub&body=$ubody |
    __ per E-Mail verschicken
\end{verbatim}

Zuerst wird ein Suchbegriff abgefragt. Dieser wird mittels \fcmd{UrlEncode} in die Variable \verb|uS| in codierter Form abgelegt. F�r den mailto-Link\index{mailto-Link} werden nacheinander Betreff (\emph{subject}) und Inhalt (\emph{body}) der E-Mail\index{E-Mail} in gleicher Weise in die beiden Variablen \verb|usub| und \verb|ubody| codiert und am Ende in den Link eingebaut.

Das Beispiel zeigt zum einen den Einsatz von Zeilenumbr�che\index{Zeilenumbr�che}, die man mit dem Befehl \cidxvar{CRLF} erzeugen kann, was die Zeichenkombination \emph{Carriage Return + Line Feed} ausgibt, und zum anderen, dass auch kodierte Texte wie die Suchlinks wieder Bestandteil von Codierungen sein k�nnen. 

Analog zu \emph{mailto} verf�gen auch manche Ticketsysteme �ber Adressen, mit denen sich Vorg�nge �ber parametrisierte Aufrufe vollst�ndig vorbereiten lassen, so dass alle relevanten Felder ausgef�llt sind. In gleicher Weise k�nnen auch spezialisierte Suchen im selben Ticketsystem in Form von Links bereitgestellt werden, so dass direkt innerhalb der richtigen Projekt- oder Aufgabenkategorie gesucht wird. Ein klassisches Beispiel ist die Einrichtung einer Metasuche\index{Metasuche}, die f�r einen eingegebenen Suchbegriff �hnlich wie das Beispiel oben verschiedene Links in den verschiedenen Firmensystemen wie Wiki, Programmdokumentation und dem Ticketsystem bereitstellt, und in letzterem noch zus�tzlich unter verschiedenen Suchparametern.

Oft m�ssen solche Tickets oder andere Vorg�nge durch Personen au�erhalb des Fachbereichs erstellt werden, die sich besonders schwer damit tun alle Eigenschaften in der f�r die Bearbeitung gew�nschten Art und Weise zu setzen. Wenn hier mit einem Skript soweit unterst�tzt werden kann, dass am Ende die Erstellung des Vorgangs direkt �ber einen Link m�glich ist, hat man viel gewonnen. Idealerweise kombiniert man so ein Erstellungsskript mit einem Expertensystem (siehe Abschnitt~\ref{Expertensystem}) und der angeleiteten Suche nach eventuell schon vorhandenen Tickets zum Thema.

Eine weitere Anwendungsm�glichkeit f�r parametrisierte Links sind nat�rlich die Aufrufe anderer \FPZb-Skripte. Wie schon beschrieben, m�ssen f�r komplexe Aufgaben oftmals Werte und Informationen abgefragt werden, die sich nicht so ohne weiteres herausfinden lassen, und f�r deren Bestimmung eigene Anleitungen hilfreich w�ren. Diese k�nnten dann ebenfalls als Skripte bereitgestellt und �ber die Hilfezeilen der entsprechenden Felder mit Parametern verlinkt werden.

\section{Systementwicklung}
In diesem Abschnitt kommen wir endlich an die Kernt�tigkeit der professionellen Softwareentwicklung, also das Schreiben von Programmcode. Auch wenn man den Eindruck bekommt, dass hier die KI-Systeme, insbesondere die LLM-Chatbots zuk�nftig alles dominieren werden, gibt es trotzdem noch F�lle, in denen das nicht so bald der Fall sein wird.

\subsection{Zielsetzung der Systementwicklung}
Die �berschrift lautet bewusst Systementwicklung\index{Systementwicklung}, da viele Softwareh�user �ber Jahre hinweg ein umfangreiches System aufbauen und weiterentwickeln, dass auf einem ebenso umfangreichen und durchdachten Framework aufbaut, das wiederum durchg�ngig auf einem Designsystem gr�ndet. Jeder Teil der Anwendung sieht damit einheitlich aus, folgt den gleichen Bedienprinzipien sowie der f�r den Dom�nenkontext festgelegten Terminologie.

Die Erweiterung des Frameworks f�r ein solches System ist aufwendig, da all diese Punkte ber�cksichtigt werden m�ssen. Die Erweiterung der Anwendung auf Basis des vorhandenen Frameworks ist dagegen recht einfach, oder sollte es zumindest sein, da in diesem Fall keine technischen Probleme mehr gel�st werden m�ssen, sondern nur noch die richtigen Komponenten in der richtigen Art und Weise zusammengebaut, und nach den Vorgaben des Product Owners konfiguriert werden m�ssen. Und je besser das Framework ist, desto weniger Spielraum gibt es beim Anwenden und damit auch weniger Fragen, weniger Abweichungen und weniger Fehler.

Wer bis hierher durchgehalten hat, wird zustimmen, dass der Aufbau neuer Programmbereiche auf Basis eines vorhandenen Frameworks geradezu pr�destiniert ist f�r den Einsatz von \FPZb-Skripten und dass diese Methode hier sowohl von der Effizienz als auch von der Qualit�t her punkten kann.

Zielsetzung ist ein System aus Anleitungen und Unterst�tzungssystemen, das alle Komponenten des Frameworks in jeweils allen Phasen der Softwareentwicklung repr�sentiert, angefangen bei der Auswahl der Komponenten zur Umsetzung bestimmter Anforderungen, �ber die Verwendung im Programmcode, bis hin zur Testphase. Alle Abh�ngigkeiten, die sich zwischen den verschiedenen Entwicklungsphasen ergeben, sollten weitestgehend in den Skripten abgebildet sein und automatisch Ber�cksichtigung finden.

\subsection{Entwicklungsanleitungen}\label{Entwicklungsanleitungen}
Entwicklungsanleitungen\index{Entwicklungsanleitungen} sind, wie der Begriff schon sagt, Anleitungen f�r die Entwickler. Sie beschreiben, wie Anforderungen mit Hilfe der zur Verf�gung stehenden Mittel umgesetzt werden k�nnen oder sollen. 

Idealerweise gibt es f�r jede wiederkehrende Entwicklung eine dokumentierte Best Practice\index{Best Practice}, die vorgibt, wie diese umzusetzen ist, so dass unabh�ngig vom ausf�hrenden Entwickler immer dieselbe L�sung herauskommen sollte, die sowohl Qualit�t, als auch Wartbarkeit garantiert.�Die Anleitung f�r ein wiederkehrendes Entwicklungsmuster ist ein guter Ort, um so eine Best Practice direkt zu hinterlegen und aktuell zu halten, denn die Tatsache, dass einem eine solche Anleitung die Arbeit leichter macht, stellt auch praktisch sicher, dass diese verwendet wird und damit die Best Practice Anwendung findet.

Auch wenn die eigentliche Arbeit an einer Entwicklung beim Product Owner beginnt, ist es sinnvoll, beim Aufbau solcher Anleitungen bei der Entwicklung zu beginnen. Zum einen sind die Entwickler diejenigen, die am wenigsten Ber�hrungs�ngste mit der Skriptentwicklung haben d�rften, zum anderen ist durch das Vorhandensein von Framework-Komponenten schon eine sehr weit fortgeschrittene Systematisierung der wiederkehrenden Entwicklungen vorhanden, die nur noch in Skriptform gebracht werden muss.

Die Vorgehensweise hierf�r ist relativ einfach und folgt dem Top-Down-Ansatz. Idealerweise nutzt man einen konkreten Anwendungsfall, um eine Anleitung zu erstellen, also die angeforderte Umsetzung einer entsprechenden Aufgabe. Ohne Anleitung w�rde man sich hierbei vermutlich an einem fr�heren Beispiel derselben Art orientieren und versuchen, daraus mit der Methode \emph{kopieren, einf�gen, anpassen} die gew�nschte neue Instanz zu erzeugen. Die Erstellung der Anleitung geht im Prinzip denselben Weg.

Im ersten Schritt erzeugt man zun�chst ein Skript, dass nur 1-zu-1 den Programmcode erzeugt, der das fr�here Beispiel ausmacht, also ohne Verallgemeinerungen und ohne Platzhalter. Nun geht man diesen St�ck f�r St�ck durch und pr�ft, was man f�r die geforderte Instanz anpassen w�rde, z.B.\ kontextbezogene Benennungen, bestimmte IDs, �nderungen aufgrund anderer Anforderungen oder Voraussetzungen.

Die Idee ist, diese Dinge nun so zu systematisieren, dass sie durch das Skript auf jeden gegebenen Anwendungsfall hin angepasst werden k�nnen, und der Programmcode den Anwendungsfall am Ende so gut wie m�glich vollst�ndig abbildet. Bezeichnungen, sowie das Vorhandensein bestimmter Voraussetzungen oder Anforderungen k�nnen abgefragt werden und ben�tigte IDs k�nnen nach Anleitung ermittelt oder mittels Zufallsgenerator generiert werden. St�ck f�r St�ck wird so aus einem Skript, das konkreten Programmcode ausgibt, ein Skript, das die vorliegende Situation abfragt, und f�r diese den passenden Programmcode samt Einbauanleitung generiert.

Die gro�e Kunst dabei ist es, die Abfrage der Situation so niederschwellig wie m�glich zu gestalten, so dass die gestellten Fragen und geforderten Eingaben mit so wenig technischem Wissen wie m�glich get�tigt werden k�nnen. Wenn also eine Eigenschaft $x$ nur in bestimmten F�llen gesetzt werden kann, und je nach Situation einen von drei verschiedenen Werten annehmen kann, sollte zun�chst verstanden werden, was diesen Fall und diese Situationen ausmachen und was die entsprechenden Werte bedeuten. Die Fragestellung des Skriptes sollte auf diesem Wissen aufbauen. Man sieht, dass der Entwickler des Skriptes ein �hnlich tiefgreifendes Wissen ben�tigt, wie der Entwickler des Frameworks, zu dessen Verwendung es anleitet.

N�hern  wir uns dieser Vorgehensweise mit einem einfachen Beispiel. Entwickelt werden soll eine Archivierungsklasse f�r die Archivierung eines Dozenten und damit verbunden auch gleich die \FPZb-Anleitung f�r weitere Archivierungsklassen. Zur�ckgreifen k�nnen wir auf den Programmcode der Archivierungsklasse f�r Kurse, der schon geschrieben wurde. 

Wir beginnen also damit, eine Anleitung f�r genau diese Klasse zu schreiben, die dann wie folgt aussieht:

\Beispiel{Archivklasse 1}
\begin{verbatim}
@Anleitung >> Ordner Archivierer �ffnen.
@Anleitung >> Neue cs-Datei erstellen mit Bezeichnung
    >|KursArchivierer.cs
@Anleitung >> Folgenden Code einf�gen:
    >|// Klasse f�r die Archivierung von Kurs-Objekten
    >|public class KursArchivierer : BasisArchivierer<Kurs>
    >|{
    >|    // Kennung der Archivierer-Klasse
    >|    public override Guid ArchiviererKennung()
    >|    {
    >|        return Guid("14df603a-cb42-4560-bf6c-a740356532b6");
    >|    }
    >|
    >|    // Methode zum Archivieren eines Kurs-Objektes
    >|    public override void Archiviere(Kurs k)
    >|    {
    >|      var konflikt = PruefeAufKonflikte(k);
    >|        if (konflikt.vorhanden)
    >|        {
    >|            ALogger.Note($"Konflikt f�r Kurs '{k}', 
    >|                {konflikt.Details}");
    >|            return;
    >|        }
    >|        base.StarteArchivierung(k);
    >|    }
    >|    
    >|    // Konfliktpr�fung f�r die Kurs-Archivierung
    >|    private Konflikt PruefeAufKonflikte(Kurs k)
    >|    {
    >|        if (k.Status == KStatus.Aktiv)
    >|            return new Konflikt("Kurs ist noch aktiv");
    >|        
    >|        // Ansonsten
    >|        return Konflikt.KeinKonflikt();
    >|    }
    >|}
\end{verbatim}

Um einen Codeblock\index{Codeblock} auszugeben, der einen vorgegebenen Programmcode erzeugt, muss man lediglich vor jeder Zeile des Codes das Ausgabezeichen \verb=>|= einf�gen. Bessere Editoren wie z.B.\ Notepad++, bzw.\ Visual Studio Code erm�glichen dies sehr einfach dadurch, dass man den Cursor mit der Tastenkombination \emph{Shift + Alt + Cursor down}, bzw.\ \emph{Strg + Alt + Cursor down} �ber mehrere Zeilen ausdehnen kann, so dass nachfolgend eingetippte Zeichen gleichzeitig in jeder einzelnen Zeile eingef�gt werden.

Da wir ja nicht die Anleitung f�r die schon vorhandene Kurse-Archivierungsklasse ben�tigen, sondern eine allgemeine Anleitung haben wollen, geht es im n�chsten Schritt darum, die variablen Anteile zu identifizieren und das Skript dahingehend anzupassen.

Die erste Abh�ngigkeit ist die Bezeichnung \emph{Kurs}, die hier gleich in verschiedenen Bedeutungen auftritt, zum einen als lesbare Bezeichnung des Objekttyps zum anderen aber auch als Name einer Klasse, die an den Generik-Parameter �bergeben wird, und zuletzt auch noch als Bestandteil des Klassennamens. Letzteres ist dahingehend zu beachten, da als Objekttypbezeichnung eventuell auch Bezeichnungen mit Umlauten oder Bindestrichen auftreten k�nnen, die in einem Klassennamen nichts zu suchen haben. Ein weiterer Unterschied besteht darin, dass die Bezeichnung der Archivierungsklasse hier festgelegt werden kann, wohin gegen die Kurs-Klasse schon im System vorhanden ist und hier nur referenziert wird.

Wir fragen also die Bezeichnung des Objekttyps, sowie den Namen der dazugeh�renden Klasse ab und k�nnen daraus die �brigen Verwendungsvarianten ableiten.

F�r die Umwandlung eines Textes in eine Zeichenkette, die weder Umlaute, noch Sonderzeichen enth�lt gibt es in \FPZ\ den Befehl \fcmd{CamelCase}, der  die dazu notwendigen Ersetzungen durchf�hrt. Die Zeichenkette bleibt insgesamt lesbar und kann z.B.\ als Klassenname, Eigenschaft oder Variable im Programmcode verwendet werden.

Um es ganz sch�n zu haben, leiten wir auch die Variable \verb|k| aus dem Objekttyp ab und ersetzen diese an den verschiedenen Stellen. Dazu suchen wir mit einem regul�ren Ausdruck den ersten Gro�buchstaben in der Bezeichnung und wandeln diesen mit dem Befehl \fcmd{ToLower} in einen Kleinbuchstaben um. Hier muss man nat�rlich darauf achten, dass die Variable nicht in Konflikt mit anderen Variablen kommen kann. Die Umwandlung einer Zeichenkette in Gro�buchstaben w�re analog mit dem Befehl \fcmd{ToUpper} m�glich.

Die in der Klasse hart codierte Guid-Kennung kann frei gew�hlt werden, und wird vom Skript selbst als Zufallswert vom Typ \emph{Guid} erzeugt. Daf�r wird der \cidxvar{NewGuid}-Befehl verwendet.

Der Programmcode in der Routine f�r die Konfliktpr�fung ist sehr spezifisch auf die Kurs-klasse bezogen und l�sst sich nicht verallgemeinern. Wir ersetzen diesen durch eine ToDo-Aufforderung\index{ToDo-Aufforderung} im Code und erg�nzen die Anleitung um einen Punkt. Wir gehen sogar zus�tzlich davon aus, dass es nicht in jedem Fall Konflikte geben kann und fragen diese Voraussetzung explizit ab. F�r den Fall, das ist keine Konflikte geben kann, lassen wir die entsprechenden Programmzeilen und Anleitungsschritte weg. 

Die beiden Stufen der Aufz�hlung lassen sich bei Anleitungen sehr gut so verwenden, dass man in der ersten Ebene die durchzuf�hrenden Punkte benennt und die zweite Ebene dazu verwendet, diese genauer auszuf�hren und zus�tzliche Informationen bereitzustellen.

Am Ende bekommen wir die folgende allgemeine Anleitung:

\Beispiel{Archivklasse 2}\label{Archivklasse2}
\begin{verbatim}
~Input Bez: Bezeichnung Objekttyp (Singular)
    ~AddHelpLine Die Bezeichnung des Objekttyps,
    __ f�r den Objekte archiviert werden sollen.
~Input Klasse: Datenklasse f�r diesen Objekttyp
    ~AddHelpLine Der Name der Datenklasse, mit der Objekte
    __ dieses Typs im Programm abgebildet werden.
?QKonf: Kann es Konflikte geben?
    ~AddHelpLine Gemeint sind Konflikte, die einer Archivierung
    __ im Wege stehen w�rden.
    #j: Ja
        ~Set KonflikteMoeglich=ja
    #n: Nein
~CamelCase cBez=$Bez
~Set AKennung=$NewGuid
~RegExMatch vmatch = $BezX|([A-Z])
~ToLower v=$vmatch(1)
@Anleitung >> Ordner Archivierer �ffnen.
@Anleitung >> Neue cs-Datei erstellen mit Bezeichnung
    >|$cBezArchivierer.cs
@Anleitung >> Folgenden Code einf�gen:
    >|// Klasse f�r die Archivierung von $Bez-Objekten
    >|public class $cBezArchivierer : BasisArchivierer<$Klasse>
    >|{
    >|    // Kennung der Archivierer-Klasse
    >|    public override Guid ArchiviererKennung()
    >|    {
    >|        return Guid("$AKennung");
    >|    }
    >|
    >|    // Methode zum Archivieren eines $Bez-Objektes
    >|    public override void Archiviere($Klasse $v)
    >|    {
    ~If $KonflikteMoeglich==ja
        >|      var konflikt = PruefeAufKonflikte($v);
        >|        if (konflikt.vorhanden)
        >|        {
        >|            ALogger.Note($"Konflikt f�r $Bez '{$v}', 
        >|                {konflikt.Details}");
        >|            return;
        >|        }
    >|        base.StarteArchivierung($v);
    >|    }
    ~If $KonflikteMoeglich==ja
        >|    
        >|    // Konfliktpr�fung f�r die $Bez-Archivierung
        >|    private Konflikt PruefeAufKonflikte($Klasse $v)
        >|    {
        >|        //ToDo: Konfliktpr�fung implementieren
        >|        
        >|        // Ansonsten
        >|        return Konflikt.KeinKonflikt();
        >|    }
    >|}
~If $KonflikteMoeglich==ja
    @Anleitung >> Konfliktpr�fung implementieren
        > in $cBezArchivierer.cs in
        ~AddCode  PruefeAufKonflikte(...)
        > bei
        ~AddCode  //ToDo: Konfliktpr�fung implementieren
        > Auflistung der Konflikte nach folgendem Muster:
        >|        if (Konfliktbedingung)
        >|            return new Konflikt("Konfliktbeschreibung");
\end{verbatim}

Beim Aufbau eines Anleitungsskriptes f�r eine Standardentwicklung stechen die Schw�chen der eigenen Framework-Klassen besonders hervor, etwa zwingend erforderliche Methoden, die jedoch nicht erzwungen werden, unsch�ne Abh�ngigkeiten von Aufrufreihenfolgen oder generell ungewollte Codewiederholungen. Eine ideal aufgebaute Klasse vermeidet Wiederholungen und f�hrt den Entwickler allein mit Hilfe des Compilers an allen relevanten Fragestellungen vorbei und sorgt so selbst�ndig f�r einen vollst�ndigen und musterg�ltigen Einbau. Da dieser Anspruch schwierig aufrecht zu halten, und noch schwieriger im Nachhinein herzustellen ist, und da die Interaktion mit einem gut gemachten Anleitungsskript gef�lliger sein d�rfte, als die mit dem C\#-Compiler, bleiben trotzdem genug Gr�nde, um solche Anleitungen zu erstellen.

\subsection{Testvorbereitung}
Jede Entwicklung sollte getestet werden, im Idealfall sogar durch automatisierte Tests. Standardentwicklungen sind davon nicht ausgeschlossen auch wenn sie dadurch, dass sie auf schon mehrfach genutztem Framework aufbauen, ein geringeres Fehlerrisiko aufweisen, wie neue Individualentwicklungen.

�hnlich wie die Implementierung, variieren auch die Testpunkte\index{Testpunkte} nur in Abh�ngigkeit der Bezeichnungen und Optionen innerhalb der Entwicklung und k�nnen so ebenfalls durch das Anleitungsskript erstellt werden. Unabh�ngig davon, ob die Entwickler ihre Entwicklungen gegenseitig testen, oder ob daf�r eine eigene Testabteilung bereitsteht, ist letztlich der Ersteller des Programmcodes daf�r verantwortlich, den Testumfang so zu formulieren, dass jede neu entwickelte oder ge�nderte Zeile zumindest einmal durchlaufen wird. Hierf�r muss er die Perspektive des Anwenders annehmen, und die dazugeh�renden Schritte der Programmbedienung so formulieren, dass im Hintergrund die Zust�nde und F�lle auftreten, bei denen die implementierte Funktionalit�t durchlaufen wird.

In gleicher Weise muss er auch den Effekt der Funktionen so beschreiben, dass dieser an der Programmoberfl�che �berpr�ft werden kann. Das kann in manchen F�llen herausfordernd sein, insbesondere wenn es darum geht, bewusst Konflikte und Ausnahmezust�nde zu provozieren, die im normalen Ablauf eher selten auftreten.

Im Anleitungsskript f�r den Entwickler k�nnen wir auf jeden Fall schon einmal die Integrationstests\index{Integrationstests} ausformulieren, die sich unmittelbar aus dem Programmcode aus der Anleitung ergeben. F�r den Fall, dass Konflikte m�glich sind, erg�nzen wir im ersten Testpunkt den Zusatz \emph{sofern keine Konflikte auftreten}. Zus�tzlich erg�nzen wir die Anleitung dahingehend, eine Testaufgabe zu erstellen und m�gliche Konfliktarten in den Testpunkten aufzuz�hlen. Die Testpunkte werden in einem eigenen Abschnitt aufgelistet und k�nnen so parallel zu den dazugeh�renden Anleitungsschritten gesammelt werden.

Damit kommen folgende Zeilen in unserer Anleitung hinzu:

\Beispiel{Archivklasse 3}\label{BspArchivklasse3}
\begin{verbatim}
...
@Testpunkte >> Ein $Bez-Objekt kann archiviert werden
~If $KonflikteMoeglich==ja
    ~AddText , sofern keine Konflikte auftreten
~AddText .
@Testpunkte >> Die Archivierungskennung f�r $Bez-Objekte 
    __ lautet $AKennung.
    > Kann im Protokoll nachgeschlagen werden.
...
@Anleitung >> Testaugabe erstellen
    > Titel: Test der $Bez-Archivierung
    > Testpunkte aus Abschnitt "Testpunke" �bernehmen
~If $KonflikteMoeglich==ja
    @Anleitung >> Konflikte in Testpunkten auflisten
        > Bei "ToDo: Konflikte auflisten"
    @Testpunkte >> Folgende Konflikte werden erkannt:
        > ToDo: Konflikte auflisten
    @Testpunkte >> Wenn ein Konflikt erkannt wird,
        ># wird die Archivierung abgebrochen,
        ># protokolliert, wo der Konflikt aufgetreten ist,
        ># Details zum Konflikt protokolliert.
\end{verbatim}

Was mit Integrationstests in Textform m�glich ist, ist je nach Art der Entwicklung auch noch in Form von Unittests\index{Unittests} m�glich.�Unitest sind im Programmcode hinterlegte Klassen und Funktionen, die nur daf�r da sind, die im Produktivcode verwendeten elementaren Klassen auf Korrektheit zu testen. Die Tests selbst geh�ren nicht zum Produkt und haben keine Funktionen in der Produktivumgebung. Unittests werden in der Regel durch spezielle Attribute gekennzeichnet und in eigenen Projekten verwaltet.

Im Fall unserer Archivierungsklasse k�nnte z.B. die Archivierer-Kennung sehr einfach �ber einen Unittest gepr�ft werden. Das Skript k�nnte dazu wie folgt den entsprechenden Testcode erzeugen:

\Beispiel{Archivklasse 4}\label{BspArchivklasse4}
\begin{verbatim}
...
@Anleitung >> Ordner ArchiviererTest �ffnen.
@Anleitung >> Neue cs-Datei erstellen mit Bezeichnung
    >|$cBezArchiviererTests.cs
@Anleitung >> Folgenden Code einf�gen:
    >|public class $cBezArchiviererTests
    >|{
    >|    [Fact]
    >|    public void ArchiviererKennung_ShouldReturnExpectedGuid()
    >|    {
    >|        // Arrange
    >|        var archivierer = new $cBezArchivierer();
    >|
    >|        // Act
    >|        var result = archivierer.ArchiviererKennung();
    >|
    >|        // Assert
    >|        Assert.Equal(Guid.Parse("$AKennung"), result);
    >|    }
    >|}
\end{verbatim}

Denkbar w�re auch ein Skript, dass direkt von den Testern verwendet wird, um sich die eigene Arbeit zu vereinfachen. Dieses m�sste dann auf Basis der Informationen ausgef�llt werden, die von den Entwicklern bereitgestellt werden, und es m�sste Mehrwert generieren, indem es daraus hilfreiche Informationen f�r die Durchf�hrung von Tests ableitet.

Es ist schwierig, in diesem Bereich Beispiele zu finden, die nicht besser in der oben genannten Form als Teil der Testvorbereitung in der Entwicklungsanleitung aufgehoben sind. Ein denkbares Anwendungsgebiet w�ren auf jeden Fall die automatisierten Oberfl�chentests. Hierbei wird mit speziellen Tools eine Programmbedienung �ber die Benutzeroberfl�che simuliert, d.h.\ es werden Schaltfl�chen gedr�ckt und Eingaben get�tigt und dabei wird immer wieder die Reaktion des Programms mit einem vorgegebenen Sollverhalten abgeglichen.

Solche Oberfl�chentests\index{Oberfl�chentests} werden in der Regel ebenfalls als Programmcode verwaltet und folgen oft �hnlichen Strukturen, wie der Produktivcode, also mit wiederverwendbaren Anteilen und Mustern und einem eigenen Framework. Insbesondere das Vorbereiten einer Testsituation und das Aufr�umen nach Durchf�hrung des Tests sind meist aufwendige und oft ben�tigte Funktionen, die einen gut organisierten Baukasten voraussetzen.

Ein Unterst�tzungsskript k�nnte sehr gut die Voraussetzungen f�r einen geplanten automatisierten Oberfl�chentest abfragen und damit schon mal den Test soweit wie m�glich vorbereiten.

\subsection{Product-Owner-Unterst�tzung}\label{POUnterstuetzung}
Der Product Owner\index{Product Owner} oder kurz PO ist die zentrale Figur im Entwicklungsprozess. Er kennt das Produkt wie seine Westentasche, wei�, wie es sich entwickelt hat und wohin es sich entwickeln soll, und mit welchen Bausteinen und nach welchen Prinzipien diese Entwicklung abl�uft. Er steht im direkten Austausch mit den Kunden, kennt die Anwender und hat gelernt, dass auch die schlimmsten Entscheidungen aus guten Absichten getroffen werden.

Die Menge an Dingen, die bei der professionellen Softwareentwicklung zu beachten sind, sind ungeheuer vielschichtig, und die permanente Weiterentwicklung eines gro�en bestehenden Systems bringt noch zahlreiche eigene Aspekte dazu. Auch wenn der PO bei dieser Aufgabe nicht allein steht, und sowohl auf das ganze Entwicklungsteam, als auch auf die ihm zuarbeitenden Spezialisten f�r Usability und Oberfl�chendesign bauen kann, h�ngt es nachher doch an ihm, aus dem ganzen Wissen und der Erfahrung ein gro�es Ganzes zu machen, das individuelle Vorlieben und W�nsche zur�ckstellt und auf f�r alle nachvollziehbaren Regeln aufbaut.

Unterst�tzungssysteme f�r die Product Owner waren mit die ersten Anwendungsbeispiele, die mit \FPZ\ umgesetzt wurden. Gerade dort, wo Standardentwicklungen haupts�chlich in einem sehr engen Rahmen konfiguriert werden mussten, konnten entsprechende Skripte all die ben�tigten Punkte abfragen und daraus eine weitestgehend vollst�ndige Beschreibung erstellen, die alle entscheidenden Aspekte der Entwicklung in hohem Detailgrad umfasste.

Wir bauen nun ein Beispiel f�r ein solches System auf und beginnen mit einem Teilaspekt, der in vielen PO-Unterst�tzungssystemen ben�tigt wird, n�mlich der Beschreibung der Stelle innerhalb der Anwendung, an der die Entwicklung greifen soll. Wie schon in Abschnitt~\ref{Strukturabbildung} beschrieben, l�sst sich die Struktur einer Anwendung gut als Entscheidungsbaum aufbauen, so dass sich ein beliebiger Programmbereich oder ein Registerblatt mit wenigen Klicks ausw�hlen l�sst.

Wir speichern dieses Skript f�r die Auswahl in einer eigenen Funktionsdatei\index{Funktionsdatei} mit der Dateiendung\index{Dateiendung} \emph{.fps} und k�nnen diese dann in verschiedenen anderen Skripten mit Hilfe des \fcmd{Include}-Befehls aufrufen. Durch die andere Dateiendung wird das Skript nicht in der Auswahl von \FPZ\ angezeigt, so dass es nicht eigenst�ndig aufgerufen werden kann. Der Befehl verf�gt �hnlich wie der \fcmd{GoSub}-Befehl �ber einen \cidxfrag{BaseKey}-Parameter, so dass man in der Datei die \verb|$BaseKey|-Variable nutzen kann und mehrere Aufrufe innerhalb desselben Skriptes m�glich sind, ohne dass es zu Konflikten bei den Schl�sseln f�r Auswahlabfragen und Eingaben kommt.

Da die Auswahl eines Programmbereichs f�r verschiedene Zwecke verwendet werden kann, schaffen wir �ber die Variable \verb|PHinweis| die M�glichkeit, einen Text zu �bergeben, der die Fragestellung um den Verwendungszweck erweitert. Die getroffene Auswahl wird dann nach Ausf�hrung �ber die beiden Variablen \verb|PAusw| und \verb|UAusw| an das aufrufende Skript zur�ck �bergeben.

\Beispiel{Programmauswahl.fps}\label{BspProgrammauswahl}
\begin{verbatim}
?P: Programmbereich $PHinweis
    #D: Dozenten
        ~Set PAusw=Dozenten
        ?U: Unterbereich in Dozenten $PHinweis
            #A: Auswahliste Dozenten
                ~Set UAusw=Auswahliste Dozenten
            #S: Registerblatt Stammdaten
                ~Set UAusw=Registerblatt Stammdaten
            #F: Registerblatt F�cher
                ~Set UAusw=Registerblatt F�cher
    #K: Klassen
        ~Set PAusw=Klassen
        ?U: Unterbereich in Klassen $PHinweis
            #A: Auswahliste Klassen
                ~Set UAusw=Auswahliste Klassen
            #S: Registerblatt Sch�ler
                ~Set UAusw=Registerblatt Sch�ler
            #U: Registerblatt Unterricht
                ~Set UAusw=Registerblatt Unterricht
\end{verbatim}

Unsere Anwendungsbeispiel selbst unterst�tzt bei der Beschreibung einer Listenfunktion, womit in diesem Zusammenhang eine Programmfunktion gemeint ist, die f�r ausgew�hlte Zeilen einer Liste aufgerufen werden kann.

Diese Aufgabenstellung klingt zun�chst recht einfach, aber wie wir sehen werden, gibt es viel zu beachten, an das man nicht unmittelbar denkt. Genau hier liegt der Vorteil eines solchen Skiptes, das diese Punkte systematisch abfragt.

\Beispiel{Listenfunktion}\label{BspListenfunktion}
\begin{verbatim}
~Input Bez: Bezeichnung der Funktion
~Set PHinweis= (Wo soll die Funktion eingebaut werden?)
~Include Programmauswahl.fps; BaseKey=E

@Story >>_ Einbau der Listenfunktion "$Bez"
    __ in die Liste des Bereichs "$PAusw - $UAusw".

@Story >> Die Funktion ist per Kontextmen� aufrufbar.
    > Bezeichnung Men�punkt: $Bez
?Pos1: Gibt es schon einen Funktionen-Abschnitt im
    __ Kontextmen�?
    #j: Ja
        ?Pos2: An welcher Stelle soll der Men�punkt in
        __ das Kontextmen� eingebaut werden?
            #O: Ganz oben
                ~Set wogenau=ganz oben
            #U: Ganz unten
                ~Set wogenau=ganz oben
            #D: Irgendo dazwischen
                ~Input Pos3: Einbau unterhalb von...
                ~Set wogenau=unterhalb von "$Pos3"
        > ...in vorhandenen Funktionen-Abschnitt $wogenau
    #n: Nein
        > Funktionen-Abschnitt anlegen nach Std.-funktionen

@Story >> Die Funktion
?QRechte: Welche Rechte setzt die Funktion voraus?
    ~AddHelpLine Schreibrechte bei schreibenden Funktionen
    #L: Leserechte auf die Liste
        ~AddText  setzt Leserechten auf die Liste voraus.
    #S: Schreibrechte auf die Liste
        ~AddText  setzt Schreibrechte auf die Liste voraus.
    #E: Die Funktion hat eigene Rechte
        ~AddText  hat eigene Rechte.
        > Die Berechtigungsstruktur muss erweitert werden.
        > Eingliederung unter Listenfunktionen - $Bez
        > Das neue Recht ist zun�chst nur f�r Admins gesetzt.

@Story >> Die Funktion ist
?AKont: Die Funktion ist
    #E: ...beschr�nkt auf Einfachauswahl
        ~AddText  beschr�nkt auf Einfachauswahl.
        > Bei Mehrfachauswahl Men�punkt deaktivieren
    #M: ...auch verf�gbar f�r Mehrfachauswahl
        ~AddText  auch verf�gbar f�r Mehrfachauswahl.

@Story >> Die Funktion
?Wirk: Die Funktion
    #A: ...�ffnet einen Assistenten
        ~AddText  �ffnet einen Assistenten.
        > Beschreibung und Skizze Assistent: xxx
        @ToDo >> Assisten skizzieren und beschreiben
    #R: ...wird nach Dialog-R�ckfrage ausgef�hrt
        ~AddText  wird nach Dialog-R�ckfrage ausgef�hrt.
        > Typ: R�ckfrage-Dialog vor Ausf�hrung
        > Text: xxx
        @ToDo >> Text Dialog-R�ckfrage ausformulieren
    #O: ...wird ohne Dialog-R�ckfrage ausgef�hrt
        ~AddText  wird ohne Dialog-R�ckfrage ausgef�hrt.

@Story >> Die Funktion bewirkt folgendes:
        > xxx
@ToDo >> Wirkung der Funktion beschreiben.

@Story >> Das Benutzerfeedback
?Feed: Das Benutzerfeedback
    #E: ...ist unmittelbar erkennbar
        ~AddText  ist unmittelbar erkennbar.
        > Kein Feedback-Dialog notwendig.
    #D: ...wird explizit �ber einen Dialog gegeben
        ~AddText  wird explizit �ber einen Dialog gegeben.
        > Typ: Feedback-Dialog nach Ausf�hrung
        > Text: xxx
        @ToDo >> Text Feedback-Dialog ausformulieren
\end{verbatim}

Zun�chst muss die Funktion benannt werden und dann muss die Liste beschrieben werden, f�r die die Funktion eingebaut werden soll. Dies ist mit einer Texteingabe und dem Aufruf der als Funktionsdatei ausgelagerten Auswahlfunktion schnell umgesetzt.

Die Festlegung, dass eine Listenfunktion �ber das Kontextmen� der Liste aufgerufen wird, ist fest im Skript hinterlegt. Der Product Owner hat hier keine Auswahlm�glichkeit und soll diese auch nicht haben. Er bestimmt, was er haben m�chte und das Designsystem\index{Designsystem} legt Form und Aussehen fest.

Dadurch, dass die grundlegende Design-Entscheidung fest im Skript hinterlegt ist, kann diese wiederum explizit in die Beschreibung einflie�en, sodass diese durchg�ngig ber�cksichtigt wird und niemand das Designsystem auswendig im Kopf haben muss.

Das Designsystem legt in unserem Fall ebenfalls fest, dass innerhalb des Kontextmen�s ein Abschnitt f�r solche Funktionen angelegt werden muss, der beim Einbau einer neuen Funktion auch schon vorhanden sein kann. Wenn der Abschnitt schon vorhanden ist, wird der Product Owner aufgefordert, die Position f�r den neuen Men�punkt explizit vorzugeben. Auch derartig kleine Details sollten nicht dem Zufall �berlassen werden.

Ein Funktionsaufruf erfordert Berechtigungen, die sich entweder aus den Berechtigungen f�r die Liste ableiten lassen oder explizit angelegt werden m�ssen. Wenn die Funktion Daten ver�ndert, sollte sie in jedem Fall Schreibrechte auf die Liste voraussetzen, wenn die Daten nur lesend verwendet werden, reichen auch Leserechte. Solche Hinweise auf allgemein festgelegte Regeln sollten im Skript als Hilfezeilen bei den damit verbundenen Fragen und Texteingaben erg�nzt werden.

Im n�chsten Schritt wird gefragt, ob die Funktion auf einzelne Zeilen beschr�nkt ist, oder auch f�r eine Mehrfachauswahl von Zeilen aufgerufen werden kann. Es ist wichtig, sich �ber solche Punkte im Vorfeld Gedanken zu machen, da oft beides m�glich, der allgemeine Fall jedoch meist deutlich schwieriger umzusetzen ist und manchmal auch noch weitere Fragen nach sich zieht. Oft wird diese Entscheidung erst nach R�cksprache mit den Entwicklern getroffen, wenn der Mehraufwand abgesch�tzt ist. Die R�cksprache mit den Entwicklern k�nnte bei solchen Punkten sogar als eigene Antwortoption zur Auswahl gestellt werden.

Nun wird festgelegt, wie die Funktion bei Auswahl reagiert. Soll ein Assistent ge�ffnet werden oder soll die Funktion direkt ausgef�hrt werden? Soll die Ausf�hrung der Funktion �ber einen R�ckfrage-Dialog abgesichert werden, in Sinne von \emph{sind Sie sicher?} Auch hier bezeichnet R�ckfrage-Dialog ein Element aus dem Designsystem, dessen Aussehen weitestgehend festgelegt ist. Symbol und die zur Auswahl stehenden Schaltfl�chen werden daher nicht abgefragt, nur der Text.

In gleicher Weise wird am Ende sichergestellt dass der Benutzer auch ein Feedback f�r die Ausf�hrung der Funktion bekommt. Wenn dieses unmittelbar �ber die Programmoberfl�che erkennbar ist, muss daf�r nichts zus�tzliches getan werden. Ansonsten sollte eine explizite R�ckmeldung in Form eines Dialogs gegeben werden.

Die eigentliche Beschreibung, was die Funktion genau macht und wie der Assistent aussieht, der eventuell aufgerufen wird, ist ebenfalls Teil der Beschreibung. Wie alles, was nur einmalig in Textform beschrieben werden muss, werden diese Informationen nicht �ber Texteingaben abgefragt, sondern am Ende durch Platzhalter in der Beschreibung gekennzeichnet, und als noch ausstehende Punkte in einer ToDo-Liste\index{ToDo-Liste} notiert.

Nachdem alles ausgef�llt ist, kann der Product Owner die Beschreibung in einen Vorgang �bertragen und die noch ausstehenden Punkte dort mit Leben f�llen, wobei er sich an dieser ToDo-Liste orientieren kann. Durch das Skript wurde sichergestellt, dass alle von ihm zu treffenden Entscheidungen ber�cksichtigt wurden und explizit als Teil des Vorgangs aufgelistet sind. 

Das explizite Auflisten solcher Details ist ein wichtiger Punkt zur Sicherstellung der Qualit�t. Die Verschriftlichung hilft zuallererst dem Produkt Owner selbst, der dadurch nochmal eine bessere Vorstellung von der von ihm beauftragten Entwicklung bekommt. In gleicher Weise hilft sie dem Entwickler der sich bei der gesamten Entwicklung klar an diesen Punkten orientieren kann, weniger R�ckfragen stellen muss und die Liste auch nach Fertigstellung zur Gegenpr�fung und zum Abhaken aller Punkte verwenden kann. Auch bei der Abnahme durch den Product Owner kann die explizite Auf\-listung der einzelnen Detailanforderungen unterst�tzend eingesetzt werden.

Diese Form von Entwicklungsunterst�tzung scheint alles andere als agil zu sein und passt nicht zu den blumig beschriebenen User Stories im Srum-Prozess und dem permanenten Austausch zwischen Entwicklung und Product Owner. Hier muss man wieder ganz klar unterscheiden zwischen der Entwicklung von Framework\index{Framework} und der Verwendung des entwickelten Frameworks. Die Entwicklung von neuem Framework kann und sollte agil verlaufen, da Anforderungen und technische M�glichkeiten permanent aufeinander abgestimmt werden m�ssen. Hier ist viel Aufwand und viel Austausch notwendig, damit das neu Geschaffene am Ende zum Bestehenden passt und systematisch, ohne gro�en Abstimmungsbedarf verwendet werden kann. Die Erweiterung des Designsystems, die Aufstellung von Regeln f�r die Verwendung und idealerweise auch die Erstellung von Unterst�tzungskripten sollten fester Bestandteil bei der Entwicklung neuer Framework-Komponenten sein.

\subsection{Qualit�tprotokolle}
Qualit�tsprotokolle sind �hnlich wie die Unterst�tzungskripte Hilfsmittel, die den am Entwicklungsprozess beteiligten Mitarbeitern helfen, ihre Arbeit einfacher und nach einer definierten Vorgehensweise durchzuf�hren. Letzteres steht hierbei im Vordergrund, denn das Protokoll macht Vorgaben, die zwingend einzuhalten sind, und die auch Gegenstand von internen Richtlinien sein k�nnen.

Insbesondere dann, wenn man als Unternehmen eine Zertifizierung\index{Zertifizierung} nach den Normen ISO 9001\index{ISO 9001} oder ISO 25001\index{ISO 25001} anstrebt und die damit verbundenen Audits bestehen m�chte, reicht es nicht, Qualit�t, Datenschutz und Datensicherheit einfach nur zu wollen, man muss Ma�nahmen im Management zur permanenten Aufrechterhaltung diese Aspekte definieren, dokumentieren, einhalten und auch nachweisen k�nnen. Dies gelingt langfristig nur dadurch, dass man diese Dinge tats�chlich in den Arbeitsalltag integriert und lebt, und den Nutzen zu sch�tzen lernt, den diese f�r den Kunden und auch das Unternehmen selbst haben. 

Je mehr man diese Integration auf den tats�chlich beabsichtigten Effekt ausrichtet, anstatt nur nach einer formalen Erf�llung der verlangten Punkte zu streben, desto einfacher f�llt es am Ende, diese Dinge in der Routine zu ber�cksichtigen und einen wirklichen Mehrwert daraus zu generieren.

Qualit�tsprotokolle\index{Qualit�tsprotokolle} sind ein einfaches Mittel daf�r und lassen sich an vielen Stellen einrichten, am besten dort, wo regelm��ige T�tigkeiten auftreten, bei denen es wichtige Dinge zu beachten gibt.

Als Fortsetzung von Abschnitt~\ref{POUnterstuetzung} beginnen wir mit einem Protokoll f�r den Product Owner, das dabei hilft, bestimmte Wirkungen der von ihm beauftragten Entwicklungen zu beleuchten, und sich damit auseinanderzusetzen.

Mit einem Fehler im Kontext der Softwareentwicklung assoziieren wir gemeinhin eine fehlerhafte Abweichung vom Sollzustand, die bei der Entwicklung entstanden ist, aber es gibt auch Fehler die schon vorher entstehen und die als Teil der Anforderungen den gesamten Entwicklungsprozess bis hin zum Test durchlaufen k�nnen, ohne dass jemand merkt, dass es falsch ist. Oft geschieht so etwas durch einen unbedachten Umgang mit personenbezogenen Daten und f�hrt am Ende zu einem meist kleinen, aber dennoch ungewollten Datenschutzversto� der h�tte vermieden werden k�nnen.

Ein klassisches Beispiel daf�r ist der Wunsch, das Geburtsdatum als Spalte f�r einen Auswahldialog f�r Personen hinzuzunehmen. Der Grund daf�r ist gut nachvollziehbar, und mit Sicherheit erleichtert diese zus�tzliche Information die Auswahl an verschiedenen Stellen, und hilft besonders beim Vorhandensein von gleichnamigen Personen. Andererseits wird die personenbezogene Information des Geburtsdatums an Stellen im Programm abrufbar, wo es prim�r nicht um die Stammdaten der Personen geht, und damit erweitert sich der Zugriff auf diese Information eventuell auf Benutzer, die diesen zuvor nicht hatten.

Die Gefahr, solche Dinge zu �bersehen, ist recht gro�, da sie oft nebenl�ufig zu den eigentlich ein Anforderungen stehen. Insofern ist es gut, ein Protokoll zu haben, dass solche Aspekte einzeln und jedes mal abfragt.

Der nachfolgende PO-Assistent\index{PO-Assistent} enth�lt beispielhaft drei Kontrollfragen aus den Bereichen Datenschutz und Datensicherheit, die der Product Owner nach Beendigung seiner Story-Formulierung durchgehen sollte.

\Beispiel{PO-Assistent}\label{BspPOAssi}
\begin{verbatim}
?Q1: Werden durch die Entwicklung Kundendaten
    __ au�erhalb der Programmumgebung gespeichert?
    #j: Ja
        @Wirkungen >> Speicherung von Kundendaten 
        __ au�erhalb der Programmumgebung: Ja
        @Erg�nzungsma�nahmen >> Speicherung von Kundendaten 
        __ au�erhalb der Programmumgebung:
        > Sicherstellen, dass der Speicherort so gew�hlt 
        __ werden kann, dass kein unautorisierter Zugriff 
        __ m�glich ist
        > Sicherstellen, dass Speichervorgang und Speicherort
        __ f�r den Anwender gut erkennbar sind
    #n: Nein
        @Wirkungen >> Speicherung von Kundendaten 
        __ au�erhalb der Programmumgebung: Nein
?Q2: Werden durch die Entwicklung personenbezogene Daten an 
    __ Stellen im Programm sichtbar wo sie es bisher nicht waren? 
    #j: Ja
        @Wirkungen >>  Erweiterte Sichtbarkeit 
        __ personenbezogener Daten: Ja
        @Erg�nzungsma�nahmen >> Erweiterte Sichtbarkeit 
        __ personenbezogener Daten:
        > Sicherstellen, dass der Zugriff auf diese Stellen durch
        __ die Vergabe von Berechtigungen abgesichert ist
        > Sicherstellen dass nicht mehr Daten angezeigt werden
        __ als in diesem Bearbeitungskontext notwendig ist
        > Ggf. Datenschutzbeauftragten mit einbeziehen
    #n: Nein
        @Wirkungen >>  Erweiterte Sichtbarkeit 
        __ personenbezogener Daten: Nein
?Q3: Werden durch die Entwicklung neue benutzerbezogene
    __ Daten protokolliert oder verarbeitet? 
    #j: Ja
        @Wirkungen >> Protokollierung von Nutzerdaten: Ja
        @Erg�nzungsma�nahmen >> Protokollierung von Nutzerdaten:
        > Sicherstellen dass die Protokollierung zweckgebunden
        __ und durch die Benutzervereinbarung abgedeckt ist
        > Sicherstellen, dass die Protokollierung in das 
        __ Datenblatt f�r Protokollierung aufgenommen wird
    #n: Nein
        @Wirkungen >> Protokollierung von Nutzerdaten: Nein
\end{verbatim}

Die Ausgabe besteht aus zwei Abschnitten. Der Abschnitt \emph{Erg�nzungsma�nahmen} beschreibt Ma�nahmen, die im Zusammenhang mit der Story sichergestellt werden sollten, z.B.\ indem entsprechende Anforderungen erg�nzt werden. Der Abschnitt \emph{Wirkungen} dient als Protokoll, dass in jedem Fall in der Story hinterlegt wird. Damit wird nachweisbar sichergestellt, dass der Assistent f�r alle eingeplanten Stories durchlaufen wurde, und man kann im Falle eines Fehlers besser nachvollziehen, wo dieser entstanden ist, und wie man diesem zuk�nftig vorbeugen kann.

Wie schon erw�hnt, kann man solche Protokolle an vielen Stellen verankern, und nat�rlich gibt es auch auf technischer Ebene sehr viele Dinge, die vom Entwickler ber�cksichtigt werden m�ssen.

Viele Teams praktizieren Verfahren wie das Vier-Augen-Prinzip\index{Vier-Augen-Prinzip} oder eine Form von Peer-Review\index{Peer-Review}, bei der abgeschlossene Entwicklungen einem Kollegen zur Pr�fung vorgelegt werden. Die versch�rfte Varianten davon nennt sich \emph{Smash It}\index{Smash It}. Hier versucht der angefragte Kollege gezielt, und unter Nutzung seines Systemwissens, Fehler und Schwachstellen der Entwicklung aufzudecken, damit diese vor Herausgabe der Version behoben werden k�nnen. Dies sind gute Erg�nzungen zu den Integrationstests, die ja prim�r auf das Abpr�fen der neuen Funktionalit�t ausgerichtet sind und nicht auf das Auffinden von Schwachstellen in deren Umgebung. Trotzdem h�ngt es am Ende hauts�chlich von der Erfahrung und der investierten Zeit des Einzelnen ab, ob ein kritischer Fehler oder eine Verwundbarkeit entdeckt wird oder nicht.

Im Laufe der Zeit und vor allem durch die systematische Untersuchung der leider doch nicht rechtzeitig erkannten schwerwiegenden Fehler stellt man fest, dass jedes Entwicklungsmuster seine eigenen Risikofaktoren\index{Risikofaktoren} mit sich bringt, und es in den meisten F�llen systematische Vorgehensweisen gibt, um die jeweiligen Risiken zu verringern.

Das Risiko, dass Daten unbeabsichtigt gel�scht werden, wird dann auftreten, wenn �berhaupt Daten gel�scht werden, Kompatibilit�tsprobleme bekommt man dann, wenn man mit versionsabh�ngigen Formaten hantiert. Probleme mit Drittkomponenten zeigen sich oft erst Monate oder Jahre sp�ter, entstehen aber in dem Moment wo man sich nicht gen�gend Zeit nimmt, um den Einbau einer neuen Drittkomponente sorgsam auszuf�hren.

Das nachfolgende Quali-Assistent\index{Quali-Assistent}-Protokoll ist �hnlich aufgebaut wie Beispiel~\ref{BspPOAssi} und zeigt anhand der oben genannten Punkte, wie ein Entwickler seine Entwicklung durch zus�tzliche Ma�nahmen absichern kann, wenn sie bestimmte Aspekte erf�llt.

\Beispiel{Quali-Assistent}
\begin{verbatim}
?Q1: Umfasst die Entwicklung Funktionen, die Datens�tze l�schen?
    #j: Ja
        @Wirkungen >> L�schen von Datens�tzen: Ja
        @Erg�nzungsma�nahmen >> L�schen von Datens�tzen:
        > Sicherstellen, dass keine Datens�tze ungewollt 
        __ gel�scht werden.
        > Abgrenzungstests formulieren
        > L�schweitergaben pr�fen
    #n: Nein
        @Wirkungen >> L�schen von Datens�tzen: Nein
?Q2: Umfasst die Entwicklung die Verwendung
    __ neuer Drittkomponenten?
    #j: Ja
        @Wirkungen >> Verwendung von Drittkomponenten: Ja
        @Erg�nzungsma�nahmen >> Verwendung von Drittkomponenten:
        > Checkliste f�r Drittkomponenten durchgehen
        > Drittkomponente in eigene Klasse kapseln
    #n: Nein
        @Wirkungen >> Verwendung von Drittkomponenten: Nein
?Q3: Umfasst die Entwicklung die Verarbeitung von 
    __ versionsabh�ngigen Formaten?
    #j: Ja
        @Wirkungen >> Versionsabh�ngige Formate: Ja
        @Erg�nzungsma�nahmen >> Versionsabh�ngige Formate:
        > Sicherstellung der Versionsbestimmbarkeit
        > Sicherstellung von Abw�rtskompatibilit�t
        > Sicherstellung von Aufw�rtskompatibilit�t
        > Kennzeichnung der Versionsabh�ngigkeit 
        __ im Datenblatt.
    #n: Nein
        @Wirkungen >> Versionsabh�ngige Formate: Nein
\end{verbatim}

Bei der Implementierung einer Funktion, die Datens�tze l�scht, wird meist nur gepr�ft, ob das zu L�schende auch wirklich gel�scht wird. Viel wichtiger ist die �berpr�fung, ob das, was nicht gel�scht werden soll, auch wirklich nicht gel�scht wird. Hierf�r kann man Abgrenzungstests\index{Abgrenzungstests} formulieren und man kann man sich -- auch mit \FPZ\ -- einfache Werkzeuge bauen, die Abweichungen der Datensatz-Anzahl tabellen�bergreifend erkennbar machen. Oft wird auch �bersehen, dass eine L�schung durch festgelegte L�schweitergaben weitere L�schungen nach sich zieht. Diese sollten immer explizit gepr�ft werden.

Der Umgang mit Drittkomponenten umfasst viele Aspekte, �ber die man einen eigenen Abschnitt schreiben k�nnte. Gerade im kommerziellen Umfeld gibt es viel zu beachten, angefangen von der Vertr�glichkeit der Lizenz, bis hin zur Dokumentation und Planung notwendiger �berwachungsaufgaben f�r die Sicherstellung der Aktualit�t und regelm��ige Pr�fung auf Verwundbarkeiten. Komponenten, die f�r sicherheitsrelevante Funktionalit�t wie Verschl�sselung eingesetzt werden, m�ssen diesbez�glich besondere Anspr�che erf�llen. Im Beispiel oben wird dementsprechend auch nur auf eine Checkliste f�r Drittkomponenten verwiesen, die man nat�rlich ebenfalls mit \FPZ\ abbilden kann. 

Durch die Kapselung einer Drittkomponente in eine eigenen Klasse l�sst sich gut sicherstellen, dass nur die ben�tigte Funktionalit�t verwendet wird und die Konfiguration auf die Eigenschaften beschr�nkt bleibt, die man zulassen m�chte.

Beim Umgang mit versionsabh�ngigen Formaten sollte von Anfang an ber�cksichtigt werden, dass sich diese �ndern k�nnen, selbst wenn sie Teil des Produktes sind. Wenn ein Format um neue Eigenschaften erweitert wird, sollten vorhandene Daten in einem �lteren Format immer noch verwendbar bleiben. Insbesondere sollte ein Format immer explizit bestimmbar sein, und nicht durch das Vorhandensein oder Nicht-Vorhandensein von Eigenschaften hergeleitet werden m�ssen.

Auch hier ist das Skript wieder so aufgebaut, dass parallel zu den erg�nzenden Ma�nahmen ein Protokoll mit den Wirkungen der Entwicklung ausgegeben wird, das zum Nachweis der Auseinandersetzung in der Aufgabe hinterlegt werden k�nnte.

Die Liste solcher Aspekte lie�e sich noch um viele Punkte verl�ngern, und man muss aufpassen, dass die Relevanzkl�rung noch mit �berschaubarem Aufwand m�glich ist. Besonders sollte man darauf achten, dass die Fragestellungen nach bestimmten Aspekten grunds�tzlich unmittelbar, also ohne langes �berlegen durch den Entwickler beantwortet werden kann, zumindest wenn er sich gerade ausf�hrlich mit der Aufgabenstellung auseinandergesetzt hat.

\section{Organisationswerkzeuge}
Im Drumherum der t�glichen Arbeit gibt es viele kleine Aufgaben, die sich oft mehrfach t�glich wiederholen, und bei denen man froh ist, wenn man auf passende Hilfswerkzeuge zur�ckgreifen kann. In diesem Abschnitt beschreiben wir einige solche Werkzeuge, die sich mit wenig Aufwand als \FPZb-Skript anfertigen lassen.

\subsection{Zeitmessung}\label{Zeitmessung}
Zeitmessung\index{Zeitmessung} ist ein m�chtiges Instrument, und meist unerl�sslich, insbesondere wenn man direkt f�r einzelne Kunden arbeitet. Da Zeitmessung oftmals mit Leistungsmessung assoziiert wird, und damit auch als Argument f�r Kritik an der eigenen Leistungsf�higkeit herangezogen werden kann, tut man sich oft schwer, diese f�r verschiedene T�tigkeiten zu etablieren.

Umgekehrt finden wir durch Zeitmessung heraus, in welche Arten von Aufgaben unsere Arbeitszeit flie�t, und k�nnen auf dieser Grundlage beurteilen, ob sie dort auch in ausreichendem Ma�e zur Wertsch�pfung beitr�gt. Stellen wir etwa fest, dass viele Stunden Zeit f�r Supportunterst�tzung aufgewendet werden m�ssen, weil die Benutzerf�hrung missverst�ndlich ist, k�nnen wir gezielt gegensteuern und mehr Zeit f�r dahingehende Ma�nahmen investieren. 

Dar�ber hinaus erlaubt uns die Zeitmessung das sogenannte Controlling\index{Controlling}, also die permanente �berpr�fung ob j�ngere Anpassungen der Arbeitsweise den beabsichtigten Effekt haben.

Aufgrund der gro�en Bedeutung ist Zeitmessung in vielen Systemen fest integriert und muss nicht durch ein zus�tzliches Hilfsmittel erg�nzt werden. Das nachfolgende Beispiel zeigt eine einfache Zeitmessung �ber ein \FPZb-Skript. Dieses erm�glicht es, die Zeitmessung zu pausieren und fortzusetzen und gibt am Ende ein minutengenaues Zeitprotokoll aus.

Dadurch, dass man das Skript parallel in mehreren Browser-Tabs �ffnen kann, kann man so auch zwischen der Zeitmessung f�r verschiedenen Aufgaben hin und her wechseln.

\Beispiel{Zeitmesser}\label{BspZeitmesser}
\begin{verbatim}
~Input A: Aufgabe
~SetInputTitle Zeitmessung f�r Aufgabe $A
~Set idx=0
~GoSub NotiereZeit
~Set Status=l�uft
~Set Aktion=pausieren
~Set Summe=0
~Set fertig=Nein
~DoWhile $fertig!=Ja
    ~Set pri=$idx
    ?A($idx): Zeitmessung $Status (Summe: $Summe Minuten)
        #P: Zeitmessung $Aktion    
        #A: Zeitmessung abschlie�en    
            ~Set fertig=Ja
    ~Execute
    ~GoSub NotiereZeit    
    ~If $Status==l�uft
        ~Replace ZTerm=$ZP($idx) - ($ZP($pri))|:-> * 60 +
        ~CalculateExpression Diff = $ZTerm
        ~AddTo Summe += $Diff
        @Zeitprotokoll f�r Aufgabe $A >>|$ZP($pri)-$ZP($idx) Uhr
            __ - $Diff Minuten
        ~Set Status=pausiert
        ~Set Aktion=fortsetzten        
    ~Else
        ~Set Status=l�uft
        ~Set Aktion=pausieren
~Loop
@Zeitprotokoll f�r Aufgabe $A >>|Summe: $Summe Minuten

~DefineSub NotiereZeit
    ~AddTo idx+=1
    ~SetDateTime Zeit = HH:mm
    ~AddNewKey ZP($idx) = $Zeit
~Return
\end{verbatim}

Das Skript funktioniert wie folgt: nach Abfrage der Aufgabe, damit diese als Titel dargestellt werden kann, wird wiederholt eine Schleife durchlaufen, bis die Zeitmessung abgeschlossen wird. In jedem Schleifendurchlauf hat man die Wahl zwischen \emph{Zeitmessung pausieren} und \emph{Zeitmessung abschlie�en} oder \emph{Zeitmessung fortsetzen} und \emph{Zeitmessung abschlie�en}, je nachdem ob die Zeitmessung gerade l�uft oder pausiert. Die Implementierung entspricht also einem Kippschalter\index{Kippschalter} mit dem zwischen diesen beiden Zust�nden hin und her geschaltet werden kann.

Bei jeder Aktion wird die Funktion \verb|NotiereZeit| durchlaufen, die mit Hilfe des Befehls \fcmd{SetDateTime} die Uhrzeit notiert und diese mit \fcmd{AddNewKey} als Teil der URL speichert, wobei der fortlaufenden Schl�ssel \verb|ZP(...)| verwendet wird.

Der Befehl \fcmd{SetDateTime} speichert die Systemzeit\index{Systemzeit} in einem vorgegebenen Format in einer Variable und der Befehl \fcmd{AddNewKey} funktioniert so, dass er den Wert nur dann als Teil der URL speichert, wenn dort noch kein Wert f�r diesen Schl�ssel vorhanden ist. Er verh�lt sich damit analog zu den Auswahlabfragen und Texteingaben, die auch nur dann eine Eingabe verlangen, wenn der Wert noch nicht in der URL vorhanden ist. Die Zeitpunkte aller Aktionen werden damit in der URL gespeichert.

Bei aktuell laufender Zeitmessung wird aus den Uhrzeiten zwischen letzter und aktueller Aktion die Zeitdifferenz in Minuten berechnet und zu einer Summe aufaddiert. F�r die Differenzberechnung wird dabei etwas trickreich der Doppelpunkt der Uhrzeiten einfach durch die Zeichenfolge \verb|* 60 +| ersetzt, sodass sich beispielsweise aus dem Ausdruck \verb|12:36 - (11:51)| der Rechenterm $12\cdot 60+36 -(11\cdot 60+51)$ ergibt, der die Differenz der Uhrzeiten in Minuten, also $45$ berechnet.

Bei Beendigung des Skriptes mittels \emph{Zeitmessung abschlie�en} wird dann das vollst�ndige Zeitprotokoll mit den Uhrzeiten, den Differenzminuten und der Gesamtzeitdauer ausgegeben.

\subsection{Meeting-Vorbereitung}\label{MeetingVorbereitung}
Auch das n�chste Beispiel verwendet Befehle zur Abfrage der Systemzeit um eine Ausgabe zu erzeugen, in diesem Fall eine �berschrift f�r das Protokoll des n�chsten Jour fixe\index{Jour fixe}, der an jedem dritten Dienstag stattfindet.

Es ist sinnvoll, Themen f�r regelm��ige Besprechungen im Vorfeld zu sammeln, was man am besten in einem Wiki-Bereich macht, in dem man f�r jedes Meeting eine eigene Protokollseite anlegt, die nach einer festen Konvention mit dem Datum als Teil der �berschrift benannt wird.

Wenn einem nun aktuell ein Thema einf�llt, das man f�r das kommende Meeting dort notieren m�chte, muss dort gegebenenfalls zuerst die entsprechende Seite mit �berschrift angelegt werden. Das dazugeh�rende Datum findet man sicher im Kalender, wo das n�chste Meeting als Teil einer Terminserie eingetragen ist, und f�r die Konvention der �berschrift kann man sich an den letzten Eintr�gen im Wiki orientieren. Einfacher ist es jedoch, das nachfolgen des Skript aufzurufen, das Datum und �berschrift f�r das aus heutiger Sicht n�chste Meeting anhand der Terminregeln automatisch berechnet und bei Bedarf auch noch direkt auf den entsprechenden Wiki-Bereich verlinken k�nnte.

\Beispiel{Jour-fixe-Planer}\label{BspJourFixePlaner}
\begin{verbatim}
~SetCulture de-DE
~SetDateTime dheute=yyyy-MM-dd
~SetDateTime Jahr = yyyy
~DateSet tag = 06.01.$Jahr|dd.MM.yyyy
~Set z=0
~Set nr=0
~DoWhile true
    ~DateAdd tag = $tag|1|d    
    ~DateFormat dvgl=$tag|yyyy-MM-dd
    ~If $dvgl == $dheute
        ~Set zukunft=ja
    ~DateFormat wotag=$tag|ddd
    ~If $wotag==Di
        ~AddTo z+=1
        ~If $z==3
            ~Set z=0
            ~AddTo nr+=1
            ~If $zukunft==ja
                ~ExitLoop
~Loop
~DateFormat erg=$tag|dd.MM.yyyy
@jour fixe WM-XX >>* Wiki-Bereich
    ~AddLink ... | jour fixe WM-XX
    ~AddText  �ffnen.
>>* Falls nicht vorhanden, neue Seite anlegen:
    >|$erg - Jour fixe WM-XX (Nr. $nr/$Jahr)
>>* (findet an jedem 3ten Dienstag statt)
\end{verbatim}

Das Beispiel des \emph{Jour-fixe-Planers}\index{Jour-fixe-Planer} zeigt auch noch weitere Befehle f�r den Umgang mit Datumsangaben\index{Datumsangaben}, angefangen mit \fcmd{SetCulture}, mit dem man die Lokalisierungseinstellungen\index{Lokalisierungseinstellungen} f�r die Ausf�hrung setzen kann, die sich insbesondere auch auf die Datums\index{Datumsformate}- und Uhrzeitformate\index{Uhrzeitformate} auswirken.

In unserem Beispiel werden mit \verb|de-DE| deutschsprachige Einstellungen gesetzt. Danach wird aus der Systemzeit das heutige Datum und zus�tzlich auch noch die Jahreszahl abgeleitet. Im Anschluss wird mit dem \fcmd{SetDateTime}-Befehl die Variable \verb|tag| auf den 06.01.\ des aktuellen Jahres gesetzt, wobei das Format, in dem der zu setzende Wert vorliegt, ebenfalls angegeben werden muss.

Die eigentliche Berechnung erfolgt in einer Schleife, in der die \verb|tag|-Variable mit dem Befehl \fcmd{DateAdd} sukzessive um einen Tag hochgez�hlt wird. F�r den Vergleicht wird dieses Datum mit dem \fcmd{DateFormat}-Befehl auf das gleiche Format gebracht wie \verb|dheute|. Sobald das aktuelle Datum erreicht wird, wird die Variable \verb|zukunft| auf \verb|ja| gesetzt, danach wird der Wochentag �ber das Format \verb|ddd| bestimmt.

In der verschachtelten If-Abfrage wird zuerst gepr�ft, ob es sich um einen Dienstag handelt, und wenn ja, wird ein Z�hler \verb|z| bis auf maximal drei hochgez�hlt. Sobald ein dritter Dienstag gefunden ist, der in der Zukunft liegt, wird die Schleife mit \fcmd{ExitLoop} verlassen. Die Termine werden zus�tzlich �ber das Jahr hinweg durchgez�hlt.

Das gefundene Datum wird dann zusammen mit der laufenden Nummer in die �berschrift f�r die Protokollseite eingebaut und als Codeblock ausgegeben, so dass alles direkt in die Zwischenablage kopiert werden kann.

Auch dieses Skript l�sst sich nach Belieben anpassen z.B.\ durch Variation der Terminregeln, die zus�tzliche Ausgabe der letzten vergangenen oder weiterer zuk�nftigen Termine, durch Ausgabe eines Such-Links, �hnlich wie in Beispiel~\ref{BspSuchlink}, der direkt nach der entsprechenden Seite im Wiki sucht, oder durch die Ausgabe einer Standard-Agenda, die man auf die neu erstellte Seite �bernehmen kann. Diese k�nnte sogar als Jahres-Agenda in einer Datei verwaltet werden und die wiederkehrenden Punkte zum richtigen Zeitpunkt auf den Plan bringen.

Umso einfacher Dinge sind, umso eher und �fter wird gebraucht davon gemacht. Meist sind es die kleinen H�rden, die nicht genommen werden.

\subsection{Asset-Dokumentation}\label{AssetDokumentation}
In einem Softwareunternehmen sollten verschiedene Arten von Assets\index{Assets} dokumentiert werden, um die Ressourcen, das geistige Eigentum und die Infrastruktur des Unternehmens zu verwalten und zu sch�tzen. Die Dokumentation sollte gleicherma�en einen �berblick �ber die vorhandenen Assets geben, als auch die damit verbundenen Aufgaben beschreiben. Je nach Art des Assets sind unterschiedliche Eigenschaften relevant, die jedoch insgesamt in einheitlicher Form dokumentiert werden sollten.

Sofern daf�r nicht eine spezialisierte Asset-Datenbank verwendet wird, ist das Firmenwiki der richtige Ort, um diese Informationen zu verwalten. \FPZ\ kann dabei helfen, die je nach Art ben�tigten Detailinformationen interaktiv abzufragen und daraus das Datenblatt und die Liste der regelm��igen oder aktuell noch ausstehenden Aufgaben zu erstellen.

Das dazugeh�rende Beispiel beginnt mit der Abfrage der allgemeinen Rahmendaten des Assets, zu denen Bezeichnung, Verwendungszweck und vor allem die Art geh�ren, von der die Menge der weiteren Eigenschaften abh�ngt. Das am Ende erstellte Dokument wird Teil des integrierten Managementsystems\index{integriertes Managementsystems} (IMS\index{IMS}) und sollte demzufolge auch einen Dokumenteigent�mer haben, der das Dokument regelm��ig sichtet und aktuell h�lt.

\Beispiel{Asset-Dokumentation (Teil 1)}\label{BspAsset1}
\begin{verbatim}
~Input Bez: Bezeichnung des neuen Assets
?Asset: Welche Art von Asset soll dokumentiert werden?        
    #: Drittanbieterkomponente
    #: Dienst (produktiv)
    #: Dienst (intern)
        ~Set ArtAsset=Dienst (intern)
        ~Set WirdProduktivEingesetzt=Nein        
        ~Set HatLizenz=Ja
        ~Set HatAufrufURL=Ja
        ~Set HatLokaleAdministratoren=Ja
        ~Set HatLokaleInstallation=Nein
        ~Set IstEinCloudDienst=Ja
        ~Set IstInformationsquelle=Nein
        ~Set IstEinGeraet=Nein
        ~Set SetztFachwissenVoraus=Ja
    #: Server (produktiv)
    #: Server (intern)    
    #: Datenblatt
    #: Software
    #: Anleitung/Beschreibung
    #: Wissensdokumentation
    #: Manueller Prozess
    #: Mobiles Endger�t

~Input DEig: Dokumenteigent�mer (Mitarbeiterk�rzel)
~Input VZweck: Wozu wird dass Asset verwendet?
~SetTitle Asset "$Bez"
@Datenblatt >>* Dokumenteigent�mer: $DEig
@Offboarding >> Dokumenteigent�mer pr�fen
@Datenblatt >>* Art: $ArtAsset
@Datenblatt >>* Verwendungszweck
    >* $VZweck
@J�hrl. Aufgaben >> Verwendung/Bedarf pr�fen
@Au�erbetriebenahme >> Asset-Dokumentation archivieren
~Execute
\end{verbatim}

Je nach Art gibt es bestimmte Eigenschaften oder Aspekte, die f�r die Dokumentation relevant sind oder sein k�nnen, und die abgefragt werden m�ssen. Ob eine Eigenschaften f�r eine Art relevant ist oder nicht, wird �ber eine entsprechend benannte Variable gesetzt, die im sp�teren Verlauf abgefragt werden. Durch die Verwendung dieser Schalter\index{Schalter} muss jede Eigenschaft oder jeder Aspekt nur einmalig im Skript implementiert werden und die Reihenfolge bleibt einheitlich.

Im Beispiel oben ist nur der Fall des intern genutzten Dienstes ausimplementiert. F�r alle anderen Typen w�rde man entsprechende Schalter setzen, z.B.\ f�r den Typ \emph{Software} die Eigenschaft \verb|HatLokaleInstallation=Ja|.

Die Implementierung der einzelnen Aspekte sieht prinzipiell so aus, dass man die daf�r ben�tigten Eigenschaften abfragt und im Datenblatt auflistet. Hier k�nnen wieder untergeordnete Aspekte eine Rolle spielen, wie z.B.\ verschiedene Arten von Lizenzen, die man ebenfalls mit Schaltern abbilden kann. F�r Assets mit Lizenz sind im Beispiel schon mal zwei Arten hinterlegt:

\Beispiel{Asset-Dokumentation (Teil 2)}\label{BspAsset2}
\begin{verbatim}
~If $HatLizenz==Ja    
    ?Liz: Welche Art von Lizenz liegt vor?
        #: Frei, kommerziell nutzbar
            ~Set ArtLizenz=Frei, kommerziell nutzbar
            ~Set KommerziellNutzbar=Ja
            ~Set Benutzerbezogen=Nein
            ~Set Kostenpflichtig=Nein
        #: Abo-Lizenz, benutzerbezogen
            ~Set ArtLizenz=Abo-Lizenz, benutzerbezogen
            ~Set KommerziellNutzbar=Ja
            ~Set Benutzerbezogen=Ja
            ~Set Kostenpflichtig=Ja
    @Datenblatt >>* Lizenz
        >* Art: $ArtLizenz
        >* Kommerziell nutzbar: $KommerziellNutzbar
        >* Benutzerbezogen: $Benutzerbezogen
        >* Kostenpflichtig: $Kostenpflichtig
    ~If $Benutzerbezogen==Ja
        @J�hrl. Aufgaben >> Lizenzumfang abgleichen
        @Offboarding >> Lizenzumfang anpassen
    ~If $Kostenpflichtig==Ja
        @J�hrl. Aufgaben >> Preis-Leistung pr�fen
        @Au�erbetriebenahme >> Lizenz k�ndigen
        ~Set ALGegenstand=Lizenzverwaltung
        ~Set ALKey=LIV
        ~GoSub AssetLink
\end{verbatim}

Die meisten Informationen landen im Datenblatt. Bei benutzerbezogenen und kostenpflichtigen Lizenzen leiten sich zus�tzlich Aufgaben f�r die j�hrliche �berpr�fung, das Offboarding und die Au�erbetriebnahme ab. So kann z.B.\ eine Lizenz zur�ckgegeben oder stillgelegt werden, wenn der nutzende Mitarbeiter das Unternehmen verl�sst. Der Funktionsaufruf \verb|AssetLink| am Ende wird in diesem Skript noch h�ufiger verwendet. Er ruft die folgende Funktion auf:

\Beispiel{Asset-Dokumentation (Teil 3)}
\begin{verbatim}
~DefineSub AssetLink
    ?QAL$ALKey: Ist ein Asset f�r "$ALGegenstand" vorhanden?
    #j: Ja
        ~Input AL$ALKey: Link auf Asset f�r "$ALGegenstand"
        @Datenblatt > $ALGegenstand:
        ~AddLink $AL$ALKey| als Asset erfasst
    #n: Nein
        @Datenblatt > $ALGegenstand:
        @ToDos >> Asset f�r $ALGegenstand erfassen
~Return
\end{verbatim}

Diese fragt eine bestimmte variable Eigenschaft des Assets ab, die vorzugsweise wieder als eigenes Asset angelegt sein sollte, so dass mittels Link darauf verwiesen werden kann. In diesem Fall ist es die Lizenzverwaltung, also die Information �ber die Anzahl, Art und Zuteilung der vorhandenen, kostenpflichtigen Lizenzen. Um die Komplexit�t �berschaubar zu halten, ist es sinnvoll, m�glichst viele Teilsysteme als eigene Assets zu erfassen und jeweils darauf zu verweisen. Auf diese Weise spart man sich die verschachtelte Abfrage der damit verbundenen Eigenschaften und Aufgaben und bekommt eine sehr homogene Dokumentation.

Wenn man mit einer derartigen Dokumentation beginnt, muss man davon ausgehen, dass einige der abgefragten Dinge zun�chst nicht vorhanden sind und der Ist-Zustand des Assets nicht alle Anforderungen erf�llt. Entsprechend sollte man diese M�glichkeit in den Auswahlabfragen grunds�tzlich ber�cksichtigen. Wenn dann auf diese Weise festgestellt wird, dass wichtige Dinge fehlen, kann daraus eine Aufgabe in einem ToDo-Abschnitt abgeleitet werden. In gleicher Weise leiten sich aus bestimmten Eigenschaften oder Antworten Aufgaben ab, die regelm��ig, z.B.\ j�hrlich im Zusammenhang mit dem Asset durchgef�hrt werden sollten, und solche, die dann notwendig werden, wenn die mit dem Asset verbundenen Mitarbeiter das Unternehmen verlassen. Am Ende gibt es auch noch Aufgaben, die bei der Au�erbetriebnahme des Assets durchzuf�hren sind, wie die Archivierung der Asset-Dokumentation selbst.

F�r die restlichen oben gesetzten Schalter k�nnte die Erfassung in sehr vereinfachter Form so aussehen:

\Beispiel{Asset-Dokumentation (Teil 4)}\label{BspAsset4}
\begin{verbatim}
~If $HatAufrufURL==Ja
    ~Input AufrufURL: Unter welcher URL wird der Dienst aufgerufen?
    @Datenblatt >>* Aufruf
    >* URL:
    ~AddLink $AufrufURL | $AufrufURL
    @J�hrl. Aufgaben >> Verf�gbarkeit pr�fen
~If $HatLokaleAdministratoren==Ja
    ~Input Admins: Mitarbeiter mit Admin-Zugang und -kenntnissen:
    @Datenblatt >>* Administration
        > Mitarbeiter mit Zugang und Wissen: $Admins
    ~Set ALGegenstand=Admin-Wissen
    ~Set ALKey=ADW
    ~GoSub AssetLink
    @Offboarding >> Admin-Zugang �bergeben
    ~Set ALGegenstand=Admin-Notfallzugang
    ~Set ALKey=ANZ
    ~GoSub AssetLink
~If $IstEinCloudDienst==Ja
    ~Input AnbCloudName: Anbieter des Cloud-Dienstes
    ~Input AnbCloudWeb: Internetseite Anbieter
    @Datenblatt >>* Anbieter
    >* Name: $AnbCloudName
    >* Webseite:
    ~AddLink $AnbCloudWeb| $AnbCloudWeb
    @J�hrl. Aufgaben >> Weitere Verf�gbarkeit pr�fen
    @Datenblatt >>* Risiko: Wegfall des Dienstes
    ~Set ALGegenstand=Wegfall-Notfallplan
    ~Set ALKey=WNP
    ~GoSub AssetLink
    @Au�erbetriebenahme >> Konfigurationsdaten sichern
~If $SetztFachwissenVoraus==Ja
    ~Set ALGegenstand=Nutzungswissen
    ~Set ALKey=WFN
    ~GoSub AssetLink

~MoveSection J�hrl. Aufgaben->J�hrliche Aufgaben
~MoveSection Offboarding->Offboarding-Aufgaben
~MoveSection Au�erbetriebenahme->Aufgaben bei Au�erbetriebenahme
~MoveSection ToDos->Aktuelle ToDos
\end{verbatim}

Die Menge der abgefragten Informationen sollte sich prim�r daran orientieren, was in Ausnahmesituationen ben�tigt wird, um den Kernbetrieb am Laufen zu halten. Auch das Ausfallrisiko auf Mitarbeiterseite ist nicht zu untersch�tzen. Mitarbeiter, die zentrales Wissen oder Zug�nge zu essentiellen Diensten haben, k�nnen verunfallen und f�r l�ngere Zeit vollst�ndig ausfallen. Insofern sollte immer darauf geachtet werden, dass relevantes Wissen festgehalten wird und dass essentielle Zug�nge notfallm��ig hinterlegt werden.

Die vier \fcmd{MoveSection}-Befehle am Ende stellen sicher, dass die einzelnen Abschnitte immer in der gleichen Reihenfolge aufgelistet werden. Zus�tzlich wird dabei die Abschnitts�berschrift komplett ausgeschrieben.

Das Ergebnisdokument kann mit �berschrift und allen Abschnitten auf eine eigene Wiki-Seite innerhalb der Asset-Dokumentation �bernommen werden. Der vorzugsweise nach Art gegliederte Unterbereich innerhalb der Dokumentation und die genaue Vorgehensweise k�nnten hierbei in einem eigenen Anleitungsabschnitt beschrieben werden. Ebenso sollte man die Abfragen und Texteingaben im Skript noch durch zus�tzliche Hilfezeilen selbsterkl�render gestalten, so dass man am Ende ein Skript bekommt, mit dem die jetzt schon f�r das entsprechende Asset verantwortlichen Mitarbeitern direkt in die Lage versetzt werden, diese in der gew�nschten Form zu dokumentieren.

Diese Dokumentation erm�glicht es dann, j�hrlich die vorhandenen Assets zu sichten und auf Notwendigkeit zu pr�fen. Nicht selten werden Vorgehensweisen oder Softwareprodukte eingef�hrt, die sich dann irgendwann von selbst ausschleichen, ohne dass eine Au�erbetriebnahme stattgefunden hat. Das ist vor allem dann unsch�n, wenn weiterhin Geld daf�r bezahlt wird. Die j�hrliche Sichtung der Assets, zusammen mit den notwendigen Wartungsaufgaben ist wichtig, um z.B.\ rechtzeitig dar�ber informiert zu sein wenn, etwa aufgrund einer Abk�ndigung, Nutzungseinschr�nkungen absehbar werden.

Gerade solche Rechercheaufgaben nach neuen oder abgek�ndigten Versionen lassen sich heute -- sofern es die Unternehmensrichtlinien zulassen -- mit einer einfachen Anfrage an einen Chatbot erledigen, z.B.\
\emph{F�r welche Versionen des Windows Server Betriebssystems l�uft der Support als n�chstes aus und gibt es schon Ank�ndigungen f�r neue Versionen?} Die Erzeugung des passenden Eingabe-Promptes auf Basis von Komponentenname und Komponententyp l�sst sich dabei sehr gut in ein \FPZb-Skript f�r diese Art von Wartungsaufgaben integrieren, dessen Aufruf wiederum als Link direkt in der Asset-Dokumentation hinterlegt ist. Das Formulieren effektiver Prompts f�r KI-Anfragen, auch \emph{Prompt Engineering}\index{Prompt Engineering} genannt, ist eine zunehmend wichtige und auch anspruchsvolle Aufgabe, und wenn man einmal einen guten Prompt im Zusammenhand mit einer wiederholt auftretenden Aufgabenstellung gefunden hat, sollte man diesen unbedingt in die dazugeh�rende Anleitung integrieren.

Durch die hohe Flexibilit�t eines solchen Skriptes, lassen sich die abgefragten Aspekte so anpassen, dass sie punktgenau auf die f�r ein Unternehmen bedeutsamen Assets und die daf�r angelegten Strukturen ausgerichtet sind. Die Dokumentation im Wiki kann dann bei �nderungen unabh�ngig von dem Skript, mit dem sie erstellt wurde, angepasst werden.

Der entscheidende Vorteil bei einer interaktiven Anleitung f�r die Dokumentation von Dingen ist jedoch die M�glichkeit, Fragestellungen und daraus abgeleitete Ma�nahmen direkt damit zu verbinden, und deren Durchf�hrung gut zu unterst�tzen. So kann ein Skript, das urspr�nglich nur f�r die Dokumentation gedacht war, und nur Informationen abgefragt hat, schrittweise bis zu einem Protokoll erweitert werden, dass am Ende garantiert, dass das Dokumentierte definierte Standards erf�llt.

\subsection{Terminplanung}\label{Terminplanung}
Unser letztes Beispiel im Abschnitt Organisationswerkzeuge ist ein kleines Terminplanungswerkzeug, das sch�n zeigt, wie sich mit ein bisschen Programmierung Aufgaben so l�sen lassen, dass es f�r alle Beteiligten von Vorteil ist.

Die Ausgangssituation d�rfte in einem Unternehmen recht h�ufig, aber mindestens einmal pro Jahr und Abteilung vorkommen: f�r eine Menge von Mitarbeiter sollen innerhalb eines Zeitraums Einzeltermine geplant werden, so dass am Ende f�r jeden Mitarbeiter der Menge ein Termin mit Raum eingeplant ist.

Die normale Vorgehensweise w�re nun, mit Hilfe des eigenen Terminkalenders und denen der Mitarbeiter und R�ume einen Termin nach dem anderen zu planen, was bedeutet, dass man mindestens immer drei Kalender offen haben muss. Unsch�n wird das ganze dann gegen Ende der Planung, wenn man f�r die verbleibenden Mitarbeiter keine m�glichen Termine mehr findet, und man somit Anfang und Mittelteil der Planung wieder umstellen muss.

Die deutlich einfachere Vorgehensweise besteht darin, zun�chst nur Termine f�r sich selbst und den Raum einzuplanen, daf�r eventuell auch einen mehr als ben�tigt wird. Anschlie�end ver�ffentlicht man die Liste dieser Termine im Wiki und bittet die Mitarbeiter, sich dort Ihre Wunschtermine auszuw�hlen, also je nach Gesamtzahl drei bis f�nf Wunschtermine pro Mitarbeiter. Hierzu tr�gt sich der Mitarbeiter einfach in die Tabelle mit seinem K�rzel hinter dem jeweiligen Termin ein.

Bei der Auswahl kann zus�tzlich noch eine Pr�ferenz angegeben werden, um auszudr�cken, ob der jeweilige Termin die erste, zweite oder dritte Wahl f�r den Mitarbeiter ist. Dies wird einfach durch die entsprechende Ziffer hinter dem K�rzel angegeben.

F�r eine Auswahl von f�nf Terminen k�nnten die W�nsche von vier Mitarbeitern mit den K�rzeln AB, CD, EF und GH wie folgt aussehen:

\begin{center}
\begin{tabular}{|c|l|}\hline
Termin 1& AB2\\
Termin 2& CD, EF1\\
Termin 3& EF2, AB1\\
Termin 4& GH, CD\\
Termin 5& CD, EF3\\ \hline
\end{tabular}
\end{center}

In diesem Fall hat der Mitarbeiter AB nur zwei Termine angegeben und GH nur einen, eventuell weil er von ausw�rts anreisen muss. Der Mitarbeiter CD hat keine Pr�ferenzen, die drei von ihm gew�hlten Termine sind f�r ihn gleichwertig.

Nach Ablauf der Wahlfrist werden nun einfach die gemachten Angaben in die Eingabefelder des folgenden Skriptes �bertragen:

\Beispiel{Terminplaner}\label{BspTerminplaner}
\begin{verbatim}
~SetStopCounter 100000; 3000000
~Input Anz: Anzahl Termine
~Execute
@Terminpr�ferenzen >>_ $Anz Termine
~Set idx=0
~DoWhile $idx<$Anz
    ~AddTo idx+=1
    ~Input Prf$idx: Termin $idx - Pr�ferenzen (kommagetrennt)
    >>* Termin $idx: $Prf$idx    
~Loop
~Set BestPSum=10000
~Set BestAnzV=0
~Set TidxS=1
~Set MAVS=
~Set AnzVS=0
~Set PSumS=0
~GoSub SucheLoesung; BaseKey=S
@Ergebnis >>_ $BestAnzV Termine verplant, $BestPSum Punkte
~Set idx=0
~DoWhile $idx<$Anz
    ~AddTo idx+=1
    >>* Termin $idx: $FErg($idx)
~Loop

// Tidx$BaseKey = N�chster Terminindex
// MAV$BaseKey = Liste der schon verplanten Mitarbeiter
// PSum$BaseKey = Aktuelle Pr�ferenz-Summe
// R�ckgabe: FErg(1..Anz) = Belegung der Termine
~DefineSub SucheLoesung    
    ~Split VerplanteMA=$MAV$BaseKey|/
    ~Calculate AnzV$BaseKey=$VerplanteMA(0)-1
    ~If $Tidx$BaseKey<=$Anz
        ~Split pr$BaseKey=$Prf$Tidx$BaseKey|,        
        ~ForEach pi in pr$BaseKey
            ~RegExMatch pw = $pi|([A-Z][A-Z][a-z]*)([0-9]?)
            ~If $pw(0)
                ~Set ku=$pw(1)
                ~Set va=$pw(2)
                ~If $va==
                    ~Set va=1
                ~If $MAV$BaseKey/!~/$ku/
                    ~Set Erg($Tidx$BaseKey)=$ku                    
                    ~Calculate Tidx$BaseKeyX$ku=$Tidx$BaseKey+1
                    ~Calculate PSum$BaseKeyX$ku=$PSum$BaseKey+$va
                    ~Set MAV$BaseKeyX$ku=$MAV$BaseKey/$ku
                    ~GoSub SucheLoesung; BaseKey=$BaseKeyX$ku
        ~Loop
        ~CalculateExpression restpot=$Anz-$BestAnzV
            __-($Tidx$BaseKey-$AnzV$BaseKey)
        ~If $restpot>=0
            ~Set Erg($Tidx$BaseKey)=Nicht besetzt
            ~Calculate Tidx$BaseKey0=$Tidx$BaseKey+1
            ~Set PSum$BaseKey0=$PSum$BaseKey
            ~Set MAV$BaseKey0=$MAV$BaseKey
            ~GoSub SucheLoesung; BaseKey=$BaseKey0        
    ~Else
        ~If $AnzV$BaseKey>$BestAnzV || $AnzV$BaseKey==$BestAnzV
            __ && $PSum$BaseKey<$BestPSum            
            ~Set BestAnzV=$AnzV$BaseKey
            ~Set BestPSum=$PSum$BaseKey    
            ~Set idx=0
            ~DoWhile $idx<$Anz
                ~AddTo idx+=1
                ~Set FErg($idx)=$Erg($idx)
            ~Loop        
~Return
\end{verbatim}

Da das Skript zur L�sungsfindung exponentiell viele Variationen durchprobiert, werden gleich zu Beginn die Stop-Z�hler\index{Stop-Z�hler} f�r die maximale Anzahl an Schleifendurchl�ufen und Befehlsausf�hrungen mit dem Befehls \fcmd{SetStopCounter} auf sehr gro�e Werte gesetzt.

Zuerst wird die Anzahl der Termine abgefragt, dann f�r jeden Termin die Pr�ferenzen, wie �blich kommagetrennt.

Die Suche nach der passenden L�sung erfolgt mit Hilfe der Funktion \verb|SucheLoesung|, die als �bergabeparameter den Index des n�chsten zu planenden Termins, die Liste der schon verplanten Mitarbeiter und die Zwischensumme der Pr�ferenzwerte �bergeben bekommt.  Die Anzahl der schon verplanten Termine wird mit Hilfe des \fcmd{Split}-Befehls aus der Aufz�hlung der schon verplanten Mitarbeiter berechnet. 

Immer wenn der Index des letzten Termins erreicht ist, wird die aktuelle L�sung mit der bisher besten L�sung verglichen, wobei prim�r die Anzahl der geplanten Termine z�hlt, und bei gleicher Anzahl eine m�glichst geringe Pr�ferenzsumme. Ist eine neue beste L�sung gefunden, wird diese in das Feld \verb|FErg(...)| �bertragen.

Ansonsten werden alle Mitarbeiter durchlaufen die sich f�r den als n�chstes zur Planung anstehenden Termin eingetragen haben. Die K�rzel werden mit Hilfe eines regul�ren Ausdrucks vom Pr�ferenzwert getrennt. Wenn keiner angegeben ist, wird der Wert 1 gesetzt. Mit Hilfe des Enth�lt-nicht-Operators\index{Enth�lt-nicht-Operator} \verb|!~| wird dann gepr�ft, ob der entsprechende Mitarbeiter noch nicht in der Menge der bisher verplanten Mitarbeiter enthalten ist, wobei das Trennzeichen \verb|/| verwendet wird. In diesem Fall werden die Parameter f�r einen weiteren Funktionsaufruf gesetzt, wobei die \verb|BaseKey|-Variable um den Buchstaben X und das Mitarbeiterk�rzel erweitert werden.

Nach Durchlauf aller Mitarbeiter wird auch noch die Variante gepr�ft, bei der der anstehende Termin nicht besetzt wird. Gerade wenn mehr Termine angeboten werden, als Mitarbeiter vorhanden sind, m�ssen Termine unbesetzt bleiben.

Da dadurch die Menge der zu untersuchende L�sungen noch mal deutlich erh�ht wird, wird hier noch zus�tzlich das Restpotenzial $R$ abgepr�ft. Dieses berechnet sich �ber die Formel $R=A-A_B-(i-A_T)$, wobei $A$ die Anzahl der angebotenen Termine ist, $A_B$, die Anzahl der verplanten Termine der bislang besten L�sung, $i$ der aktuelle Index und $A_T$ die Anzahl der bis jetzt verplanten Termine. Ist $R\ge0$, kann trotz nicht besetzen des anstehenden Termins potentiell noch eine bessere L�sung gefunden werden, ansonsten ist das ausgeschlossen und diese Variante muss nicht weiterverfolgt werden.

Am Ende wird die beste gefundene L�sung als Ergebnis ausgegeben. Das Skript kann auch mit den Daten aus unserem Beispiel oben umgehen und findet zielsicher die beste L�sung mit vier verplanten Terminen und einer Pr�ferenzsumme von 4, d.h.\ jeder Mitarbeiter bekommt einen Termin seiner ersten Wahl:

\begin{center}
\begin{tabular}{|c|l|}\hline
Termin 1& nicht besetzt\\
Termin 2& EF\\
Termin 3& AB\\
Termin 4& GH\\
Termin 5& CD\\ \hline
\end{tabular}
\end{center}

Der Vorteil dieser Methode liegt darin, dass die l�stige Suche nach freien Terminen auf die einzelnen Mitarbeiter aufgeteilt wird, die aber daf�r nur jeweils ihren eigenen Kalender sichten m�ssen. Als Ausgleich daf�r bekommen sie die M�glichkeit, mit der Pr�ferenz Einfluss auf die Terminauswahl zu nehmen.

Dadurch, dass auch die Reihenfolge, in der die W�nsche eingetragen werden so gut wie keine Rolle spielt, ist dieser Ansatz auch besser als die Startschuss-Variante, bei dem sich jeder nur f�r genau einen Termin eintragen darf, der noch nicht gew�hlt wurde, so dass die Mitarbeiter benachteiligt werden, die nicht unmittelbar auf die Aufforderung reagieren k�nnen.

\section{Mitgestaltung}
Wie schon mehrfach beschrieben, sollte man Skripte nicht einfach nur als Programmcode sehen, der durch die Umgebung ausgef�hrt werden kann und mit dem sich bestimmte Aufgaben einfach l�sen lassen, sondern man sollte Skripte vor allem auch als ein Medium betrachten, das es erm�glicht, fachliches Wissen so zu hinterlegen, dass es bei der Anwendung des Skriptes automatisch ber�cksichtigt wird. Es ist nicht notwendig, dass der Anwender wei�, wie man die Aufgabe richtig umsetzt, so lange ihm das Skript keine M�glichkeit gibt, sie falsch umzusetzen.

Ausgehend davon sollte die Zielsetzung sein, m�glichst viel in den K�pfen vorhandenes Wissen in so eine Form zu bringen. In diesem Abschnitt zeigen wir anhand von Beispielen, wie das m�glich ist, auch ohne dass man jedem Mitarbeiter die \FPZb-Programmierung ans Herz legen muss.

\subsection{Angeleitete Skript-Erweiterungen}
Bei Auswahlabfragen hat man die Auswahl zwischen verschiedenen Werten. Was aber kann der Anwender machen, wenn der von ihm gew�nschte Wert nicht in der Liste aufgef�hrt ist? F�r diesem Fall k�nnte man einen entsprechenden zus�tzlichen Auswahlwert \emph{Nicht in Liste}\index{nicht in Liste} erg�nzen, und versuchen, diesen neuen Fall so gut wie m�glich im Skript zu behandeln. 

Meist werden in Abh�ngigkeit von Auswahlwerten nur bestimmte Variablen oder Schalter gesetzt, die dann in einem anderen, schon vorhandenen Teil des Skriptes, weiterverarbeitet werden. In diesem Fall kann man die entsprechenden Werte f�r die zu setzenden Variablen direkt abfragen und hat so den allgemeinen Fall behandelt. Sch�ner w�re es jedoch, wenn die verschiedenen Eingaben, die eventuell mit einigem Aufwand zusammengetragen werden m�ssen, f�r den n�chsten Anwender direkt im Skript als ausw�hlbarer Wert eingebaut w�ren.

\FPZ\ kann selbst keine �nderung an den Skripten durchf�hren, aber es kann sehr gut dazu anleiten.

Betrachten wir die Fragestellung, unter welcher Lizenz ein Softwareprodukt bereitgestellt wird, um daraus bestimmte Eigenschaften abzuleiten, wie die Bezeichnung, eine Link auf eine Beschreibung der Lizenz, die Information ob diese Lizenz den Einsatz im kommerziellen Umfeld erlaubt und ob sie kostenpflichtig ist.

Hier bietet es sich an, die Auswahlabfrage zun�chst mit einem Wert zu beginnen, und weitere Lizenzen dann hinzuzunehmen, wenn sie tats�chlich ben�tigt werden. Diese Abfrage k�nnte man wie folgt aufbauen:

\Beispiel{Lizenzabfrage}\label{BspLizenzabfrage}
\begin{verbatim}
?Lizenz: Unter welcher Lizenz wird die Software bereitgestellt?
    #: MIT-Lizenz
        ~Set BezeichnungLizenz=MIT-Lizenz
        ~Set LinkAufLizenzbeschreibung=
            __https://de.wikipedia.org/wiki/MIT-Lizenz
        ~Set ErlaubtKomVerw=Ja
        ~Set IstKostenpflichtig=Nein
    // Oberhalb von hier neue Eintr�ge einf�gen
    #x: Nicht in Liste
        @Skript-Erweiterung >>_ Es wird empfohlen, das Skript
            __ um die nicht vorhandenen Eintr�ge zu erweitern. 
            __ Gehe dazu wie folgt vor:
        >> �ffne die Skriptdatei in einem Texteditor:
            >|$ScriptFilePath
        >> F�ge vor Zeile $LineNumber-7 folgenden Codeblock ein:
            ~Input ILizenz: Unter welcher Lizenz wird die 
                __ Software bereitgestellt?
            >|$Chr(9)#: $ILizenz
            ~Input IBez: Bezeichnung Lizenz
            >|$Chr(9)$Chr(9)~Set BezeichnungLizenz=$IBez
            ~Set BezeichnungLizenz=$IBez
            ~Input ILink: Link auf Lizenzbeschreibung
            >|$Chr(9)$Chr(9)~Set LinkAufLizenzbeschreibung=$ILink
            ~Set LinkAufLizenzbeschreibung=$ILink
            ?QKom: Erlaubt kommerzielle Verwendung?
                #j: Ja
                    >|$Chr(9)$Chr(9)~Set ErlaubtKomVerw=Ja
                    ~Set ErlaubtKomVerw=Ja
                #n: Nein
                    >|$Chr(9)$Chr(9)~Set ErlaubtKomVerw=Nein
                    ~Set ErlaubtKomVerw=Nein
            ?QKost: Ist kostenpflichtig?
                #j: Ja
                    >|$Chr(9)$Chr(9)~Set IstKostenpflichtig=Ja
                    ~Set IstKostenpflichtig=Ja
                #n: Nein
                    >|$Chr(9)$Chr(9)~Set IstKostenpflichtig=Nein
                    ~Set IstKostenpflichtig=Nein
        >> Speichere die Datei ab.
@Ausgabe >>* Die Software wird unter der Lizenz "$BezeichnungLizenz"        
        __ bereitgestellt.
    >* Link: 
    ~AddLink $LinkAufLizenzbeschreibung | $LinkAufLizenzbeschreibung
    >* Kommerzielle Verwendung m�glich? - $ErlaubtKomVerw
    >* Kostenpflichtig? - $IstKostenpflichtig
\end{verbatim}

Die Auswahl von \emph{Nicht in Liste} fragt die ben�tigen Werte ab und setzt die Variablen in der gleichen Weise, wie es die anderen Eintr�ge machen, so dass das Skript seine Aufgabe erf�llen kann, z.B.\ die Dokumentation der abgefragten Lizenz mit den ganzen Eigenschaften.

Es passiert aber noch mehr. Als Teil der Ausgabe wird in einem eigenen Abschnitt \emph{Skript-Erweiterung}\index{Skript-Erweiterung} eine Schritt-f�r-Schritt-Anleitung gegeben, die beschreibt, wie man die neu erfasste Lizenz permanent in das Skript einbauen kann. Hierbei wird beschrieben, wo man die Skript-Datei findet, und welche Zeilen man wo einf�gen muss. Der einzuf�gende Codeblock wird vollst�ndig ausgegeben und kann direkt in die Zwischenablage �bernommen werden. Die Erweiterung kann damit ohne jede Kenntnis der \FPZb-Programmiersprache umgesetzt werden.

F�r die Zusammenstellung der ben�tigten Informationen wird zum einen die Variable \cidxvar{ScriptFilePath} verwendet, die den Dateipfad des aktuellen Skriptes ausgibt. Dieser ist aus Sicht des Servers, auf dem \FPZ\ l�uft, und muss gegebenenfalls gegen einen ausgeschriebenen Pfad ausgetauscht werden, wenn der Zugriff auf die Datei f�r den Anwender �ber einen anderen Pfad erfolgen soll. Die Zeilennummer\index{Zeilennummer} wird mit Hilfe der Variablen \cidxvar{LineNumber} abgefragt, die stets die Nummer der entsprechenden Zeile zur�ckgibt. M�chte man sich auf eine Zeile $n$ Zeilen dar�ber beziehen, kann f�r diese Variable direkt ein Korrekturwert von $-n$ angegeben werden.

Das Tabulatorzeichen\index{Tabulatorzeichen} f�r die Einr�ckung des auszugebenden Skriptcodes wird �ber die Sequenz \cidxvar{Chr(9)} angegeben, da unmaskierte Tabulatorzeichen beim Einlesen von Skripten durch vier Leerzeichen ersetzt werden.

\subsection{Metaskripting}\label{Metaskripting}
Jetzt kann man zu Recht bem�ngeln, dass der Scriptcode, der f�r die \emph{Nicht in Liste}-Auswahlm�glichkeit erstellt werden muss, selbst sehr komplex und relativ umfangreich ist. Die H�rde eine solche M�glichkeit zu schaffen scheint damit recht hoch zu sein, und selbst wenn man sich daf�r bei einer Vorlage bedienen kann, sind immer noch einige Anpassungen notwendig, da sich die abzufragenden Variablen ja von Fall zu Fall unterscheiden.

Die L�sung f�r dieses Problem ist, man ahnt es vermutlich schon, ein \FPZb-Skript, das einem diese Arbeit abnimmt.

Das sogenannte Metaskripting\index{Metaskripting}, also das Erstellen von Skriptcode der neuen Skriptcode erzeugt, ist ein m�chtiges Instrument und l�sst sich an vielen Stellen einsetzen, sogar �ber verschiedene Programmiersprachen hinweg. Es gibt Beispiele f�r produktiv eingesetzte \FPZb-Skripte, die SQL-Anweisungen erzeugen, der wiederum C\#-Code ausgegeben.

F�r unseren Anwendungsfall der angeleiteten Skript-Erweiterung ben�tigen wir folgende Eingaben: die Fragestellung, einen Schl�ssel, eine Liste von Werten, die entweder als Texteingabe oder Schalter abgefragt werden sollen, und die Information in welcher Tiefe der Code einger�ckt sein soll.

Die Abfrage erfolgt entsprechend mit vier Texteingaben, wobei die Menge der abzufragenden Werte �ber eine kommagetrennte Aufz�hlung gemacht wird und die Unterscheidung zwischen Textangaben und Schaltern durch die Angabe eines Fragezeichens bei jedem Wert getroffen wird. Die Einr�ckung wird durch eine Sequenz von T-Zeichen angegeben, wobei sp�ter jedes T durch Tabulatorzeichen ersetzt wird.

Das Ergebnis sieht wie folgt aus:

\Beispiel{Meta-Erweiterung}\label{BspMetaErweiterung}
\begin{verbatim}
~Input Frage: Welche Frage soll ausgegeben werden?
~Input Key: Schl�ssel f�r die Frage
~Input WListe: Abzufragende Werte (kommagetrennt, @ wenn leer)
~AddHelpLine Werte mit "?" werden als Schalter abgebildet.
~AddHelpLine Der Wert mit "=" wird auf den Auswahlwert gesetzt.
~Input Tabs: Einr�ckung (T pro Tab, @ sonst)
~CamelCase CKey=$Key
~Set T=$Chr(9)
~Set ST=$Chr(36)Chr(36)Chr(9)
~Replace EE=$Tabs|@->
~Replace FF=$EE|T->$Chr(36)Chr(36)Chr(9)
~Replace EE=$EE|T->$Chr(9)
@Erweiterung >>* Auswahlabfrage
>|$EE?$CKey: $Frage
>|$EE$T// Oberhalb von hier neue Eintr�ge einf�gen
>|$EE$T#x: Nicht in Liste
>|$EE$T$T@Skript-Erweiterung >>_ Es wird empfohlen, das
    __ Skript um die nicht vorhandenen Eintr�ge zu erweitern.
    __ Gehe dazu wie folgt vor:
>|$EE$T$T>> �ffne die Skriptdatei in einem Texteditor:
>|$EE$T$T$T>|$Chr(36)ScriptFilePath
>|$EE$T$T>> F�ge vor Zeile $Chr(36)LineNumber-5 folgenden
    __ Codeblock ein:
>|$EE$T$T$T~Input I$CKey: $Frage
>|$EE$T$T$T>|$FF$ST#: $I$CKey
~If $WListe!=@
    ~Split Werte=$WListe|,
    ~ForEach widx in Werte
        ~CamelCase vidx=$widx
        ~If $widx~?
            @Erweiterung >|$EE$T$T$T?Q$vidx: $widx
            >|$EE$T$T$T$T#j: Ja
            >|$EE$T$T$T$T$T>|$FF$ST$ST~Set $vidx=Ja
            >|$EE$T$T$T$T$T~Set $vidx=Ja
            >|$EE$T$T$T$T#n: Nein
            >|$EE$T$T$T$T$T>|$FF$ST$ST~Set $vidx=Nein
            >|$EE$T$T$T$T$T~Set $vidx=Nein
            ~If $ErsterSchalter!=Ja
                @IfBlocks >>* If-Abfragen f�r die Schalter
                ~Set ErsterSchalter=Ja
            @IfBlocks >|$EE~If $$vidx==Ja
            >|$EE$T//$widx
        ~ElseIf $widx~=
            @Erweiterung >|$EE$T$T$T>|$FF$ST$ST~Set $vidx=$I$CKey
            >|$EE$T$T$T~Set $vidx=$I$CKey
        ~Else            
            @Erweiterung >|$EE$T$T$T~Input I$vidx: $widx
            >|$EE$T$T$T>|$FF$ST$ST~Set $vidx=$I$vidx
            >|$EE$T$T$T~Set $vidx=$I$vidx
    ~Loop
@Erweiterung >|$EE$T$T>> Speichere die Datei ab.
~MoveSection IfBlocks->Erweiterung
\end{verbatim}

Um den Code schmal zu halten, werden am Anfang die Hilfsvariablen \verb|T| und \verb|ST| mit den h�ufig ben�tigten Sequenzen belegt. Damit die Tabulatorzeichen (\cidxvar{Chr(9)}) im innersten Code bis zum letztendlichen Einf�gen in das Zielskript noch maskiert bleiben, wird das \verb|$|-Zeichen der \cidxvar{Chr(...)}-Sequenzen teilweise wiederholt durch \cidxvar{Chr(36)} ersetzt. In gleicher Weise werden auch die \verb|$|-Zeichen in den Aufrufen der Systemvariablen \cidxvar{ScriptFilePath} und \cidxvar{LineNumber} durch \cidxvar{Chr(36)} ersetzt, so dass diese erst bei der Ausf�hrung im Zielskript aufgel�st werden und nicht bereits in der Ausgabe dieses Skriptes.

F�r die Schalter, also die Werte, die mit einem Fragezeichen gekennzeichnet wurden, werden in einem zweiten Codeblock noch If-Abfragen erstellt. Damit kann man, �hnlich wie in den Beispielen~\ref{BspAsset1} bis \ref{BspAsset4}, f�r jeden Schalter individuellen Code ausf�hren, wenn dieser gesetzt ist.

Metaskripting ist ohne Frage anspruchsvoll, aber daf�r ist der Nutzen aufgrund der Multiplikatorwirkung recht gro�, und wenn ein solches Skript am Ende dazu f�hrt, dass alle Mitarbeiter in der Lage sind, die produktiv eingesetzten Skripte permanent zu erweitern und ihr Wissen dort einzupflegen, dann ist das den Aufwand auf jeden Fall wert.

Auch hier verh�lt es sich �hnlich wie bei der Entwicklung von Anwendungscode. Wird etwas nur einmal ben�tigt, kann es der Spezialist machen. Wird etwas mehrfach ben�tigt, kann der Spezialist passendes Werkzeug daf�r anfertigen, mit dem die Aufgabe einfacher umgesetzt werden kann.

\subsection{Wissensdokumentation}\label{Wissensdokumentation}
Die Dokumentation von Wissen oder auch die Verwaltung von Wissen ist eine der wichtigsten Aufgaben in einem Unternehmen, das seine Wertsch�pfung in erster Linie aus dem Wissen der Mitarbeiter generiert.

Fast alle Abl�ufe, die im Arbeitsalltag stattfinden, setzen Wissen voraus, ebenso die meisten Anwendungen und Werkzeuge, die dabei eingesetzt werden.

Dort, wo dieses Wissen jeden Tag, oder zumindest einmal in der Woche ben�tigt wird, wird es sich von selbst verfestigen, so dass man hier keinen Aufwand betreiben muss. Anders sieht es aus bei Wissen, dass nur gelegentlich ben�tigt wird, und bei dem zus�tzlich gar nicht offensichtlich ist, ob es hier etwas zu wissen gibt oder nicht. Wenn dieses Wissen fehlt, laufen Dinge schief, was im besten Fall nur Mehraufwand bedeutet und im schlimmsten Fall zu problematischen Fehler f�hren kann.

Wer nicht wei�, wie eine bestimmte Art von Skript erstellt wird, wird keines erstellen k�nnen, wer nicht wei�, welches Risiko eine bestimmte Entwicklung mit sich bringt, wird nicht darauf achten k�nnen und die Eintrittswahrscheinlichkeit erh�hen, und wer nicht wei�, welche Richtlinien im Umgang mit Kundendaten zu beachten sind, wird m�glicherweise einen Datenschutzversto� verursachen.

Wie schon mehrfach beschrieben, lassen sich mit Hilfe von Skripten gute Anleitungen und Hilfsmittel erstellen, die das wesentliche Detailwissen �ber bestimmte Arbeitsschritte enthalten und so den einzelnen Mitarbeiter auch wissensm��ig entlasten.

Auch das Grundwissen zu allgemeinen Themengebieten, wie C\#-Programmierung oder Kenntnis der SQL-Syntax muss nicht unbedingt firmenspezifisch verwaltet werden, da es ja eine generelle Vorstellung davon gibt, was es bedeutet, dieses Wissen zu haben, und man es ggf.\ mit zahlreichen Schulungen oder Tutorials abgleichen und auffrischen kann.

Generell sollte das vorausgesetzte Wissen auf die Punkte beschr�nkt werden, die unternehmensspezifisch sind, die einem helfen, die richtigen Entscheidungen zu treffen, Fehler zu vermeiden und das f�r bestimmte Arbeitsschritte ben�tigte Detailwissen, bzw.\ die Anleitungen daf�r zielsicher aufzufinden. Das ist immer noch eine ganze Menge, und f�r jede spezielle Aufgabe und f�r jeden Verantwortungsbereich gibt es eigenes Wissen, das ben�tigt wird.

Wie aber verwaltet man dieses Wissen und wie erf�hrt ein neuer Mitarbeiter, was er alles wissen muss? Wie l�sst sich �berpr�fen ob ein Mitarbeiter alles wei�, was er wissen sollte? Die Verantwortung �ber diese Punkte liegt in der Regel beim direkten Vorgesetzten, der aber auch nicht in den einzelnen Menschen hineinschauen kann. Meist behilft man sich damit die zwingend einzuhaltenden Richtlinien schriftlich festzuhalten und f�r das restliche Wissen grobe Themengebiete aufzulisten, die zur Orientierung verwendet werden k�nnen. 

Ein im Gegensatz dazu sehr expliziter Ansatz f�r die Wissensverwaltung\index{Wissensverwaltung} ist die Auf\-lis\-tung von Wissen in Form von Reflektionsaussagen\index{Reflektionsaussagen}. Eine Reflektionsaussage ist dabei eine Aussage, die in der Form \emph{Ich wei�, \dots} formuliert wird, z.B.\ \emph{Ich wie man in einem PowerShell-Skript auf Befehlszeilenargumente zugreift}.

Die Aussage dient, wie die Bezeichnung schon sagt, dazu, das eigene Wissen zu reflektieren, also dar�ber nachzudenken, ob man diese Aussage mit Ja oder mit Nein beantworten kann. Der gro�e Vorteil bei dieser Form der Erfassung liegt darin, dass es sehr einfach ist, Wissen in dieser Form zu erfassen, da man ja nicht den Inhalt des Wissens erfassen muss, sondern nur einen Verweis darauf. Das ist sogar dann m�glich wenn man die Aussage f�r sich selbst gar nicht mit Ja beantworten kann. Dar�berhinaus erlaubt es die Formulierung, das abgefragte Wissen sehr genau abzugrenzen, so dass man in der Regel klar verstehen kann, was gemeint ist.

Genau diese Einfachheit ist wichtig, um mit m�glichst wenig Aufwand m�glichst alles an Wissen zu dokumentieren, das man f�r wichtig und relevant h�lt. Vor allem in dem Moment, wo man feststellt, dass aufgrund fehlenden Wissens etwas schief gelaufen ist, sollte man umgehend daf�r sorgen, dass dieses fehlende Wissen in die Dokumentation aufgenommen wird.

Diese Praktik l�sst sich recht gut mit \FPZ\ anwenden. Die Listen mit den Reflektionsaussagen zu einzelnen Wissensgebieten lassen sich einfach als Textdateien anlegen und in einem Ordner speichern. In unserem Beispiel gibt es einen Ordner \emph{.Wissensgebiete}, der wiederum weitere Ordner f�r bestimmte Wissensgebiete enth�lt, hier \emph{Wissen Flowprotocol2} und \emph{Wissen Softwareentwicklung}. Durch die Markierung des Ordners \emph{.Wissensgebiete} mit einem Punkt am Anfang, wird dieser nicht als Skript-Unterordner im Men�baum von \FPZ\ aufgelistet.

Im Ordner \emph{Wissen f�r FlowProtocol2} sind die beiden Textdateien \emph{Basisbefehle.txt} und \emph{Programmierung.txt} hinterlegt, die wie die Skripte auch, mit einem einfachen Texteditor erstellt wurden. Diese haben folgenden Inhalt:

\Beispiel{Basisbefehle.txt}\label{BspBasisbefehle}
\begin{verbatim}
[Ausgabe]
Ich wei�, wie man eine Ausgabe erzeugt.
Ich wei�, wie man eine zweistufige Ausgabe erzeugt.
Ich wei�, wie man die Ausgabe in Abschnitte gliedert.
Ich wei�, wie man die Nummerierung der Aufz�hlung steuert.
Ich wei�, wie man einen Codeblock erzeugt.
Ich wei�, wie man einen Link auf eine Internetseite erzeugt.

[Eingabe]
Ich wei�, wie man eine Auswahlabfrage erstellt.
Ich wei�, wie man eine Texteingabe abfragt.
Ich wei�, wie man einen eingegebenen Text wieder ausgibt.
Ich wei�, wie man Auswahlabfragen ineinander verschachtelt.
\end{verbatim}

\Beispiel{Programmierung.txt}\label{BspProgrammierung}
\begin{verbatim}
[Strukturen]
Ich wei�, wie eine If-Abfrage aufgebaut ist. 
Ich wei�, wie eine Bedingung aufgebaut ist.
Ich wei�, wie man eine Von-1-bis-10-Schleife implementiert.
Ich wei�, wie man eine Liste von Elementen durchl�uft.
Ich wei�, wie man eine Funktion definiert und aufruft.
Ich wei�, wie man eine Sprungmarke definiert und anspringt.

[Verarbeitung]
Ich wei�, wie man Variablen setzt und verwendet.
Ich wei�, wie man Texte ersetzen kann.
Ich wei�, wie man regul�re Ausdr�cke verwenden kann.
Ich wei�, wie man Berechnungen durchf�hren kann.
\end{verbatim}

Beide Dateien haben das Dateilisten-Format\index{Dateilisten}, d.h.\ Leerzeilen und mit \verb|\\| beginnende Kommentarzeilen werden ignoriert und Zeilen, die mit eckigen Klammern beginnen und enden werden als Abschnitts�berschriften f�r die darauffolgenden Zeilen interpretiert.

Das Skript, mit dem man auf diese Daten zugreift sieht wie folgt aus:

\Beispiel{Wissensreflektion}\label{BspWissensreflektion}
\begin{verbatim}
~ListDirectories wdir = .Wissensgebiete; Pattern=Wissen *
~DynamicOptionGroup WGeb:wdir;W�hle das Wissensgebiet
~Execute
~ListFiles wdats = .Wissensgebiete|$wdir($WGeb); Pattern=*.txt
~Set idx=0
~ForEach wdat in wdats
    ~AddTo idx+=1
    ~Replace wpak($idx)=$wdat|.txt->
~Loop
~If $idx==0
    @Hinweis >>_ Im Ordner "Wissensgebiete\$wdir($WGeb)"
    __ sind keine Dateien vorhanden. Die Ausf�hrung
    __ wird abgebrochen.
    ~End
~DynamicOptionGroup WPak:wpak;W�hle das Wissenspaket
~Execute
~Set WPfad=.Wissensgebiete|$wdir($WGeb)|$wdats($WPak)
?A: Was m�chtest du tun?
    #sa: Alle Reflektionsaussagen durchgehen
        ~ForEachLine ref in $WPfad; IndexVar=idx; SectionVar=Ab
            ~SetInputSection $Ab
            ?R$idx: $ref
            #j: Ja
                @Was ich wei� >> $ref
            #n: Nein
                @NN >> $ref
            ~Calculate mod4 = $idx % 4
            ~If $mod4==0
                ~Execute
        ~Loop
        ~MoveSection  NN->Was ich nicht wei�
    #s4: Drei Reflektionsaussagen als Stichprobe
        ~ForEachLine ref in $WPfad; Take=3; IndexVar=idx
            ?R$idx: $ref
            #j: Ja
            #n: Nein
                ~Replace refmod=$ref|Ich wei�, ->...
                @Heute sollte ich rausfinden... >> $refmod
        ~Loop
    #db: Weitere Reflektionsaussagen erfassen
        @Anleitung >> �ffne die Wissensdatei
            >|$CurrentScriptPath\Wissensgebiete\$wdir($WGeb)
            __\$wdats($WPak)
        >> F�ge neue Reflektionsaussagen ein
            > Diese sollten mit "Ich wei�, " beginnen
            > ... und nach M�glichkeit nicht die Antwort enthalten
\end{verbatim}

Es beginnt mit einer Abfrage aller Verzeichnisse im Ordner \emph{.Wissensgebiete} die mit \emph{Wissen} beginnen. Dazu wird der Befehl \fcmd{ListDirectories} verwendet, der das Ergebnis in einem Feld abspeichert.

Durch den Befehl \fcmd{DynamicOptionGroup} wird diese Auf\-lis\-tung als Auswahl einer Auswahlabfrage angeboten, sodass der Anwender ein Wissensgebiet und damit gleichzeitig einen Ordner ausw�hlen kann.

Im n�chsten Schritt werden mit dem Befehl \fcmd{ListFiles} die Dateien innerhalb des gew�hlten Ordners abgefragt, die auf \emph{.txt} enden. �hnlich wie oben wird hier eine Auswahlabfrage generiert, zuvor wird jedoch in einer Schleife eine Auf\-lis\-tung ohne die txt-Endungen erzeugt. Wenn in dem gew�hlten Ordner keine Dateien vorhanden sind, wird die Ausf�hrung mit einer Hinweismeldung abgebrochen.

Nach Auswahl einer Wissenspaket-Datei hat der Anwender nun drei M�glichkeiten, diese zu verwenden. Er kann \emph{Alle Reflektionsaustragen durchgehen}, wobei diese mit den Antwortm�glichkeiten Ja und Nein abgefragt werden. Am Ende werden dann die beiden Listen mit den �berschriften \emph{Was ich wei�} und \emph{Was ich nicht wei�} ausgegeben. Diese Auf\-lis\-tung wird mit dem schon in Abschnitt~\ref{RegUndDaten} verwendeten Befehl \fcmd{ForEachLine} erzeugt, der jeder Zeile der ausgew�hlten Datei durchl�uft. Die Zeile selbst wird dabei in der \verb|ref|-Variablen zur�ckgegeben. Zus�tzlich k�nnen Variablen f�r den Zeilenz�hler �ber das \cidxfrag{IndexVar}-, und f�r die Abschnitts�berschrift �ber das \cidxfrag{SectionVar}-Argument angegeben werden. Bei jedem Schleifendurchlauf wird �berpr�ft ob der Z�hler durch vier teilbar ist und wenn ja, der \fcmd{Execute}-Befehl ausgef�hrt, sodass die Reflektionsaussagen in Viererpaketen abgefragt werden.

Bei \emph{Drei Reflektionsaussagen als Stichprobe} werden zuf�llig drei Zeilen aus der Datei ausgew�hlt. Die Zahl drei wird mit dem \cidxfrag{Take}-Argument angegeben. Der Seed-Wert\index{Seed-Wert} des Zufallsgenerators\index{Zufallsgenerator} wird dabei als Parameter \verb|_rseed| in der URL notiert, so dass beim erneuten Aufrufen oder einer Aktualisierung der Seite wieder die genau gleiche Zufallsfolge entsteht. Hier werden nur die mit Nein beantworteten Punkte unter der �berschrift \emph{Heute sollte ich rausfinden\dots} in leicht angepasster Form aufgelistet.

Bei \emph{Weitere Reflektionsaussagen erfassen} wird eine Anleitung ausgegeben, die dabei unterst�tzt, die richtige Datei zu finden und dort weitere Aussagen einzuf�gen. Hierbei wird f�r die Zusammensetzung des Pfades die Systemvariable \cidxvar{CurrentScriptPath} verwendet, die den Pfad des aktuell ausgef�hrten Skriptes ausgibt.

Bei der Anwendung dieser Mittel gibt es verschiedene M�glichkeiten. Die Zielsetzung sollte in jedem Fall darin bestehen, dass am Ende alle Mitarbeiter die Dinge wissen, die Sie wissen m�ssen, um ihre Arbeit gut zu machen, und dabei sollte auch immer davon ausgegangen werden, dass die Mitarbeiter die Bereitschaft dazu mitbringen.

Zu sehen, wie viel Fachwissen f�r die eigene Arbeit notwendig ist, unterstreicht wie anspruchsvoll diese ist, und sollte daher ein gutes Gef�hl vermitteln. Besonders dann, wenn ein Gro�teil dieses Wissens �ber Jahre hinweg aus dem Team heraus aufgebaut wurde, ist die Sichtweise darauf eine andere, als wenn die Inhalte ausschlie�lich aus externen Quellen stammen. Der Entwickler einer Framework-Komponente hat letztendlich das gleiche Interesse daran, dass sein Kollege diese richtig verwendet, wie der Datenschutzbeauftragte daran hat, dass die Datenschutz-Richtlinien allen bekannt sind und eingehalten werden.

Das Durchgehen der Reflektionsaussagen erm�glicht es jedem einzelnen, die wei�en Flecken auf seiner Wissenslandkarte zu identifizieren und eventuell vorhandene L�cken zu schlie�en. Daf�r muss er diese L�cken auch nicht gegen�ber einem anderen offenlegen, denn \FPZ\ arbeitet diskret und gibt keine Informationen weiter. Zu wissen, was man nicht wei�, ist die Grundvoraussetzung daf�r, um fehlendes Wissen zielgerichtet zu erg�nzen. Man kann damit die ben�tigte Information recherchieren oder direkt bei einem Kollegen nachfragen. Durch die Reflektionsaussage hat man sogar eine kompakte Beschreibung, mit der man bei seiner Anfrage auf das Wissen Bezug nehmen kann..

Generell kann man davon ausgehen, dass es etliche Punkte gibt, die noch von mehreren Mitarbeitern nicht gewusst werden, oder bei denen es aus anderen Gr�nden sinnvoll ist, sie nochmal in gemeinsamer Runde zu besprechen, z.B.\ um auf den gleichen Stand zu kommen. Gleichzeitig unterliegt auch das Wissen selbst einem permanenten Wandel, es kommen neue Punkte dazu und bestehendes Wissen wird eventuell nicht mehr ben�tigt. Von daher ist auch eine regelm��ige redaktionelle �berarbeitung notwendig und es macht Sinn beides zusammenzulegen.

Unter dem Titel Wissensaktualisierung\index{Wissensaktualisierung} plant man mindestens einmal im Jahr einen Termin ein, bei dem alle Kollegen die M�glichkeit haben, sowohl ihr eigenes, als auch das gemeinsame Wissen aufzufrischen und zu aktualisieren. Idealerweise geht jeder daf�r im Vorfeld die vorhandenen Listen durch und notiert die Punkte, die er im Meeting angesprochen haben m�chte, sei es um sich selbst oder seine Kollegen auf den aktuellen Stand zu bringen, oder �ber die Herausnahme oder �berarbeitung des Reflektionsaussage zu diskutieren. 

F�r die Sammlung dieser Punkte kann eine eigene Variante im Wissensreflektion-Skript eingebaut werden. Die von den verschiedenen Mitarbeitern gesammelten Punkte werden dann vom Moderator des Meetings zusammengetragen und je nach H�ufigkeit zu einer Agenda umgesetzt. Je nach Menge der zu besprechenden Reflektionsaussagen ist es sinnvoll, diese in Vorbereitung des Meetings noch mit Antworten oder Verweisen auf eventuelle Informationsquellen anzureichern, sodass man im Meeting selbst keine Zeit mehr daf�r aufwenden muss. Die Diskussion �ber die zu wissenden Dinge ist ebenso wertvoll wie der Wissensaustausch selbst, und so wird jeder am Ende des Tages viel gewonnen haben.

Die im Beispiel gezeigte Variante mit der Stichprobe\index{Stichprobe} aus drei Reflektionsaussagen l�sst sich gut daf�r einsetzen, sich jeden Morgen oder zweimal die Woche eine solche Auswahl zu generieren, und so permanent sein eigenes Wissen auf die Probe zu stellen. Die dabei gefundenen L�cken lassen sich dann ebenso beil�ufig im Arbeitsalltag schlie�en.

Diese Form der Stichproben-Abfrage kann mit einer etwas gr��eren Anzahl auch daf�r verwendet werden, eine Art von Pr�fungsbogen zu erstellen, der dann f�r eine Kontrolle des Wissens durch eine dritte Person verwendet wird. Hier muss man sich jedoch genau fragen, ob und unter welcher Voraussetzungen man einen derartigen Kontrollmechanismus tats�chlich einsetzen m�chte.

Umgekehrt gibt es auch gesetzliche Nachweispflichten, die man in Bezug auf seine Mitarbeiter wahrnehmen muss, etwa dass diese einen F�hrerschein haben, wenn sie mit Firmenfahrzeugen unterwegs sind. Das Zusenden der mit Ja beantworteten Reflektionsaussagen ist eine einfache M�glichkeit, das Vorhandensein von Wissen gegen�ber seinen Vorgesetzten konkret zu best�tigen. Diese Vorgehensweise ist dahingehend stressfrei, da sich der Mitarbeiter ja die f�r sich notwendige Zeit nehmen kann, um eine eventuell vorgegebene Mindestquote zu erreichen. Erg�nzend dazu gibt es ja auch noch den gemeinsamen Austausch und den Wissensaktualisierungstermin.

Sowohl die Berechnung der Quote, als auch das optionale Hinterlegen eventueller Antworten und Informations-Verweise direkt bei den Reflektionsaussagen l�sst sich sehr einfach im Skript umsetzen. Ersteres ist eine einfache Berechnung und f�r den zweiten Punkt kann man f�r das Hinterlegen von Antwort-Informationen eine bestimmte Syntax vorgeben, die man mit Hilfe eines regul�ren Ausdrucks verarbeiten kann.

\subsection{Personalisierung}\label{Personalisierung}
Bei vielen Arten von interaktiven Anleitungen wird man zwangsl�ufig auf Elemente verweisen wollen, die lokal von der jeweiligen Umgebung des Mitarbeiters abh�ngig sind. Dazu geh�ren in erster Linie Datei- und Ordnerpfade, mit denen sich  etwa parametrisierte Kommandozeilenaufrufe vorbereiten lassen, so dass der Anwender diese �ber eine Konsole oder die  Tastenkombination \emph{Win-R}\index{Win-R} ausf�hren kann. Dazu kommen Benutzerk�rzel oder andere Kennzeichen, die sich z.B.\ in erzeugten Programmcode einarbeiten lassen. 

Die Pfadstruktur k�nnte man vereinheitlichen, was allerdings m�hsam und einschr�nkend ist, und die ansonsten ben�tigten Werte kann man abfragen, allerdings ist es gleicherma�en f�r Ersteller und Anwender mit viel  Aufwand verbunden, wenn dies in jedem Skript geschieht.

Die L�sung besteht darin, eine einzelne Funktionsdatei zu erstellen, die bei Auswahl eines Benutzers eine definierte Menge von Variablen benutzerbezogen setzt, so dass diese im ganzen Skript verf�gbar sind. Diese k�nnte so aussehen:

\Beispiel{Benutzervariablen.fps}
\begin{verbatim}
?BK: Welche Benutzer-/Umgebungsvariablen sollen genutzt werden?
    #AB: Anton Barium
        ~Set BenVarKuerzel=AB
        ~Set BenVarPfadAufgabenmappen=H:\Aufgabenmappen
        ~Set BenVarPfadTools=C:\Work\Tools
    #CD: Christian Deuterium
    #EF: Ernst Flour
\end{verbatim}

Im Beispiel sind exemplarisch die Variablen f�r den Benutzer AB ausformuliert, und in einer im Echtbetrieb eingesetzten Variante w�rde es vermutlich noch viel mehr benutzerbezogene Variablen geben. 

Mit den Hilfsmitteln aus Abschnitt~\ref{Metaskripting} k�nnte man auch dieses Skript ohne M�he selbst-erweiterbar machen, so dass ein Mitarbeiter, der erstmalig mit dieser Abfrage konfrontiert wird, bei den Eingaben angeleitet und unterst�tzt wird.

Zu beachten ist auch, dass alle Variablen mit einem einheitlichen Pr�fix benannt wurden, sodass versehentliche Namenskonflikte mit den im Skript selbst gesetzten Variablen ausgeschlossen werden k�nnen.

Ein Skript, das auf solchen Variablen aufbaut, k�nnte wie folgt aussehen:

\Beispiel{Patchskript}\label{BspPatchskript}
\begin{verbatim}
~Include Benutzervariablen.fps
~Input AKennung: Gib die Aufgabenkennung ein:
~Input PatchCode: Gib den Korrekturcode ein:
~SetDateTime Datum=yyyy-MM-dd
~Set DatBasis=$BenVarPfadAufgabenmappen\$Datum_$AKennung
~XmlEncode xPatchCode=$PatchCode
@Anleitung >> Erstelle eine neue Datei im Editor.
>> Gib den folgenden Code ein:
    >|<?xml version="1.0"?>
    >|<PATCHSCRIPT>
    >|  <NAME>$AKennung</NAME>
    >|  <CODE>$xPatchCode</CODE>
    >|  <CREATOR>$BenVarKuerzel</CREATOR>
    >|  <CREATIONDATE>$Datum</CREATIONDATE>
    >|</PATCHSCRIPT>
>> Speichere die Datei ab unter
    >|$DatBasis\$AKennung.xml
>> Rufe den Patch-Converter f�r die Datei auf:
    >|$BenVarPfadTools\PatchConverter\Bin\PatchConverter.exe
        __ "$DatBasis\$AKennung.xml"
>>  Das fertige Patchscript
    ~AddCode  $AKennung.psc
    ~AddText  liegt unter
    >|$DatBasis
\end{verbatim}

In diesem Beispiel wird die Erstellung eines einfachen Patchskriptes angeleitet, das in einer nach Standardkonventionen benannten Aufgabenmappe als XML-Datei erstellt und mit Hilfe einer internen Hilfsanwendung in ein anderes Format konvertiert wird.

Die Anleitung kombiniert daf�r verschiedene innerhalb der Anleitung abgefragte und berechnete Eigenschaften mit den benutzerabh�ngigen Variablen. Dazu wird die Funktionsdatei \emph{Benutzervariablen.fps} gleich zu Beginn eingelesen, womit die Auswahlabfrage nach den Benutzervariablen als Teil des Skriptes erscheint.

Ein Anwender kann den Link auf das Skript in den Lesezeichen seines Browsers nun auch so abspeichern, dass diese erste Abfrage dort schon beantwortet ist. Die URL daf�r bekommt er, indem er nur die Auswahlabfrage nach den Benutzervariablen beantwortet und auf \emph{Weiter} klickt. In unserem Fall wird dabei der URL-Parameter \verb|BK=AB| angef�gt. Damit ist beim Aufruf gar nicht mehr erkennbar, dass das Skript benutzerabh�ngige Anteile hat.

Mit dem Befehl \fcmd{XmlEncode} wird im Beispiel der eingegebene Patchcode in die XML-Codierung\index{XML-Codierung} umgewandelt, so dass die f�r die Verwendung innerhalb eines XML-Tags kritischen Zeichen wie \verb|<|, \verb|>| und \verb|&|,  sowie die Anf�hrungszeichen maskiert werden. Bei den anderen Elementen l�sst sich das Vorhandensein dieser Zeichen semantisch ausschlie�en.

\section{Befehlsreferenz}\label{Befehlsreferenz}

\subsection{Kernelemente}
\begin{grundelement}\label{defBedingungen}
Eine \emph{Bedingung}\index{Bedingungen} ist ein Ausdruck, der zu \emph{wahr} oder \emph{falsch} ausgewertet werden kann. Bedingungen werden bei verschiedenen Befehlen verwendet, z.B.\ in Do-While-Schleifen oder  bei If-Bedingungen.\par
Eine Bedingung muss in \FPZ\ in der disjunktiven Normalform angegeben werden, also als Oder-Verkn�pfung (\verb=||=) von Und-Verkn�pfungen (\verb|&&|), wobei keine Klammerung notwendig ist.

Als Literale sind die Konstanten \verb|1| und \cidxfrag{true} (wahr\index{wahr}), sowie \verb|0| und \cidxfrag{false} (falsch\index{falsch}) verwendbar, sowie die folgenden Vergleichsoperatoren\index{Vergleichsoperatoren}:
\begin{center}
\begin{tabular}{|l|c|l|l|}\hline
Ausdruck& Formel & Sprachform & Typenbeschr�nkung\\\hline\hline
\verb|$s==$t| & $s=t$ & $s$ ist gleich $t$ & f�r Zeichenketten $s$ und $t$\\
\verb|$s!=$t|  & $s\neq t$ & $s$ ist ungleich $t$ & f�r Zeichenketten $s$ und $t$\\
\verb|$x<>$y| & $x\neq y$ & $x$ ist ungleich $y$ & f�r Zahlen $x$ und $y$\\
\verb|$x<$y| & $x<y$ & $x$ ist kleiner $y$ & f�r Zahlen $x$ und $y$\\
\verb|$x<=$y| & $x\le y$ & $x$ ist kleiner gleich $y$ & f�r Zahlen $x$ und $y$\\
\verb|$x>$y| & $x>y$ & $x$ ist gr��er $y$ & f�r Zahlen $x$ und $y$\\
\verb|$x>=$y| & $x\ge y$ & $x$ ist gr��er gleich $y$ & f�r Zahlen $x$ und $y$\\
\verb|$s~$t| & $s\supseteq t$ & $s$ enth�lt $t$ & f�r Zeichenketten $s$ und $t$\\
\verb|$s!~$t| & $s\nsupseteq t$ & $s$ enth�lt $t$ nicht & f�r Zeichenketten $s$ und $t$\\
\verb|?$v| && $v$ ist gesetzt & f�r eine Variable $v$\\
\verb|!?$v| && $v$ ist nicht gesetzt) & f�r eine Variable $v$\\\hline
\end{tabular}
\end{center}
\end{grundelement}

\begin{grundelement}\label{defDateipfad}
Ein \emph{Dateipfad}\index{Dateipfad}, bzw.\  \emph{Ordnerpfad}\index{Ordnerpfad} ist eine Zeichenkette, die einen Pfad auf eine Datei, bzw.\ einen Ordner angibt. Als Trennzeichen in den Pfadangaben kann das Trennzeichen des Betriebssystems verwendet werden, auf dem \FPZ\ l�uft, also \verb|\| f�r Windows\index{Windows} und \verb|/| f�r Linux\index{Linux}, es kann aber auch \verb=|= verwendet werden, das unabh�ngig vom Betriebssystems richtig aufgel�st wird.\par
Der Zugriff von \FPZ\ ist dabei beschr�nkt auf Dateien und Ordner, die im Skripte-Verzeichnis liegen. Die Angabe \verb|.| verweist auf den Ordner, in dem das gerade ausgef�hrte Skript liegt. Angaben, die mit \verb=.|= (bzw.\ \verb|.\| oder \verb|./|) beginnen, werden ausgehend vom Skripte-Verzeichnis interpretiert und alle anderen werden wieder ausgehend von dem Ordner interpretiert, in dem das gerade ausgef�hrte Skript liegt
\end{grundelement}

\begin{grundelement}\label{defLinks}
Ein \emph{Link}\index{Link} ist ein im Browser angezeigter Text, der optisch meist blau und unterstrichen hervorgehoben ist, und der bei Anklicken eine im Link hinterlegte Adresse (URL) in einem Tab oder einen neuen Fenster �ffnet.\\
Links k�nnen in \FPZ\ sowohl in der Ausgabe, als auch in Hilfezeilen erzeugt werden, wobei stets eine URL und ein Anzeigetext angegeben werden k�nnen.\par
Um die Gefahr abzuwenden, dass �ber Links nicht erkennbar sch�dliche URL-Verweise in Skripten eingebaut werden, kann man festlegen, welche URLs ausschlie�lich durch den Anzeigetext angezeigt werden sollen. Ein Link wird im Browser ausschlie�lich mit dem Anzeigetext dargestellt, wenn\dots
\begin{itemize}
\item der Anfang der URL mit einem Eintrag der in der Konfiguration hinterlegten Whitelist (Datei appsettings.json, Eigenschaft \verb|LinkWhitelist|) �bereinstimmt oder
\item die Whitelist leer ist oder
\item auf ein anderes \FPZb-Skript verwiesen wird oder
\item der Anzeigetext den Link als Text enth�lt (ggf. auch ohne "https://").
\end{itemize}
Anderenfalls wird die URL dem Anzeigetext in Klammern nachgestellt.\par
Die in \FPZ\ angezeigten Links werden standardm��ig in einem neuen Tab ge�ffnet.
\end{grundelement}

\begin{grundelement}\label{defListen}
Eine \emph{Liste}\index{Listen} ist in \FPZ\ eine Folge von fortlaufend indizierten Variablen, z.B.\ \verb|V(1)|, \verb|V(2)|, \verb|V(3)|,\dots , die mit dem Index $1$ beginnt und mit dem letzten Index endet, f�r den eine entsprechende Variable noch gesetzt ist. Listen werden durch verschiedene Befehle erzeugt und k�nnen selbst auch als Argumente an Befehle �bergeben werden. Hierbei wird dann nur der Grundname der Variable (hier also \verb|V|) angegeben.
\end{grundelement}

\subsection{Umgebungsvariablen}

\begin{description}
\item[\cidxvar{BaseKey}] gibt innerhalb von Funktionen oder mit dem  \verb|~Include|-Befehl aufgerufene Funktionsdateien den Wert zur�ck, der mit dem gleichnamigen Argument beim Aufruf �bergeben wurde. Diese Variable gibt auch bei rekursiv verschachtelten Selbstaufrufen immer den zum jeweiligen Aufruf passenden Wert zur�ckgibt. Mit ihr k�nnen Argumente und Variablen bei rekursiven Aufrufen derselben Funktion voneinander abgetrennt werden.

\item[\cidxvar{BaseURL}] gibt die URL f�r das aktuell ausgef�hrte Skript ohne Parameter zur�ck.

\item[\cidxvar{Chr(...)}] gibt das Zeichen mit dem angegebenen ANSI-Code zur�ck

\item[\cidxvar{CRLF}] gibt die Zeichenkombination \emph{Carriage Return + Line Feed} zur�ck.

\item[\cidxvar{LineNumber}] gibt die Zeilennummer der aktuellen Zeile zur�ck. Mittels Zusatz $-n$ wird direkt der um $n$ verringerte Wert zur�ckgegeben, z.B.\ \verb|$LineNumber-3|.

\item[\cidxvar{NewGuid}] erzeugt eine Guid, z.B.\ 8bceae6b-aa8b-497e-aa3a-17a6634b7215.

\item[\cidxvar{ResultURL}] gibt die URL f�r das aktuell ausgef�hrte Skript mit den aktuell gesetzten Parametern zur�ck.

\item[\cidxvar{ScriptFilePath}] gibt den Dateipfad f�r das aktuell ausgef�hrte Skript aus Sicht der Anwendung zur�ck.

\item[\cidxvar{ScriptPath}] gibt den Dateipfad f�r das Skripte-Verzeichnis\index{Skripte-Verzeichnis} aus Sicht der Anwendung zur�ck.

\item[\cidxvar{CurrentScriptPath}] gibt den Dateipfad des Ordners f�r das aktuell ausgef�hrte Skript aus Sicht der Anwendung zur�ck.
\end{description}

\subsection{Befehle}\label{Befehle}
\begin{description}

\item[\cidxfrag{>\mbox{}>}]$[$\synvar{Format}$]$ \synvar{Text}\\
Gibt einen Text als Aufz�hlung der ersten Ebene aus. Als Format kann zwischen \verb|#| f�r nummerierte Aufz�hlung, \verb|*| f�r nicht nummerierte Aufz�hlung, \verb|_| f�r Flie�text und \verb=|= f�r Codezeile gew�hlt werden. Standardm��ig wird die nummerierte Aufz�hlung verwendet.\par
Optional kann durch Voranstellen von \verb|@|\synvar{Abschnitt} ein Abschnitt ausgew�hlt werden.

\item[\cidxfrag{>}]$[$\synvar{Format}$]$ \synvar{Text}\\
Gibt einen Text als Aufz�hlung der zweiten Ebene aus. Als Format kann zwischen \verb|#| f�r nummerierte Aufz�hlung, \verb|*| f�r nicht nummerierte Aufz�hlung, \verb|_| f�r Flie�text und \verb=|= f�r Codezeile gew�hlt werden. Standardm��ig wird die nicht nummerierte Aufz�hlung verwendet.\par
Optional kann durch Voranstellen von \verb|@|\synvar{Abschnitt} ein Abschnitt ausgew�hlt werden.

\item[\cidxfrag{>.}]$[$\synvar{Format}$]$ \synvar{Text}\\
Gibt einen Text als Aufz�hlung der dritten Ebene aus. Als Format kann zwischen \verb|#| f�r nummerierte Aufz�hlung, \verb|*| f�r nicht nummerierte Aufz�hlung, \verb|_| f�r Flie�text und \verb=|= f�r Codezeile gew�hlt werden. Standardm��ig wird die nicht nummerierte Aufz�hlung verwendet.\par
Optional kann durch Voranstellen von \verb|@|\synvar{Abschnitt} ein Abschnitt ausgew�hlt werden.

\item[\cidxfrag{?}] $[$\synvar{Schl�ssel}$]$ \verb|:| \synvar{Text}\\
Erzeugt eine Auswahlabfrage mit mehreren Optionen, von denen genau eine ausgew�hlt werden kann. Der Anzeigetext wird oberhalb der Auswahlliste angezeigt.
Der Schl�ssel wird verwendet, um die Eingabe in der URL zu speichern, und innerhalb des Skriptes abzufragen. Er muss eindeutig sein, bzw. wird bei Wiederholung als schon gesetzt betrachtet. Innerhalb einer Datei k�nnen die Schl�ssel durch Anh�ngen des \verb|'|-Zeichens automatisch durchnummeriert werden.\par
Die Optionen werden darunter mit dem \verb|#|-Befehl angegeben.
�ber die beiden Variablen \synvar{Schl�ssel}\verb|_OptionGroupPromt| und \synvar{Schl�ssel}\verb|_SelectedOptionText| lassen sich der Text der Frage und der ausgew�hlten Option abrufen.

\item[\cidxfrag{\#}] $[$\synvar{Optionsschl�ssel}$]$ \verb|:| \synvar{Anzeigetext}\\
F�gt eine Option zur �bergeordneten Auswahlabfrage hinzu. Wird der Optionsschl�ssel weggelassen, wird dieser automatisch durchnummeriert.
F�r die jeweils gew�hlte Option werden die darunter einger�ckten Befehle ausgef�hrt, f�r die anderen nicht.\\

\item[\fcmd{AddCode}] \synvar{Text}\\
H�ngt einen Text in der f�r Programmcode typischen Schriftart an die letzte Ausgabezeile an. Um dazwischen ein Leerzeichen zu erzeugen, m�ssen zwischen Befehlsname und Argument zwei Leerzeichen angegeben werden.\\
Siehe Abschnitt~\ref{LinksInline}, verwendet in Beispiel~\ref{BspLinksUndInlineCode}, \ref{Archivklasse2} und \ref{BspPatchskript}.

\item[\fcmd{AddHelpLine}] \synvar{Text}\\
F�gt eine Hilfezeile f�r die �bergeordnete Auswahlabfrage oder Texteingabe mit einem Text hinzu.\\
Siehe Abschnitt~\ref{Hilfezeilen}, verwendet in Beispiel~\ref{BspHilfezeilen}, \ref{Archivklasse2} und \ref{BspMetaErweiterung}.

\item[\fcmd{AddHelpLink}] \synvar{URL} \verb=|= \synvar{Anzeigetext}\\
F�gt einen Link mit der entsprechenden URL und dem Anzeigetext in der aktuellen Hilfezeile ein (siehe auch Kernelement~\ref{defLinks}).\\
Siehe Abschnitt~\ref{Hilfezeilen}, verwendet in Beispiel~\ref{BspHilfezeilen}.

\item[\fcmd{AddHelpText}] \synvar{Text}\\
F�gt den Text \synvar{Text} in der aktuellen Hilfezeile ein. Um dazwischen ein Leerzeichen zu erzeugen, m�ssen zwischen Befehlsname und Argument zwei Leerzeichen angegeben werden.\\
Siehe Abschnitt~\ref{Hilfezeilen}, verwendet in Beispiel~\ref{BspHilfezeilen}.

\item[\fcmd{AddLink}] \synvar{URL} \verb=|= \synvar{Anzeigetext}\\
F�gt einen Link mit der entsprechenden URL und dem Anzeigetext in der aktuellen Ausgabezeile ein (siehe auch Kernelement~\ref{defLinks}).\\
Siehe Abschnitt~\ref{Befehle}, verwendet in Beispiel~\ref{BspLinksUndInlineCode}, \ref{BspSuchlink}, \ref{BspJourFixePlaner} und \ref{BspLizenzabfrage}.

\item[\fcmd{AddNewKey}] \synvar{Schl�ssel} \verb|=| \synvar{Wert}\\
Speichert einen Wert unter einem Schl�ssel in der URL, wenn der Schl�ssel dort noch nicht enthalten ist.\\
Siehe Abschnitt~\ref{Zeitmessung}, verwendet in Beispiel~\ref{BspZeitmesser}.

\item[\fcmd{AddText}] \synvar{Text}\\
F�gt einen Text in der aktuellen Ausgabezeile ein. Um dazwischen ein Leerzeichen zu erzeugen, m�ssen zwischen Befehlsname und Argument zwei Leerzeichen angegeben werden.\\
Siehe Abschnitt~\ref{LinksInline}, verwendet in Beispiel~\ref{BspLinksUndInlineCode}, \ref{BspArchivklasse3}, \ref{BspListenfunktion}, \ref{BspJourFixePlaner} und \ref{BspPatchskript}.

\item[\fcmd{AddTo}] \synvar{Variable} \verb|+=| \synvar{Zahlenwert}\\
Addiert eine Zahlenwert zu einer Variablen dazu. Hat die Variable noch keinen Wert, so wird der Ausgangswert $0$ impliziert.\\
Siehe Abschnitte~\ref{BewSysteme} und \ref{Berechnungen}, verwendet in Beispiel~\ref{BspBewertungssystem}, \ref{BspBugPlaner}, \ref{BspPrimzahlen}, \ref{BspAbfrageschleife}, \ref{BspTuermeVonHanoi}, \ref{BspKombinator}, \ref{Passwortgenerator}, \ref{BspZeitmesser}, \ref{BspJourFixePlaner}, \ref{BspTerminplaner} und \ref{BspWissensreflektion}.

\item[\fcmd{Calculate}] \synvar{Variable} \verb|=| \synvar{Operand 1} \synvar{Operator} \synvar{Operand 2}\\
F�hrt eine einfache Rechnung aus und speichert das Ergebnis in einer Variablen. Als Operator stehen zur Auswahl \verb|+|, \verb|-|, \verb|*|, \verb|/| und \verb|%| (Modulo).\\
Siehe Abschnitt~\ref{Berechnungen}, verwendet in Beispiel~\ref{BspBerechnungen}, \ref{BspPrimzahlen}, \ref{BspZeitmesser}, \ref{BspTerminplaner} und \ref{BspWissensreflektion}.

\item[\fcmd{CalculateExpression}] \synvar{Variable} \verb|=| \synvar{Ausdruck}\\
Berechnet einen mathematischen Ausdruck, der aus den vier Grundrechenarten \verb|+|, \verb|-| , \verb|*| und \verb|/|, dem \verb|^|-Operator f�r die Potentenzierung, dem \verb|%|-Operator 
f�r Modulo-Rechnung, sowie den Operatoren \verb|sqrt|, \verb|sin|, \verb|cos|, \verb|tan|, \verb|exp| und \verb|ln| in einer beliebigen Klammerung besteht. Das Ergebnis wird einer Variablen zugewiesen. Es gilt die Punkt-vor-Strich-Rechnung\index{Punkt-vor-Strich-Rechnung}. Die Argumente der trigonometrischen Funktionen werden im Bogenma� interpretiert.\\
Siehe Abschnitt~\ref{Berechnungen}, verwendet in Beispiel~\ref{BspBerechnungen}, \ref{BspPrimzahlen}, \ref{BspZeitmesser} und \ref{BspTerminplaner}.

\item[\fcmd{CamelCase}] \synvar{Variable} \verb|=| \synvar{Text}\\
Wandelt einen Text in CamelCase-Schreibweise\index{CamelCase-Schreibweise} um und speichert das Ergebnis in einer Variablen.
Hierbei werden Whitespace, Umlaute und Sonderzeichen entfernt und durch Gro�schreibung des darauffolgenden Buchstaben gekennzeichnet.\\
Siehe Abschnitt~\ref{Entwicklungsanleitungen}, verwendet in Beispiel~\ref{Archivklasse2} und \ref{BspMetaErweiterung}.

\item[\fcmd{ClearVar}] \synvar{Variablenmuster}
L�scht alle Variablen aus den Speicher, die einem Variablenmuster entsprechen. Dies kann ein Variablenname sein oder der Anfang eines Variablennamens, gefolgt von \verb|*|, wobei im zweiten Fall alle Variablen gel�scht werden, die auf diese Weise beginnen. Die Ausf�hrungszeit f�r einzelne Befehle steigt deutlich an, wenn eine gro�e Anzahl an Variablen, also 2000 oder mehr, verwendet werden. Durch den Befehl k�nnen ganze Listen von Variablen nach Gebrauch wieder entfernt werden, so dass die Ausf�hrungszeit f�r die nachfolgenden Befehle wieder sinkt.

\item[\fcmd{DateAdd}] \synvar{Variable} \verb|=| \synvar{Ausgangsdatum} \verb=|= \synvar{Zahlenwert} \verb=|= \synvar{Intervall}\\
Addiert zu einem Ausgangsdatum eine Zeitspanne hinzu, die �ber einen ganzzahligen Zahlenwert und ein Intervall gegeben wird. Als Intervall sind folgende Zeichen zugelassen:
\verb|y| = Jahre, \verb|M| = Monate, \verb|w| = Wochen, \verb|d| = Tage, \verb|h| = Stunden, \verb|m| = Minuten\\
Das Ausgangsdatum muss im internen Format f�r Datumswerte\index{Datumswerte} angegebe werden (yyyy-MM-dd HH:mm:ss).\\
Siehe Abschnitt~\ref{MeetingVorbereitung}, verwendet in Beispiel~\ref{BspJourFixePlaner}.

\item[\fcmd{DateFormat}] \synvar{Variable} \verb|=| \synvar{Datumswert} \verb=|= \synvar{Format}\\
Wanelt einen Datumswert in ein bestimmten Format um und speichert das Ergebnis in einer Variablen. Als Format sind die in der Programmiersprache C\# m�glichen Formatzeichen zul�ssig. Der Datumswert muss im internen Format f�r Datumswerte\index{Datumswerte} angegebe werden (yyyy-MM-dd HH:mm:ss).\\
Siehe Abschnitt~\ref{MeetingVorbereitung}, verwendet in Beispiel~\ref{BspJourFixePlaner}.

\item[\fcmd{DateSet}] \synvar{Variable} \verb|=| \synvar{Wert} \verb=|= \synvar{Format}\\
Interpretiert einen Text als Datum in einem vorgegebenen Format und speichert das Ergebnis in einer Variablen. Als Format sind die in der Programmiersprache C\# m�glichen Formatzeichen zul�ssig. Der Ergebniswert wird im internen Format f�r Datumswerte\index{Datumswerte} gespeichert (yyyy-MM-dd HH:mm:ss).\\
Siehe Abschnitt~\ref{Befehle}, verwendet in Beispiel~\ref{BspJourFixePlaner}.

\item[\fcmd{DecryptText}] \synvar{Variable} \verb|=| \synvar{Text} \verb=|= \synvar{Schl�ssel} \\
Entschl�sselt einen mit \fcmd{EncryptText} verschl�sselten Text mit dem angegebenen Schl�ssel und speichert das Ergebnis in einer Variablen.

\item[\fcmd{DefineSub}] \synvar{Funktionsname}\\
Definiert eine Funktion mit dem angegebenen Funktionsnamen, die dann mit dem \fcmd{GoSub}-Befehl aufgerufen werden kann. Die Funktion muss mit \fcmd{Return} abgeschlossen werden.
%Siehe Abschnitte~\ref{}, \ref{} und \ref{}, verwendet in Beispiel~\ref{}, \ref{} und \ref{}.\\
Siehe Abschnitt~\ref{Funktionen}, verwendet in Beispiel~\ref{BspTuermeVonHanoi}, \ref{BspZeitmesser} und \ref{BspTerminplaner}.
 
\item[\fcmd{DoWhile}] \synvar{Bedingung}\\
F�hrt die nachfolgenden einger�ckten Befehle wiederholt aus, solange die Bedingung erf�llt ist. Die Schleife wird mit dem Befehl \fcmd{Loop} auf gleicher Ebene abgeschlossen. Siehe auch Kernelement~\ref{defBedingungen}.\\
Siehe Abschnitte~\ref{FormatierungDerEingabe} und \ref{Schleifen}, verwendet in Beispiel~\ref{BspBugPlaner}, \ref{BspPrimzahlen}, \ref{BspAbfrageschleife}, \ref{Passwortgenerator}, \ref{BspZeitmesser}, \ref{BspJourFixePlaner} und \ref{BspTerminplaner}.

\item[\fcmd{DynamicOptionGroup}] \synvar{Schl�ssel} \verb|:| \synvar{Optionenliste} \verb|;| \synvar{Fragestellung}\\
Erzeugt eine Auswahlabfrage mit einem Schl�ssel, einer Fragestellung und einer Menge von Optionen, von denen genau eine ausgew�hlt werden kann. Im Unterschied zum \verb|?|-Befehl werden die Optionen hier zur Laufzeit aus einer Liste ausgelesen, siehe auch Kernelement~\ref{defListen}.
Das mit dem Schl�ssel in der URL gespeicherte und als Variable abrufbare Ergebnis der Auswahlabfrage ist der Index der ausgew�hlten Option.
Der Schl�ssel der Auswahlabfrage kann wie beim \verb|?|-Befehl mit \verb|'| per Durchnummerierung im Skript festgelegt werden.
Ist die Optionenliste leer, wird der Befehl �bersprungen.\\
Siehe Abschnitt~\ref{Wissensdokumentation}, verwendet in Beispiel~\ref{BspWissensreflektion}.

\item[\fcmd{Else}]\mbox{}\\
Ist eine optionale Erg�nzung zu einer If-Abfrage und evtl.\ vorangehenden ElseIf-Abfragen. Der Befehl f�hrt die nachfolgend einger�ckten Befehle aus, wenn alle vorherigen Bedingungen nicht erf�llt waren. Pro If-Abfrage darf maximal eine Else-Abfrage angegeben werden.\\
Siehe Abschnitt~\ref{IfUndBedingungen}, verwendet in Beispiel~\ref{BspBewertungsschema}, \ref{Passwortgenerator}, \ref{BspZeitmesser}, \ref{BspTerminplaner} und \ref{BspMetaErweiterung}.

\item[\fcmd{ElseIf}] \synvar{Bedingung}\\
Pr�ft eine alternative Bedingung nach einer If-Abfrage und evtl.\ anderen vorangehenden ElseIf-Abfragen. Der Befehl f�hrt die nachfolgend einger�ckten Befehle aus, wenn die angegebene Bedingung erf�llt ist und alle Bedingungen der vorangegangenen If-, und ElseIf-Befehle nicht erf�llt waren. Es k�nnen beliebig viele ElseIf-Abfragen aneinandergereiht werden.\\
Siehe Abschnitt~\ref{IfUndBedingungen}, verwendet in Beispiel~\ref{BspBewertungsschema} und \ref{Passwortgenerator}.

\item[\fcmd{EncryptText}] \synvar{Variable} \verb|=| \synvar{Text} \verb=|= \synvar{Schl�ssel}\\
Verschl�sselt einen Text mit dem angegebenen Schl�ssel und speichert das Ergbnis in einer Variablen.

\item[\fcmd{End}]\mbox{}\\
Beendet die Skriptausf�hrung.\\
Siehe Abschnitt~\ref{Funktionen}, verwendet in Beispiel~\ref{BspTuermeVonHanoi} und \ref{BspWissensreflektion}.

\item[\fcmd{EndParagraph}]\mbox{}\\
Beendet den aktuellen Absatz im aktuellen Textblock der Ausgabe. Das Hinzuf�gen der n�chsten Zeile beginnt damit automatisch einen neuen Absatz mit einem entsprechenden Abstand.

\item[\fcmd{EvalExpression}] \synvar{Variable} $[$\verb|!|$]$ \verb|=| \synvar{Bedingung}\\
Wertet eine Bedingung aus und speichert das Ergebnis in einer Variablen als \cidxfrag{true} oder \cidxfrag{false}. Mit \verb|=| wird der Wert direkt zugewiesen, mit \verb|!=| wird der Wert bei der Zuweisung negiert. Die Ergebniswerte solcher Auswertungen k�nnen in anderen Ausdr�cken weiterverwendet werden. Siehe auch Kernelement~\ref{defBedingungen}.\\
Siehe Abschnitt~\ref{FormatierungDerEingabe}, verwendet in Beispiel~\ref{BspBugPlaner}.

\item[\fcmd{Execute}]\mbox{}\\
F�hrt alle bis zu dieser Stelle durchlaufenen Eingabebefehle aus, sofern f�r diese noch kein Wert gesetzt ist. Erst wenn alle diese Eingaben vorhanden sind, werden die nachfolgenden Befehle ausgef�hrt.
Mit dem Befehl kann die Menge der gleichzeitig abgefragten Eingaben reduziert werden, und es kann sichergestellt werden, dass die nachfolgenden Befehle erst dann ausgef�hrt werden, wenn die von diesen ben�tigten Eingaben auch vorhanden sind.\\
Siehe Abschnitte~\ref{Texteingaben}, \ref{FormatierungDerEingabe} und \ref{Wissensdokumentation}, verwendet in Beispiel~\ref{BspBugPlaner}, \ref{BspTextersetzung}, \ref{BspAbfrageschleife}, \ref{BspBewertungsschema}, \ref{BspZeitmesser}, \ref{BspTerminplaner} und \ref{BspWissensreflektion}.

\item[\fcmd{ExitLoop}]\mbox{}\\
Verl�sst eine Schleife und setzt die Ausf�hrung nach dem Loop-Befehl fort.\\
Siehe Abschnitte~\ref{Schleifen} und \ref{MeetingVorbereitung}, verwendet in Beispiel~\ref{BspPrimzahlen} und \ref{BspJourFixePlaner}.

\item[\fcmd{FileExists}] \synvar{Variable} \verb|=| \synvar{Dateipfad}\\
Pr�ft, ob eine entsprechende Datei existiert und speichert das Ergebnis als Wahrheitswert in einer Variablen. Siehe auch Kernelement~\ref{defDateipfad}.

\item[\fcmd{ForEach}] \synvar{Variable} \verb|in| \synvar{Liste}\\
F�hrt die nachfolgenden einger�ckten Befehle bis zum dazugeh�renden \fcmd{Loop}-Befehl f�r jedes Element der Liste aus, siehe auch Kernelement~\ref{defListen}. Bei jedem Durchlauf wird das jeweilige Element der Variablen zugewiesen.\\ 
Siehe Abschnitte~\ref{Schleifen} und \ref{TexteAufteilen}, verwendet in Beispiel~\ref{BspBugPlaner}, \ref{BspPrimzahlen}, \ref{BspAbfrageschleife}, \ref{BspKombinator}, \ref{BspTerminologiepruefer}, \ref{BspTerminplaner}, \ref{BspMetaErweiterung} und \ref{BspWissensreflektion}.

\item[\fcmd{ForEachLine}] \synvar{Variable} \verb|in| \synvar{Dateipfad} $[$\verb|; Take=| \synvar{Anzahl} $]$\\
$[$\verb|; IndexVar=| \synvar{Indexvariable} $][$\verb|; SectionVar=| \synvar{Abschnittsvariable} $]$\\
$[$\verb|; NoFormat| $]$\\
Liest Zeilen der �ber den Dateipfad angegebenen Datei ein (siehe auch Kernelement~\ref{defDateipfad}) und f�hrt die nachfolgenden einger�ckten Befehle bis zum dazugeh�renden \fcmd{Loop}-Befehl wiederholt aus. F�r jede durchlaufene Zeile wird die Variable mit dem Inhalt der Zeile belegt.\par
Mit dem optionalen \cidxfrag{Take}-Parameter kann die Anzahl der durchlaufenen Zeilen vorgegeben werden. Diese werden dann zuf�llig ausgew�hlt. Bei Angabe des optionalen \cidxfrag{IndexVar}-Parameters wird die laufende Nummer in der Indexvariablen zur�ckgegeben.\par
Der Befehl ist standardm��ig auf das Einlesen von Datenlisten\index{Datenlisten} ausgerichtet, also Textdateien bei denen Leerzeilen und mit \verb|\\| beginnende Kommentarzeilen �bersprungen werden, und Zeilen, die mit einer �ffnenden eckigen Klammer beginnen und mit einer schlie�enden eckigen Klammer enden als Abschnittskennung\index{Abschnittskennung} interpretiert werden. F�r alle Zeilen werden Whitespace-Zeichen an Anfang und Ende entfernt. Bei Angabe des optionalen \cidxfrag{SectionVar}-Parameters wird der jeweilige Abschnitt in der Abschnittsvariablen zur�ckgegeben.\par
Mit dem \cidxfrag{NoFormat}-Argument werden alle Zeilen der Textdatei zur�ckgegeben, auch Leerzeilen. Es werden keine Formate wie Abschnittskennungen oder Kommentare interpretiert und kein Whitespace entfernt.

Siehe Abschnitte~\ref{RegUndDaten} und \ref{Wissensdokumentation}, verwendet in Beispiel~\ref{BspTerminologiepruefer} und \ref{BspWissensreflektion}.

\item[\fcmd{GoSub}]  \synvar{Funktionsname} $[$\verb|; BaseKey=| \synvar{Basisschl�ssel}$]$\\
F�hrt die angegebenen Funktion aus. Bei Angabe des optionalen \cidxfrag{BaseKey}-Pa\-ra\-me\-ters wird die gleichnamige Systemvariable auf den angegebenen Basisschl�ssel gesetzt. Nachdem die Funktion durchlaufen wurde, wird die Skriptausf�hrung in der n�chsten Zeile fortgesetzt.\\
Siehe Abschnitte~\ref{Funktionen} und \ref{POUnterstuetzung}, verwendet in Beispiel~\ref{BspTuermeVonHanoi}, \ref{BspZeitmesser} und \ref{BspTerminplaner}.

\item[\fcmd{GoTo}] \synvar{Sprungmarke}\\
Setzt die Skriptausf�hrung an der angegebenen Sprungmarke fort. Diese wird mit dem \fcmd{JumpMark}-Befehl gesetzt.\\
Siehe Abschnitt~\ref{Spruenge}, verwendet in Beispiel~\ref{BspSpruenge}.

\item[\fcmd{If}] \synvar{Bedingung}\\
F�hrt die nachfolgend einger�ckten Befehle genau dann aus, wenn die Bedingung erf�llt ist (siehe auch Kernelement~\ref{defBedingungen}). 
Ist die Bedingung nicht erf�llt, wird der n�chste Befehl auf gleicher Ebene ausgef�hrt.\par
Der If-Befehl kann mit beliebig vielen ElseIf-Befehlen und einem Else-Befehl kombiniert werden.

Siehe Abschnitt~\ref{IfUndBedingungen}, verwendet in Beispiel~\ref{BspBugPlaner}, \ref{BspPrimzahlen}, \ref{BspBewertungsschema}, \ref{BspTuermeVonHanoi}, \ref{Passwortgenerator}, \ref{BspTerminologiepruefer}, \ref{Archivklasse2}, \ref{BspArchivklasse3}, \ref{BspZeitmesser}, \ref{BspJourFixePlaner}, \ref{BspTerminplaner}, \ref{BspMetaErweiterung} und \ref{BspWissensreflektion}.
 
\item[\fcmd{Implies}] \synvar{Schl�ssel} \verb|=| \synvar{Wert}\\
Setzt einen Wert f�r einen Schl�ssel in der URL. Damit k�nnen z.B. Antworten f�r noch ausstehende Auswahlabfragen impliziert werden.
 
\item[\fcmd{Include}] \synvar{Dateipfad} $[$\verb|; BaseKey=| \synvar{Basisschl�ssel}$]$\\
F�hrt die �ber den Dateipfad angegebene Funktionsdatei aus. Die standardm��ige Dateiendung f�r Funktionsdateien ist \emph{.fps}. Die Syntax ist diesselbe wie in den Skriptdateien.
Bei Angabe des optionalen \cidxfrag{BaseKey}-Parameters wird die gleichnamige Systemvariable auf den angegebenen Basisschl�ssel gesetzt. Nachdem die Datei durchlaufen wurde, wird die Skriptausf�hrung in der n�chsten Zeile fortgesetzt.\\
Siehe Abschnitt~\ref{POUnterstuetzung}, verwendet in Beispiel~\ref{BspListenfunktion} und \ref{BspPatchskript}.

\item[\fcmd{Input}] \synvar{Schl�ssel} \verb|:| \synvar{Anzeigetext}\\
Fragt eine Texteingabe ab und verwaltet diese unter dem angegebenen Schl�ssel. Der Anzeigetext wird oberhalb des Eingabefeldes angezeigt.
Der Schl�ssel wird verwendet, um die Eingabe in der URL zu speichern, und innerhalb des Skriptes abzufragen. Er muss eindeutig sein, bzw. wird bei Wiederholung als schon gesetzt betrachtet. Innerhalb einer Datei k�nnen die Schl�ssel durch Anh�ngen des \verb|'|-Zeichens automatisch durchnummeriert werden.\\
Siehe Abschnitte~\ref{Texteingaben} und \ref{TexteAufteilen}, verwendet in Beispiel~\ref{BspTextersetzung}, \ref{BspAbschnittsverschiebung}, \ref{BspBugPlaner}, \ref{BspHilfezeilen}, \ref{BspAbfrageschleife}, \ref{BspBewertungsschema}, \ref{BspKombinator}, \ref{BspTerminologiepruefer}, \ref{BspSuchlink}, \ref{Archivklasse2}, \ref{BspListenfunktion}, \ref{BspZeitmesser}, \ref{BspTerminplaner}, \ref{BspLizenzabfrage}, \ref{BspMetaErweiterung}, \ref{BspWissensreflektion} und \ref{BspPatchskript}.

\item[\fcmd{JumpMark}] \synvar{Sprungmarke}\\
Setzt eine Sprungmarke, die mit dem \fcmd{GoTo}-Befehl angesprungen werden kann.\\
Siehe Abschnitt~\ref{Spruenge}, verwendet in Beispiel~\ref{BspSpruenge}.

\item[\fcmd{ListDirectories}] \synvar{Listenvariable} \verb|=| \synvar{Verzeichnispfad} $[$\verb|; Pattern=| \synvar{Suchmuster}$]$\\
Listet alle Verzeichnisnamen im Verzeichnispfad auf, die dem Suchmuster entsprechen (siehe Kernelement~\ref{defDateipfad}). Das Ergebnis wird als Liste zu der Listenvariable gespeichert (siehe Kernelement~\ref{defListen}). Der Eintrag mit dem Index $0$ gibt die Anzahl der gefundenen Verzeichnisse zur�ck.\\
Siehe Abschnitt~\ref{Wissensdokumentation}, verwendet in Beispiel~\ref{BspWissensreflektion}.

\item[\fcmd{ListFiles}] \synvar{Listenvariable} \verb|=| \synvar{Verzeichnispfad} $[$\verb|; Pattern=| \synvar{Suchmuster}$]$\\
Listet alle Dateinamen im Verzeichnispfad auf, die dem Suchmuster entsprechen (siehe Kernelement~\ref{defDateipfad}). Das Ergebnis wird als Liste zu der Listenvariable gespeichert (siehe Kernelement~\ref{defListen}). Der Eintrag mit dem Index $0$ gibt die Anzahl der gefundenen Dateien zur�ck.\\
Siehe Abschnitt~\ref{Wissensdokumentation}, verwendet in Beispiel~\ref{BspWissensreflektion}.

\item[\fcmd{Loop}]\mbox{}\\
Beendet einen Schleifendurchlauf und beginnt den n�chsten Schleifendurchlauf. Eine Schleife kann durch verschiedene Befehle initiiert werden.\\
Siehe Abschnitt~\ref{Schleifen}, verwendet in Beispiel~\ref{BspBugPlaner}, \ref{BspPrimzahlen}, \ref{BspAbfrageschleife}, \ref{BspKombinator}, \ref{Passwortgenerator}, \ref{BspTerminologiepruefer}, \ref{BspZeitmesser}, \ref{BspJourFixePlaner}, \ref{BspTerminplaner}, \ref{BspMetaErweiterung} und \ref{BspWissensreflektion}.

\item[\fcmd{MoveSection}] \synvar{Ausgangsabschnitt} \verb|->| \synvar{Zielabschnitt}\\
�bertr�gt den Inhalt eines Ausgangsabschnitts in einen Zielabschnitt. Der Ausgangsabschnitt wird gel�scht.
Ist der Zielabschnitt noch nicht vorhanden, so wird dieser am Ende angelegt. Ist der Ausgangsabschnitt nicht vorhanden, so bleibt der Aufruf ohne Wirkung. Sind Ausgangsabschnitt und Zielabschnitt identisch, so wird dieser Aschnitt entfernt. Die beien aufeinandertreffenden Textbl�cke werden zusammengef�hrt, wenn sie die gleiche Formatierung haben.\\
Siehe Abschnitte~\ref{AusgabetitelUndAbschnittsverschiebungen} und \ref{AssetDokumentation}, verwendet in Beispiel~\ref{BspAbschnittsverschiebung}, \ref{BspMetaErweiterung} und \ref{BspWissensreflektion}.

\item[\fcmd{Random}] \synvar{Variable} \verb|=| \synvar{Untergrenze} \verb|..| \synvar{Obergrenze}\\
Erzeugt eine ganze Zufallszahl im Bereich von Untergrenze bis Obergrenze und speichert diese in einer Variablen.
Bei Aktualisierung der Seite wird jeweils ein neuer Zufallswert generiert.\\
Siehe Abschnitt~\ref{ErsetzungenUndZufallsgenerierung}, verwendet in Beispiel~\ref{Passwortgenerator}.

\item[\fcmd{RegExMatch}] \synvar{Listenvariable} \verb|=| \synvar{Text} \verb=|= \synvar{Regul�rer Ausdruck}\\
Wendet einen regur�ren Ausdruck auf einen Text an und speichert das Ergebnis in einer Liste zur angegebenen Listenvariablen. In der Variablen mit Index $0$ wird \verb|true| zur�ckgegeben, wenn der Text dem regul�ren Ausdruck entspricht, ansonsten \verb|false|. In den nachfolgenden Variablen mit den Indices $1$, $2$, \dots werden die Werte der Gruppen abgelegt, die im Ausdruck enthalten sind.\\
Siehe Abschnitt~\ref{RegUndDaten}, verwendet in Beispiel~\ref{BspTerminologiepruefer}, \ref{Archivklasse2} und \ref{BspTerminplaner}.

\item[\fcmd{RegExReplace}] \synvar{Variable} \verb|=| \synvar{Text} \verb=|= \synvar{Regul�rer Ausdruck} \verb|->| \synvar{Ersetzungstext}\\
Ersetzt in einem Text alle Vorkommen, die dem regul�ren Ausdruck entsprechen, durch einen Ersetzungstext und speichert das Ergebnis in einer Variablen.
 
\item[\fcmd{Replace}] \synvar{Variable} \verb|=| \synvar{Text} \verb=|= \synvar{Suchtext} \verb|->| \synvar{Ersetzungstext}\\
Ersetzt alle Vorkommen des Suchtextes in einem Text durch den Ersetzungstext und speichert das Ergebnis in einer Variablen.\\
Siehe Abschnitte~\ref{Texteingaben}, \ref{Funktionen} und \ref{ErsetzungenUndZufallsgenerierung}, verwendet in Beispiel~\ref{BspTextersetzung}, \ref{BspTuermeVonHanoi}, \ref{BspKombinator}, \ref{Passwortgenerator}, \ref{BspZeitmesser}, \ref{BspMetaErweiterung} und \ref{BspWissensreflektion}.

\item[\fcmd{Return}]\mbox{}\\
Beendet eine Funktion und setzt die Skriptausf�hrung in der n�chsten Zeile des \fcmd{GoSub}-Befehls fort, an dem die Funktion aufgerufen wurde.\\
Siehe Abschnitt~\ref{Funktionen}, verwendet in Beispiel~\ref{BspTuermeVonHanoi}, \ref{BspArchivklasse4}, \ref{BspZeitmesser} und \ref{BspTerminplaner}.

\item[\fcmd{Round}] \synvar{Variable} \verb|=| \synvar{Dezimalzahl} \verb=|= \synvar{Anzahl an Stellen}\\
Rundet eine Dezimalzahl auf die angegebene Anzahl an Stellen und speichert das Ergebnis in einer Variablen.\\
Siehe Abschnitt~\ref{Berechnungen}, verwendet in Beispiel~\ref{BspBerechnungen}.

\item[\fcmd{Set}] \synvar{Variable} \verb|=| \synvar{Wert}\\
Speichert einen Wert in einer Variablen. Variablenbezeichnungen k�nnen aus Buchstaben, Ziffern und runden Klammern bestehen und auch wieder andere Variablen enthalten. Die in Variablennamen verwendeten Klammern haben keine syntaktische Bedeutung, k�nnen aber die Lesbarkeit von Variablen verbessern.\\
Siehe Abschnitt~\ref{VariablenSetzenUndVerwenden}, verwendet in Beispiel~\ref{BspHalloWeltMitVariable}, \ref{BspVariablenersetzung}, \ref{BspBewertungssystem}, \ref{BspAbschnittsverschiebung}, \ref{BspBugPlaner}, \ref{BspPrimzahlen}, \ref{BspAbfrageschleife}, \ref{BspBewertungsschema}, \ref{BspTuermeVonHanoi}, \ref{BspKombinator}, \ref{Passwortgenerator}, \ref{BspTerminologiepruefer}, \ref{Archivklasse2}, \ref{BspListenfunktion}, \ref{BspZeitmesser}, \ref{BspJourFixePlaner}, \ref{BspTerminplaner}, \ref{BspLizenzabfrage}, \ref{BspMetaErweiterung}, \ref{BspWissensreflektion} und \ref{BspPatchskript}.

\item[\fcmd{SetCulture}] \synvar{Variable} \verb|=| \synvar{Kultur-Kennzeichen}\\
Setzt die Kultureinstellungen anhand des angegebenen Kultur-Kennzeichens, z.B. "de-DE". Diese werden z.B.\ bei formatierten Datumsausgaben angewendet wie der Ausgabe von Wochentagen.\\
Siehe Abschnitt~\ref{MeetingVorbereitung}, verwendet in Beispiel~\ref{BspJourFixePlaner}.

\item[\fcmd{SetDateTime}] \synvar{Variable} \verb|=| \synvar{Format}\\
Speichert das aktuelle Datum oder den aktuellen Zeitpunkt in einer vorgegebenen Formatierung in einer Variablen.\\
Siehe Abschnitte~\ref{Zeitmessung} und \ref{MeetingVorbereitung}, verwendet in Beispiel~\ref{BspZeitmesser}, \ref{BspJourFixePlaner} und \ref{BspPatchskript}.

\item[\fcmd{SetInputDescription}] \synvar{Beschreibung}\\
Setzt die Beschreibung f�r die Eingabeseiten. Diese wird dort unterhalb der �berschrift angezeigt.\\
Siehe Abschnitt~\ref{FormatierungDerEingabe}, verwendet in Beispiel~\ref{BspBugPlaner}.

\item[\fcmd{SetInputSection}] \synvar{Abschnitts�berschrift}\\
Setzt die Abschnitts�berschrift f�r die nachfolgenden Eingabebefehle.\\
Siehe Abschnitt~\ref{FormatierungDerEingabe}, verwendet in Beispiel~\ref{BspBugPlaner} und \ref{BspWissensreflektion}.

\item[\fcmd{SetInputTitle}] \synvar{Titel}\\
Setzt den Titel f�r die Eingabeseiten.\\
Siehe Abschnitt~\ref{FormatierungDerEingabe}, verwendet in Beispiel~\ref{BspBugPlaner} und \ref{BspZeitmesser}.

\item[\fcmd{SetSection}] \synvar{Abschnitts�berschrift}\\
Setzt die Abschnitts�berschrift f�r die nachfolgenden Ausgabebefehle. Der Befehl ist eine Alternative zu der \verb|@|-Syntax der Ausgabebebefehle.

\item[\fcmd{SetStopCounter}] \synvar{Schleifengrenze} \verb|;| \synvar{Befehlsgrenze}\\
Setzt die Obergrenzen f�r die Anzahl der Schleifendurchl�ufe und der durchlaufenenen Befehle pro Seitenaufruf.
Beim �berschreiten der Grenzen bricht das Programm mit einer Fehlermeldung ab.
Die Grenzen dienen zur Verhinderung von Endlosschleifen und Endlosrekursionen und stehen Werksseitig auf 1000 und 20000, was f�r die meisten Skripte ausreichen sollte.
Ist f�r ein komplexes Skript absehbar, dass diese Grenzen im regul�ren Betrieb �berschritten werden, k�nnen die Grenzen mit dem o.g.\ Befehl hochgesetzt werden.\\
Siehe Abschnitte~\ref{Schleifen} und \ref{Terminplanung}, verwendet in Beispiel~\ref{BspTerminplaner}.

\item[\fcmd{SetTitle}] \synvar{Titel}\\
Setzt den Titel f�r die Ausgabeseite.\\
Siehe Abschnitt~\ref{AusgabetitelUndAbschnittsverschiebungen}, verwendet in Beispiel~\ref{BspAbschnittsverschiebung} und \ref{BspBugPlaner}.

\item[\fcmd{Sort}] \synvar{Ergebnisliste} \verb|=| \synvar{Ausgangsliste} $[$\verb|To| \synvar{Indexliste}$]$\\
Sortiert eine Ausgangsliste alphanumerisch aufteigend und speichert das Ergebnis als Ergebnisliste (siehe Kernelement~\ref{defListen}). Wird das \verb|To|-Attribut mit einer Indexliste angegeben, so wird dort die Indexzuordung eingetragen, mit der die Umsortierung auf weitere Listen �bertragen werden kann.

\item[\fcmd{Split}] \synvar{Variable} \verb|=| \synvar{Text} \verb=|= \synvar{Trennzeichen}\\
Teilt einen Text anhand eines Trennzeichens auf und speichert das Ergebnis als Liste (siehe Kernelement~\ref{defListen}). Als Trennzeichen sind auch l�ngere Zeichenkombinationen zul�ssig.
�ber den Eintrag mit dem Index 0 kann die Anzahl der Listenelemente abgerufen werden. Bei der Aufteilung wird Whitespace vor und nach den Trennzeichen standardm��ig entfernt. Wenn dieser erhalten bleiben soll, muss er mit \cidxvar{Chr(...)}-Befehlen maskiert werden.\\
Siehe Abschnitte~\ref{TexteAufteilen} und \ref{Terminplanung}, verwendet in Beispiel~\ref{BspBugPlaner}, \ref{BspKombinator}, \ref{Passwortgenerator}, \ref{BspTerminplaner} und \ref{BspMetaErweiterung}.

\item[\fcmd{ToLower}] \synvar{Variable} \verb|=| \synvar{Text}\\
Wandelt einen Text in Kleinbuchstaben um und speichert das Ergebnis in einer Variablen.\\
Siehe Abschnitt~\ref{Entwicklungsanleitungen}, verwendet in Beispiel~\ref{Archivklasse2}.

\item[\fcmd{ToUpper}] \synvar{Variable} \verb|=| \synvar{Text}\\
Wandelt einen Text in Gro�buchstaben um und speichert das Ergebnis in einer Variablen.\\
Siehe Abschnitt~\ref{Entwicklungsanleitungen}.

\item[\fcmd{Trim}] \synvar{Variable} \verb|=| \synvar{Text}\\
Entfernt umschlie�enden Whitespace aus dem Text und speichert das Ergebnis in einer Variablen.

\item[\fcmd{UrlEncode}] \synvar{Variable} \verb|=| \synvar{Text}\\
Wandelt einen Text in die URL-Codierung um und speichert das Ergebnis in einer Variablen.\\
Siehe Abschnitt~\ref{ParametrisierteLinks}, verwendet in Beispiel~\ref{BspSuchlink}.

\item[\fcmd{XmlEncode}] \synvar{Variable} \verb|=| \synvar{Text}\\
Wandelt einen Text in die XML-Codierung um und speichert das Ergebnis in einer Variablen. Die XML-Codierung maskiert die Zeichen \verb|<|, \verb|>|, \verb|&|, sowie die Anf�hrungszeichen, so dass der Text innerhalb von XML-Tags stehen kann.\\
Siehe Abschnitt~\ref{Personalisierung}, verwendet in Beispiel~\ref{BspPatchskript}.
\end{description}
\printindex
\end{document}